\chapter{Computation, equality and normal forms}
\label{prov:ch:comp}
\label{comp:ch}

Three of the eight rules of the native scalar core $\Nat$
(\cref{base:sec:kernel}) concern numerical terms: the zero computation
equation, the successor computation equation and the transport rule. This
chapter collects the calculi in which these three rules are merged, and the
calculi in which one or both computation equations leave the rule inventory
and become part of the syntax of terms. Two ideas organize the material.

The first idea keeps the terms, the formulas and ordinary substitution
unchanged. Reflexivity and one rule \rulename{Eval}, which composes an
equation with one computation step, replace the three rules
(\cref{comp:sec:eval}).

The second idea restricts the terms to the normal forms of a computation and
lets every substitution of a rule be followed by normalization. A
computation that is normalized needs no rule. \Cref{comp:sec:normalizers}
proves once, for an arbitrary normalizer that satisfies five conditions, that
derivations transfer between $\Nat$ and the kernel on normal forms. The zero
computation alone is normalized in \cref{comp:sec:zero}; on the resulting
zero-normal terms the remaining computation and the equality reasoning are
presented by the rules of
\cref{comp:sec:reflexive,comp:sec:seed,comp:sec:step,comp:sec:seedeval}.
A weaker zero normalization, which contracts only zero iterators with a
variable or $0$ as initializer and is not coherent with substitution, is
treated in \cref{comp:sec:atomic} with rules that keep their raw instances.
Both computations are normalized in \cref{comp:sec:full}.

Every calculus $\mathsf{K}$ of the chapter is given with its complete rule
inventory and with a theorem of one of two forms: $\mathsf{K}$ derives
exactly the sequents of $\Nat$, or $\mathsf{K}$ derives the normal form of a
sequent exactly when $\Nat$ derives the sequent. In both cases the
bi-interpretation of \cref{base:thm:bridge} transfers
(\cref{comp:cor:bridge}). \Cref{comp:tab:calculi} lists the calculi; the
rule names are explained in the sections cited.

\begin{table}[ht]
\centering
\small
\begin{tabular}{@{}lp{0.17\linewidth}p{0.44\linewidth}l@{}}
\toprule
Calculus & Terms & Rules & Equivalence\\
\midrule
$\mathsf{E}_3$ & raw & \rulename{Hyp}, \rulename{Refl}, \rulename{Eval} (three cases), \rulename{QE}, \rulename{DE}, \rulename{WI} & \cref{comp:thm:eval}\\
$\mathsf{E}_2$ & raw & the same, \rulename{Eval} with two cases & \cref{comp:thm:root}\\
$\mathsf{R}^{\mathrm{raw}}$ & raw & \rulename{Hyp}, \rulename{I0}, \rulename{IS}, \rulename{Tr}, \rulename{QE}, \rulename{DE}, \rulename{RI} & \cref{comp:prop:rawreflexive}\\
$\Nat^{0}$ & zero-normal & \rulename{Hyp}, \rulename{IS}, \rulename{Tr}, \rulename{QE}, \rulename{DE}, \rulename{DI}, \rulename{QI} & \cref{comp:thm:zero}\\
$\mathsf{Z}_6$ & zero-normal & \rulename{Hyp}, \rulename{IS}, \rulename{Tr}, \rulename{QE}, \rulename{DE}, \rulename{WI} & \cref{comp:thm:zsix}\\
$\mathsf{R}$ & zero-normal & \rulename{Hyp}, \rulename{IS}, \rulename{Tr}, \rulename{QE}, \rulename{DE}, \rulename{RI} & \cref{comp:thm:reflexive}\\
$\mathsf{S}$ & zero-normal & \rulename{Hyp}, \rulename{IS}, \rulename{Seed}, \rulename{QE}, \rulename{DE}, \rulename{WI} & \cref{comp:thm:seed}\\
$\mathsf{U}$ & zero-normal & \rulename{Hyp}, \rulename{Refl}, \rulename{Step}, \rulename{QE}, \rulename{DE}, \rulename{WI} & \cref{comp:thm:step}\\
$\mathsf{C}_{\mathrm{Step}}$ & zero-normal, clauses & \rulename{Use}, \rulename{Refl}, \rulename{Step}, \rulename{DE}, \rulename{WI} & \cref{comp:thm:clausal}\\
$\mathsf{C}_{\mathrm{Seed}}$ & zero-normal, clauses & \rulename{Use}, \rulename{Refl}, \rulename{SeedEval}, \rulename{DE}, \rulename{WI} & \cref{comp:thm:seedeval}\\
$\mathsf{C}_{\mathrm{App}}$ & unit-normal, clauses & \rulename{Use}, \rulename{Refl}, \rulename{App}, \rulename{DE}, \rulename{WI} & \cref{comp:thm:app}\\
$\mathsf{A}_6$ & atomic-seed normal, raw metadata & \rulename{Hyp}, \rulename{IS}, \rulename{Tr}, \rulename{QE}, \rulename{DE}, \rulename{WI} & \cref{comp:thm:asix}\\
$\Nat^{\mathrm{f}}$ & full normal & \rulename{Hyp}, \rulename{Refl}, \rulename{Tr}, \rulename{QE}, \rulename{DE}, \rulename{DI}, \rulename{QI} & \cref{comp:thm:full}\\
$\mathsf{F}$ & full normal & \rulename{Hyp}, \rulename{Refl}, \rulename{Tr}, \rulename{QE}, \rulename{DE}, \rulename{WI} & \cref{comp:thm:fsix}\\
\bottomrule
\end{tabular}
\caption{The calculi of the chapter.}
\label{comp:tab:calculi}
\end{table}

%%%%%%%%%%%%%%%%%%%%%%%%%%%%%%%%%%%%%%%%%%%%%%%%%%%%%%%%%%%%%%%%%%%%%%%%%%%%
\section{Conventions and shared lemmas}
\label{comp:sec:conv}

This section fixes the notation common to the ten packages of the chapter
(\pkg{computational\_equality}, \pkg{equality\_from\_iteration},
\pkg{uniform\_step}, \pkg{seed\_evaluation}, \pkg{zero\_normal},
\pkg{reflexive\_base}, \pkg{initializer\_equality},
\pkg{full\_normalization}, \pkg{computation\_normal},
\pkg{atomic\_seed\_refinement}) and proves once the
arguments that several of them repeat.

\subsection{Notation}
\label{comp:sec:notation}

\emph{Raw terms} are the terms of \cref{base:sec:terms}, with ordinary
capture-free substitution; no computation is performed by substitution, so
$\It{0}{a}{u}{h}$ is a raw term distinct from $a$. The relation $s\equiv t$ is
literal identity of syntax, up to the renaming of bound variables; in a
de Bruijn representation it is equality of finite trees. Every test in a
rule of this chapter is a test of literal identity and never a test of
provable equality. A formula is a \emph{literal equation} if its outermost
constructor is equality; an equation below $\Q{y}$ or $\D{x}{y}$ is not a
literal equation. The numeral $1$ is $\Suc 0$.

A \emph{template} is a term $h$ with a distinguished variable $v$; it may
contain $v$ several times, not at all, or below iterator binders, and it may
contain other variables. Its instance at $s$ is written $h(s)$. A
\emph{matrix} is a formula $B$ with one or two distinguished variables;
$B(r)$ and $B(t,r)$ are its instances. Instances are formed with the
substitution of the calculus at hand: ordinary substitution on raw terms and
the normalizing substitution of \cref{comp:def:nusubst} on normal forms.

The rules of $\Nat$ are those of \cref{base:sec:kernel}, with their names:
\rulename{Hyp}, \rulename{I0}, \rulename{IS}, \rulename{Tr} (from $s=t$ and
$h(s)=r$ infer $r=h(t)$), \rulename{QE}, \rulename{DE}, \rulename{DI} (by
cases) and \rulename{QI} (by progress, for literal equations). The first
premise of \rulename{Tr} is its \emph{major} premise and the second its
\emph{minor} premise. Weakening and substitution of terms in derivations
(\cref{base:lem:struct}), cut (\cref{base:lem:cut}), the equality rules
(\cref{base:lem:eq}) and the Leibniz rule (\cref{base:thm:leibniz}) are
admissible in $\Nat$.

\begin{definition}[Witness-selected introduction]
\label{comp:def:wi}
\rulename{WI} is the rule $\rulename{DI}_{\mathrm{w}}$ of
\cref{kern:def:W}, with the function $\eta_{B,t}$ defined there, read with
the terms and the substitution of the calculus at hand; the test $t\equiv0$
is applied to the witness as it is supplied in that syntax. An instance with
$\eta_{B,t}(z)=z$ is in \emph{progress mode}; every other instance is in
\emph{case mode} and is an instance of \rulename{DI}.
\end{definition}

On raw terms \rulename{WI} replaces \rulename{DI} and \rulename{QI} by
\cref{kern:thm:witness}, whose proof repairs the one missing instance of
\rulename{DI} at the term $\It{0}{0}{u}{0}$ of \cref{kern:lem:alias}. That
term is not available on normal forms. \Cref{comp:lem:single} states the
argument for an arbitrary syntax and an arbitrary term that is provably $0$
and not literally $0$; the sections below supply the term.

\subsection{Equality from reflexivity and crossed transport}

\begin{lemma}
\label{comp:lem:eqalg}
Let $\mathsf{K}$ be a calculus that has \rulename{Tr} for every template,
with a substitution for which the instance of the identity template at $s$
is $s$.
\begin{enumerate}[label=(\alph*)]
\item If $\Gamma\proves_{\mathsf K}\ell=r$ for some term $\ell$, then
 $\Gamma\proves_{\mathsf K}r=r$.
\item If $\mathsf{K}$ derives $r=r$ for every term $r$, then the following
 rules are admissible in $\mathsf{K}$, by derivations that use only
 reflexivity and \rulename{Tr}: symmetry; transitivity; congruence, from
 $s=t$ infer $h(s)=h(t)$; and transport in a literal equation, from $s=t$
 and $a(s)=b(s)$ infer $a(t)=b(t)$.
\end{enumerate}
\end{lemma}

\begin{proof}
(a) is the derivation of \cref{base:lem:eq}(i) with the equation $\ell=r$
in place of the instance of \rulename{I0}: \rulename{Tr} with the identity
template and $\ell=r$ as both premises concludes $r=r$. (b) The derivations
of \cref{base:lem:eq}(ii)--(v) use reflexivity and \rulename{Tr} only, one
application for each of symmetry, transitivity and congruence and two for
transport in an equation; they apply to $\mathsf{K}$ with the instances of
the templates formed by the substitution of $\mathsf{K}$.
\end{proof}

For a calculus on raw terms that contains \rulename{I0} the lemma is
\cref{kern:lem:local}(iii). It is stated here for the calculi in which
\rulename{I0} is not a rule and reflexivity has another source.

\subsection{The single introduction}

\begin{lemma}[The single introduction]
\label{comp:lem:single}
Let $\mathsf{K}$ be a calculus on native formulas over a term syntax with a
substitution for which instances compose (ordinary substitution on raw
terms, or a normalizing substitution, by \cref{comp:lem:laws}). Assume that
$\mathsf{K}$ has \rulename{Hyp}, \rulename{QE} and \rulename{DE}; that
weakening, injective renaming of free variables and cut on an assumption are
admissible in $\mathsf{K}$; and that symmetry of equality and transport in
literal equations are admissible in $\mathsf{K}$.
\begin{enumerate}[label=(\roman*)]
\item If $\mathsf{K}$ has \rulename{WI}, then \rulename{QI} is admissible in
 $\mathsf{K}$, and so is the introduction of a closure by cases, for every
 matrix $B$:
 \[
  \infrule{\Gamma,\,B(0)\proves B(\Suc z)}{\Gamma\proves\Q{y}B(y)}
  \qquad(z\text{ fresh}).
 \]
 Every instance of \rulename{DI} is an instance of \rulename{WI}, except
 when $B$ is a literal equation and $t\equiv0$.
\item If $\mathsf{K}$ has \rulename{WI}, $E(x,y)$ is a literal equation and
 $Z$ is a term with $Z\not\equiv0$ and $\Gamma\proves_{\mathsf K}Z=0$, then
 \[
  \Gamma,\,E(0,0)\proves_{\mathsf K}E(0,\Suc z)\quad(z\text{ fresh})
  \qquad\text{implies}\qquad
  \Gamma\proves_{\mathsf K}\D{x}{y}E(x,y).
 \]
\item If $\mathsf{K}$ has \rulename{DI} and \rulename{QI}, then
 \rulename{WI} is admissible in $\mathsf{K}$.
\end{enumerate}
\end{lemma}

\begin{proof}
The three parts are the three constructions in the proof of
\cref{kern:thm:witness}, which use of the calculus only the rules and the
admissible operations assumed here.

(i) For \rulename{QI}, the construction given there applies: \rulename{WI}
for the matrix $E(y)$ with an unused witness coordinate, at the witness
$0$, is in progress mode and has the given premise, and \rulename{DE} with
\rulename{Hyp} returns $\Q{y}E(y)$. For the closure by cases use the same
construction with $B(y)$ in place of $E(y)$ and the witness $1$: since
$1\not\equiv0$, \rulename{WI} is in case mode for every $B$, and its premise
is $\Gamma,B(0)\proves B(\Suc z)$. In $\Nat$ this rule is
\cref{base:lem:qI}, with the witness $0$; with \rulename{WI} the witness
must not be the literal $0$. The last assertion is the definition of the
two modes.

(ii) This is the repair in the proof of \cref{kern:thm:witness}, with $Z$ in
place of the term $c$ of \cref{kern:lem:alias}, after $z$ is renamed so that
it is fresh for $Z$: in the context $\Gamma,E(Z,0)$, transport along $Z=0$,
cut with the given derivation, and transport along $0=Z$ give
$E(Z,\Suc z)$, and \rulename{WI} at the witness $Z$ is in case mode. Only
transport in literal equations is used, and no term is substituted into the
given derivation.

(iii) An instance of \rulename{WI} in case mode is an instance of
\rulename{DI}. An instance in progress mode, with the premise
$\Gamma,E(0,z)\proves E(0,\Suc z)$, is derived as in
\cref{base:cor:combined}(i), by \rulename{QI}, \rulename{QE} under the
assumption $E(0,0)$, and \rulename{DI} at the witness $0$.
\end{proof}

\begin{remark}
\label{comp:rem:single}
The term $Z$ of part (ii) is $\It{0}{0}{u}{0}$ on raw terms
(\cref{kern:lem:alias}), the terms $\It{1}{0}{u}{0}$ or $\It{1}{0}{u}{u}$ on
zero-normal terms (\cref{comp:lem:alias}), and a variable $b$ opened by
\rulename{DE} on full normal forms (\cref{comp:lem:freshzero}). As on raw
terms (\cref{kern:rem:wsubst}), no calculus with \rulename{WI} is shown
closed under the substitution of terms by a direct induction on
derivations; on normal forms a substitution can in addition create a redex
whose contraction turns a witness into the literal $0$. Substitution is
transferred from an equivalent calculus with \rulename{DI} and
\rulename{QI}. Weakening, injective renaming and cut leave the witness and
the root of the matrix unchanged and are proved directly, as in
\cref{kern:lem:struct}.
\end{remark}

\subsection{Congruence below binders and provability of conversion}
\label{comp:sec:conversion}

Transport in $\Nat$ replaces a term in a template by capture-free
substitution. It does not replace the body of an iterator by a pointwise
equal body, and it does not replace a matrix below $\Q{y}$ or $\D{x}{y}$.
Both replacements are derived rules of $\Nat$: the congruence of the body
of an iterator is \cref{base:cor:rot}(ii), which rests on an equation
induction, and the congruence of the binders is
\cref{base:lem:bindercong}. Every comparison of $\Nat$ with a calculus on
normal forms depends on them, because a normalizer rewrites below binders.

\begin{definition}[Computation reduction]
\label{comp:def:reduction}
The relation $\to$ on raw terms is the closure, under every term context
including the bodies of iterators, of the two contractions
\[
 \It{0}{a}{u}{h}\to a,\qquad
 \It{\Suc n}{a}{u}{h}\to h\subst{\It{n}{a}{u}{h}}{u}.
\]
The first is the \emph{zero contraction} and the second the \emph{successor
contraction}; $\to^{*}$ is the reflexive and transitive closure.
\end{definition}

\begin{lemma}[Conversion is provable]
\label{comp:lem:conversion}
If $t\to^{*}t'$, then $\proves_{\Nat}t=t'$, with all free variables of $t$
as parameters. If the formula $A'$ is obtained from the formula $A$ by
replacing each term by a reduct of it, then $\Gamma,A\proves_{\Nat}A'$ and
$\Gamma,A'\proves_{\Nat}A$ for every $\Gamma$.
\end{lemma}

\begin{proof}
For one step $t\to t'$ argue by induction on the position of the contracted
redex. At the root the equation is an instance of \rulename{I0} or
\rulename{IS}. Below a successor, in the counter or in the initializer of an
iterator, it follows from the induction hypothesis by congruence
(\cref{base:lem:eq}(iv)) in the template with one hole at that position. In
the body of
$\It{n}{a}{u}{h}$, with $h\to h'$, the induction hypothesis is
$\proves_{\Nat}h=h'$ with $u$ free, in the empty context, and
\cref{base:cor:rot}(ii) gives $\It{n}{a}{u}{h}=\It{n}{a}{u}{h'}$.
Transitivity extends the statement to $\to^{*}$.

For formulas argue by induction on $A$. At a literal equation use the two
term equations and the equality rules. At $\Q{y}B$ and $\D{x}{y}B$ the
induction hypothesis holds for the matrix with the bound variables free, in
both directions, and \cref{base:lem:bindercong} applies.
\end{proof}

\begin{remark}
The congruence of the body is restated as a lemma in \pkg{zero\_normal},
\pkg{initializer\_equality}, \pkg{full\_normalization} and
\pkg{computation\_normal}, each time by a direct equation induction on
$\It{z}{a}{u}{h}=\It{z}{a}{u}{h'}$ in place of the uniqueness lemma used
for \cref{base:cor:rot}; its version in $\Teq$ is formalized as
\Lean{Q.iter\_body\_congr} and \Lean{Q.iter\_congr}. The pointwise equation
of the bodies is derived in a context without $u$; it is not a hypothesis
that may vary with the state of the iteration. The equality reasoning of
the calculi below is recovered without induction, whereas the comparison of
each calculus on normal forms with $\Nat$ uses the equation induction
inside \cref{base:cor:rot}(ii). The packages \pkg{uniform\_step} and
\pkg{seed\_evaluation} state this dependency explicitly.
\end{remark}

%%%%%%%%%%%%%%%%%%%%%%%%%%%%%%%%%%%%%%%%%%%%%%%%%%%%%%%%%%%%%%%%%%%%%%%%%%%%
\section{The rule Eval on raw terms}
\label{comp:sec:eval}

This section draws on \pkg{computational\_equality}, which introduces the
rule \rulename{Eval} with three cases, and on
\pkg{equality\_from\_iteration}, which shows that the first case is
redundant and that neither of the other two is. The terms, the formulas and
substitution are those of $\Nat$; nothing is normalized.

\subsection{The operation \texorpdfstring{$\Phi$}{Phi} and the calculi}

\begin{definition}
\label{comp:def:phi}
For raw terms $s,t,b$ the partial operation $\Phi(s,t,b)$ is given by the
rows
\[
\begin{array}{l|l|l}
 \text{shape of }b & \text{required source }s & \Phi(s,t,b)\\ \hline
 b\equiv s & s & t\\
 b\equiv\It{0}{a}{u}{h} & s\equiv a & t\\
 b\equiv\It{\Suc c}{a}{u}{h} & s\equiv\It{c}{a}{u}{h} & h\subst{t}{u}
\end{array}
\]
and is undefined if no row applies.
\end{definition}

The rows are disjoint. In the second row $s$ is a proper subterm of $b$, and
in the third row $b$ has exactly one node more than $s$, so neither overlaps
the first; the second and third rows have different constructors at the root
of the counter. Hence $\Phi$ is a deterministic partial operation on syntax.
The third row requires an equation for the whole preceding iterate
$\It{c}{a}{u}{h}$; it does not compare initializers, and it applies
unchanged when $c\equiv0$.

\begin{definition}[The calculi $\mathsf{E}_3$ and $\mathsf{E}_2$]
\label{comp:def:e3}
The calculus $\mathsf{E}_3$ has raw terms, native formulas, ordinary
substitution, finite lists of reusable assumptions as contexts, and exactly
the six rules \rulename{Hyp}, \rulename{QE}, \rulename{DE}, \rulename{WI}
(\cref{comp:def:wi}) and
\[
 \infrule{}{\Gamma\proves s=s}\quad\text{(\rulename{Refl})},
 \qquad
 \infrule{\Gamma\proves s=t\qquad\Gamma\proves b=r}
         {\Gamma\proves r=\Phi(s,t,b)}\quad\text{(\rulename{Eval})},
\]
the latter whenever $\Phi(s,t,b)$ is defined. The three rows of $\Phi$ give
the three cases of \rulename{Eval}:
\begin{align*}
 &\infrule{\Gamma\proves s=t\qquad\Gamma\proves s=r}{\Gamma\proves r=t}
   &&\text{(\rulename{Root})}\\[4pt]
 &\infrule{\Gamma\proves a=t\qquad\Gamma\proves\It{0}{a}{u}{h}=r}
          {\Gamma\proves r=t}
   &&\text{(\rulename{Eval}$_0$)}\\[4pt]
 &\infrule{\Gamma\proves\It{c}{a}{u}{h}=t\qquad
           \Gamma\proves\It{\Suc c}{a}{u}{h}=r}
          {\Gamma\proves r=h\subst{t}{u}}
   &&\text{(\rulename{Eval}$_{\Suc}$)}
\end{align*}
The calculus $\mathsf{E}_2$ is $\mathsf{E}_3$ without the case
\rulename{Root}. In \rulename{DE} the eigenvariable is fresh for the
context, the conclusion and the parameters of the matrix. No other logical,
equality, successor, computation, induction or structural rule is primitive.
\end{definition}

The conclusion of \rulename{Eval} is crossed: its left side is the right
side of the second premise. The rule takes no template, no occurrence path,
no counter and no body as a parameter; these are read off the two premises,
which determine the conclusion or the failure of the rule.

\begin{remark}[Counting]
\label{comp:rem:counting}
$\mathsf{E}_3$ has six rules if \rulename{Eval} counts once and eight if its
cases count separately; $\mathsf{E}_2$ has six and seven. The Lean
formalization \Lean{EPrf} has the two cases of $\mathsf{E}_2$ as
constructors (\Lean{evalZero}, \Lean{evalSucc}) and \rulename{Root} as the
derived rule \Lean{EPrf.evalId}. The calculus $\mathsf{U}$ of
\cref{comp:sec:step} has one case and six rules under either convention. The
number of arithmetic and equality schemas, including reflexivity, is four
for $\mathsf{E}_3$, three for $\mathsf{E}_2$ and two for $\mathsf{U}$.
\end{remark}

\subsection{The identity frame}

\begin{lemma}[Identity frame]
\label{comp:lem:frame}
Suppose that a calculus has reflexivity and, for every term $s$, a term
$K(s)$ such that the inference
\[
 \infrule{s=t\qquad K(s)=r}{r=t}
\]
is a rule. Then \rulename{Root} is derivable, by two instances of that rule
and one reflexivity. The equation $K(s)=s$ is not needed as an axiom.
\end{lemma}

\begin{proof}
Let $p$ derive $s=t$ and $q$ derive $s=r$. The rule applied to $q$ and to
the reflexivity $K(s)=K(s)$ gives $K(s)=r$. The rule applied to $p$ and to
$K(s)=r$ gives $r=t$. Each of $p$ and $q$ is used once.
\end{proof}

\begin{theorem}[\rulename{Root} is redundant]
\label{comp:thm:root}
For every context $\Gamma$ and native formula $A$,
\[
 \Gamma\proves_{\mathsf{E}_2}A\quad\Longleftrightarrow\quad
 \Gamma\proves_{\mathsf{E}_3}A .
\]
The translations leave every term and formula of the sequent unchanged. A
derivation of $\mathsf{E}_3$ with $N$ inference nodes, $k$ of which are
instances of \rulename{Root}, is translated to a derivation of
$\mathsf{E}_2$ with exactly $N+2k$ nodes.
\end{theorem}

\begin{proof}
From right to left, apply \cref{comp:lem:frame} with
$K(s)\defeq\It{0}{s}{u}{0}$: its rule is \rulename{Eval}$_0$ with the body
$0$. Explicitly, an instance of \rulename{Root} with premises $p\colon s=t$
and $q\colon s=r$ is replaced by
\[
 \mathrm{RC}(p,q)\defeq
 \mathrm{Ev}_0\bigl(p,\ \mathrm{Ev}_0(q,\ \mathrm{R}_{K(s)})\bigr),
\]
where $\mathrm{R}_{w}$ is the reflexivity leaf $w=w$ and $\mathrm{Ev}_0$ is
\rulename{Eval}$_0$ with its premises in the displayed order. The
replacement is carried out recursively; every other inference is copied.
Each replaced node becomes two instances of \rulename{Eval}$_0$ and one
reflexivity, and its premises occur once, which gives the count. From left
to right the translation is the inclusion.
\end{proof}

The count is of inference nodes. The added reflexivity leaf contains the
term $K(s)$, whose size depends on $s$, so the theorem bounds neither the
size of a serialized derivation nor the time of checking it.

\subsection{Computation, equality and transport recovered}
\label{comp:sec:evalrecover}

The derivations of this subsection are valid in $\mathsf{E}_2$ over an
arbitrary context and use only \rulename{Refl}, \rulename{Eval}$_0$ and
\rulename{Eval}$_{\Suc}$, apart from the given premises. No binder rule, no
numerical case distinction and no induction is used. Write
$\mathrm{Ev}(p,q)$ for \rulename{Eval} with first premise $p$ and second
premise $q$.

\begin{lemma}[Equality rules]
\label{comp:lem:evalalg}
Symmetry and transitivity are derivable in $\mathsf{E}_2$: for
$p\colon s=t$ put $\mathrm{Sym}(p)\defeq\mathrm{RC}(\mathrm{R}_s,p)$, which
derives $t=s$; for $p\colon s=t$ and $q\colon t=v$ put
$\mathrm{Trans}(p,q)\defeq\mathrm{RC}(q,\mathrm{Sym}(p))$, which derives
$s=v$.
\end{lemma}

\begin{proof}
$\mathrm{RC}$ derives $r=t$ from $s=t$ and $s=r$
(\cref{comp:thm:root}). With $s=s$ and $s=t$ it derives $t=s$. With $t=v$
and $t=s$ it derives $s=v$.
\end{proof}

\begin{lemma}[Both computation equations]
\label{comp:lem:evalcomp}
For all raw $c,a,h$ the equations
\[
 \It{0}{a}{u}{h}=a,\qquad
 \It{\Suc c}{a}{u}{h}=h\subst{\It{c}{a}{u}{h}}{u}
\]
have derivations in $\mathsf{E}_2$ with three nodes.
\end{lemma}

\begin{proof}
Put $A\defeq\It{0}{a}{u}{h}$, $C\defeq\It{c}{a}{u}{h}$ and
$B\defeq\It{\Suc c}{a}{u}{h}$. Then $\mathrm{Ev}(\mathrm{R}_a,\mathrm{R}_A)$
is an instance of \rulename{Eval}$_0$ and concludes $A=a$, and
$\mathrm{Ev}(\mathrm{R}_C,\mathrm{R}_B)$ is an instance of
\rulename{Eval}$_{\Suc}$ and concludes $B=h\subst{C}{u}$. The second
derivation is the same when $c\equiv0$.
\end{proof}

\begin{theorem}[Congruence in every template]
\label{comp:thm:evalcong}
For every template $h(u)$ the rule
\[
 \infrule{\Gamma\proves s=t}{\Gamma\proves h\subst{s}{u}=h\subst{t}{u}}
\]
is admissible in $\mathsf{E}_2$. The derivation adds eight inference nodes to
the derivation of $s=t$, which occurs once.
\end{theorem}

\begin{proof}
Let $p$ derive $s=t$ and put $A_s\defeq\It{0}{s}{u}{h}$ and
$B_s\defeq\It{1}{s}{u}{h}$. Four instances of \rulename{Eval} give
\begin{align*}
 z_s&\defeq\mathrm{Ev}(\mathrm{R}_s,\mathrm{R}_{A_s}) &&\colon A_s=s,\\
 c_s&\defeq\mathrm{Ev}(z_s,\mathrm{R}_{B_s}) &&\colon B_s=h\subst{s}{u},\\
 z_{s,t}&\defeq\mathrm{Ev}(p,\mathrm{R}_{A_s}) &&\colon A_s=t,\\
 \mathrm{Cong}_h(p)&\defeq\mathrm{Ev}(z_{s,t},c_s)
   &&\colon h\subst{s}{u}=h\subst{t}{u}.
\end{align*}
The first and third lines are instances of \rulename{Eval}$_0$. In the
second and fourth lines the second premise has the left side $B_s$, whose
preceding iterate is literally $A_s$, the left side of the first premise;
these are instances of \rulename{Eval}$_{\Suc}$. The substitutions into $h$
are the ordinary substitutions performed by \rulename{Eval} itself, so no
congruence is assumed inside the derivation. The tree has four reflexivity
leaves, of which $\mathrm{R}_{A_s}$ occurs twice, and four instances of
\rulename{Eval}.
\end{proof}

\begin{corollary}[Transport]
\label{comp:cor:evaltr}
The rule \rulename{Tr}, from $s=t$ and $h\subst{s}{u}=r$ infer
$r=h\subst{t}{u}$, is admissible in $\mathsf{E}_2$ for all templates and all
terms $s$, including the literal $0$. In $\mathsf{E}_3$ the derivation adds
nine inference nodes to the derivations of the premises; in $\mathsf{E}_2$
it adds eleven. Each premise occurs once.
\end{corollary}

\begin{proof}
Apply \rulename{Root} to $\mathrm{Cong}_h(p)$ and the second premise. In
$\mathsf{E}_2$ the instance of \rulename{Root} is the three nodes of
$\mathrm{RC}$.
\end{proof}

\begin{remark}
\pkg{equality\_from\_iteration} derives transport in $\mathsf{E}_2$ by a
different composition: from $C_0\defeq\It{0}{s}{u}{h}$ and
$C_1\defeq\It{1}{s}{u}{h}$ it forms $C_0=t$ and $C_1=r$ by
\cref{comp:lem:evalalg,comp:lem:evalcomp} and concludes with one instance of
\rulename{Eval}$_{\Suc}$. The template $h$ occurs in these derivations as
the body of an iterator; it is not a parameter of a primitive rule. The node
counts say nothing about the size of the term labels, which contain copies
of $h$, $s$ and $t$.
\end{remark}

\subsection{Equivalence with the core}

\begin{theorem}[Equivalence]
\label{comp:thm:eval}
For every finite context $\Gamma$ of native formulas and every native
formula $A$,
\[
 \Gamma\proves_{\mathsf{E}_3}A
 \quad\Longleftrightarrow\quad
 \Gamma\proves_{\mathsf{E}_2}A
 \quad\Longleftrightarrow\quad
 \Gamma\proves_{\Nat}A .
\]
The translations are effective on finite derivations and leave every term,
formula and assumption of the sequent unchanged.
\end{theorem}

\begin{proof}
The first equivalence is \cref{comp:thm:root}. For the second it suffices,
by \cref{kern:thm:witness}, to compare $\mathsf{E}_2$ and $\mathsf{E}_3$
with the calculus $\mathsf{W}$ of \cref{kern:def:W}, which has the rules
\rulename{Hyp}, \rulename{QE}, \rulename{DE} and \rulename{WI} in common
with them.

From $\mathsf{W}$ to $\mathsf{E}_2$: \rulename{I0} and \rulename{IS} are
derived by \cref{comp:lem:evalcomp} and \rulename{Tr} by
\cref{comp:cor:evaltr}; every other inference is copied, and each instance
of \rulename{WI} keeps its witness and its mode.

From $\mathsf{E}_3$ to $\mathsf{W}$: reflexivity and the equality rules are
derivable in $\mathsf{W}$ (\cref{kern:lem:local}(iii)). For an instance of
\rulename{Eval} with premises $s=t$ and $b=r$ distinguish the rows of
$\Phi$. If $b\equiv s$, \rulename{Tr} with the identity template gives
$r=t$. If $b\equiv\It{0}{s}{u}{h}$, then \rulename{I0} gives $b=s$, and the
equality rules give $r=t$ from $b=s$, $s=t$ and $b=r$. If
$b\equiv\It{\Suc c}{a}{u}{h}$ and $s\equiv\It{c}{a}{u}{h}$, then
\rulename{IS} gives $b=h\subst{s}{u}$; with $b=r$ this gives
$h\subst{s}{u}=r$, and \rulename{Tr} with the template $h$ and the major
premise $s=t$ gives $r=h\subst{t}{u}$. Neither translation uses a
structural rule or a substitution of terms in derivations.
\end{proof}

\begin{remark}
This is the proof of \pkg{computational\_equality}, which states the
theorem for $\mathsf{E}_3$ and cites \Lean{wprf\_iff\_nprf} for the step
from $\mathsf{W}$ to $\Nat$; \pkg{equality\_from\_iteration} adds
$\mathsf{E}_2$. The Lean statements are \Lean{eprf\_iff\_nprf} and
\Lean{eprf\_iff\_iprf}. Weakening, injective renaming and cut are
admissible in $\mathsf{E}_3$ and $\mathsf{E}_2$ also by the direct
induction of \cref{kern:lem:struct}: the operation $\Phi$ is a match of
literal shapes, which an injective renaming preserves and a cut does not
touch.
\end{remark}

\begin{corollary}[Linear bound on inference nodes]
\label{comp:cor:evalnodes}
A derivation with $N$ inference nodes in $\mathsf{W}$ is translated to a
derivation of $\mathsf{E}_3$ with at most $9N$ nodes.
\end{corollary}

\begin{proof}
Each leaf \rulename{I0} or \rulename{IS} becomes three nodes
(\cref{comp:lem:evalcomp}). Each instance of \rulename{Tr} becomes the
translations of its premises, each once, and nine further nodes
(\cref{comp:cor:evaltr}). Every other node is copied. Induction on the
derivation gives the bound.
\end{proof}

The bound holds for this translation. It is not a bound on the size of
formulas or of serialized derivations, and it is not a statement of
optimality.

\subsection{Substitution}

\begin{proposition}[\rulename{Eval} is stable under substitution]
\label{comp:prop:phisubst}
If $\Phi(s,t,b)$ is defined, then for every capture-free substitution
$\sigma$ the term $\Phi(s[\sigma],t[\sigma],b[\sigma])$ is defined by the
same row and
\[
 \Phi(s[\sigma],t[\sigma],b[\sigma])\equiv\Phi(s,t,b)[\sigma].
\]
\end{proposition}

\begin{proof}
The first row is immediate. In the second row substitution keeps the outer
iterator with the literal counter $0$, and its initializer is the
substituted source. In the third row
$b[\sigma]\equiv\It{\Suc(c[\sigma])}{a[\sigma]}{u}{h[\sigma]}$ and
$s[\sigma]\equiv\It{c[\sigma]}{a[\sigma]}{u}{h[\sigma]}$, where $\sigma$ acts
on the parameters of $h$ and leaves $u$ bound; the composition law of
ordinary substitution identifies $(h\subst{t}{u})[\sigma]$ with
$h[\sigma]\subst{t[\sigma]}{u}$. No computation takes place: if $c[\sigma]$
is the literal $0$, the preceding iterate is still the raw term
$\It{0}{a[\sigma]}{u}{h[\sigma]}$. The size comparison after
\cref{comp:def:phi} excludes a collapse into the first row.
\end{proof}

The proposition concerns \rulename{Refl} and \rulename{Eval}. The calculus
$\mathsf{E}_3$ keeps \rulename{WI}, so \cref{comp:rem:single} applies:
substitution of terms in derivations of $\mathsf{E}_3$, the Leibniz rule,
numerical cases and equation induction are admissible by
\cref{comp:thm:eval} and \cref{kern:prop:transfer}, and
\rulename{Eval} adds no repair of its own.

\subsection{Irredundancy}

\begin{proposition}[Neither remaining case can be deleted]
\label{comp:prop:evalsharp}
The calculus $\mathsf{E}_2$ without \rulename{Eval}$_0$ does not derive
$\It{0}{0}{u}{0}=0$. The calculus $\mathsf{E}_2$ without
\rulename{Eval}$_{\Suc}$ does not derive
$\It{1}{0}{u}{\Suc u}=\Suc(\It{0}{0}{u}{\Suc u})$.
\end{proposition}

\begin{proof}
Interpret the formulas in the standard natural numbers, with the standard
$0$ and $\Suc$, equality as equality and the readings of
\cref{base:sec:reading} for $\Q{y}$ and $\D{x}{y}$. Under any interpretation
of the iterator as a function of its counter, its initializer and its
parameters, \rulename{Hyp}, \rulename{Refl}, \rulename{QE} and \rulename{DE}
are sound; \rulename{WI} is sound in progress mode by induction on the
standard numbers and in case mode by the distinction between $0$ and
successors.

Without \rulename{Eval}$_0$: interpret $\It{n}{a}{u}{h}$ as the $n$-fold
iterate of $h$ applied to $a+1$. Then $\It{\Suc c}{a}{u}{h}$ is $h$ applied
to $\It{c}{a}{u}{h}$, so \rulename{Eval}$_{\Suc}$ is sound, while
$\It{0}{0}{u}{0}$ has the value $1$.

Without \rulename{Eval}$_{\Suc}$: interpret $\It{n}{a}{u}{h}$ as $a$.
\rulename{Eval}$_0$ is sound, and the displayed successor equation reads
$0=1$.
\end{proof}

\begin{proposition}[Reflexivity cannot be deleted]
\label{comp:prop:reflsharp}
In $\mathsf{E}_2$ without \rulename{Refl}, no literal equation is derivable
from a context that contains no literal equation. In particular no equation
is derivable from the empty context.
\end{proposition}

\begin{proof}
By induction on a derivation whose context contains no literal equation and
whose conclusion is one. \rulename{Hyp} cannot be the last rule. Either case
of \rulename{Eval} has a literal equation as a premise in the same context.
An instance of \rulename{QE} whose conclusion $B(v)$ is a literal equation
has a matrix that is a literal equation, so its second premise $B(0)$ is a
literal equation in the same context. \rulename{DE} adds an assumption with
root $\Q{y}$, so the context of its second premise still contains no literal
equation and its conclusion is the same equation. \rulename{WI} concludes a
formula with root $\D{x}{y}$. Each case contradicts the induction hypothesis.
\end{proof}

The two propositions concern the deletion of these rules from this
inventory. They exclude neither another term language nor another rule
basis.

%%%%%%%%%%%%%%%%%%%%%%%%%%%%%%%%%%%%%%%%%%%%%%%%%%%%%%%%%%%%%%%%%%%%%%%%%%%%
\section{Normalizers and the transfer of derivations}
\label{comp:sec:normalizers}

This section states once the comparison argument that
\pkg{zero\_normal}, \pkg{initializer\_equality},
\pkg{full\_normalization} and \pkg{computation\_normal} each carry out for
their own normalizer, and on which \pkg{equality\_from\_iteration},
\pkg{uniform\_step}, \pkg{seed\_evaluation} and \pkg{reflexive\_base} rely by
citation. The argument uses five properties of the normalizer and nothing
else.

\subsection{Admissible normalizers and normalizing substitution}

\begin{definition}[Admissible normalizer]
\label{comp:def:normalizer}
An \emph{admissible normalizer} is a map $\nu$ from raw terms to raw terms
such that for all raw terms $t$, all renamings $\rho$ of variables and all
capture-free simultaneous substitutions $\sigma$:
\begin{enumerate}[label=($\nu$\arabic*)]
\item $\nu(x)\equiv x$, $\nu(0)\equiv0$ and
 $\nu(\Suc t)\equiv\Suc\nu(t)$;
\item $\nu(t[\rho])\equiv\nu(t)[\rho]$ and $\FV(\nu(t))\subseteq\FV(t)$;
\item $\nu(\nu(t))\equiv\nu(t)$;
\item $\nu(t[\sigma])\equiv\nu\bigl(\nu(t)[\nu\circ\sigma]\bigr)$;
\item $\proves_{\Nat}t=\nu(t)$, with the free variables of $t$ as
 parameters.
\end{enumerate}
A term $t$ is \emph{$\nu$-normal} if $\nu(t)\equiv t$. The map $\nu$ is
extended to formulas by $\nu(s=t)\defeq(\nu(s)=\nu(t))$,
$\nu(\Q{y}B)\defeq\Q{y}\nu(B)$ and $\nu(\D{x}{y}B)\defeq\D{x}{y}\nu(B)$, and
to contexts entry by entry. A formula is $\nu$-normal if $\nu(A)\equiv A$.
\end{definition}

Condition ($\nu$4) is the \emph{coherence of normalization with
substitution}: the normal form of an instance depends only on the normal
forms of the term and of the substituted terms. Condition ($\nu$5) says that
$\nu$ is the identity inside $\Nat$; it is a statement about each fixed term
and does not say that $\Nat$ or $\iPRA$ proves a uniform theorem about
$\nu$. By ($\nu$2) the extension to formulas does not depend on the names of
bound variables, and $\nu$ never changes the root of a formula; in
particular $\nu(A)$ is a literal equation exactly when $A$ is.

\begin{definition}[Normalizing substitution]
\label{comp:def:nusubst}
For a raw term or formula $T$ and a substitution $\sigma$,
\[
 T\langle\sigma\rangle\defeq\nu(T[\sigma]).
\]
For a template $h$ and a matrix $B$ of $\nu$-normal syntax, the instances
$h(s)$, $B(r)$ and $B(t,r)$ are formed with this substitution.
\end{definition}

\begin{lemma}[Laws of normalizing substitution]
\label{comp:lem:laws}
Let $\nu$ be an admissible normalizer, $T$ a raw term or formula, and
$\sigma,\tau$ substitutions.
\begin{enumerate}[label=(\alph*)]
\item $\nu(\nu(T)[\sigma])\equiv\nu(T[\sigma])$.
\item $(T\langle\sigma\rangle)\langle\tau\rangle\equiv
 T\langle x\mapsto\sigma(x)\langle\tau\rangle\rangle\equiv
 \nu(T[\sigma][\tau])$.
\item If $T$ is $\nu$-normal, then $T\langle\rho\rangle\equiv T[\rho]$ for
 every renaming $\rho$; in particular the identity substitution leaves $T$
 unchanged.
\end{enumerate}
In particular instances compose: for a matrix $B(x,y)$ of $\nu$-normal
syntax, the instance at $t$ of the template $B(v,r)$ is $B(t,r)$.
\end{lemma}

\begin{proof}
For terms, (a) follows from ($\nu$4) applied to $\nu(T)$, from ($\nu$3) and
from ($\nu$4) applied to $T$:
$\nu(\nu(T)[\sigma])\equiv\nu(\nu(\nu(T))[\nu\circ\sigma])\equiv
\nu(\nu(T)[\nu\circ\sigma])\equiv\nu(T[\sigma])$. For (b), by (a),
$(T\langle\sigma\rangle)\langle\tau\rangle\equiv\nu(\nu(T[\sigma])[\tau])
\equiv\nu(T[\sigma][\tau])$. Let $\theta(x)\defeq\sigma(x)[\tau]$, so that
$T[\sigma][\tau]\equiv T[\theta]$ by the composition law of ordinary
substitution. By ($\nu$4), ($\nu$3) and ($\nu$4) again,
$\nu(T[\nu\circ\theta])\equiv\nu(\nu(T)[\nu\circ\nu\circ\theta])\equiv
\nu(\nu(T)[\nu\circ\theta])\equiv\nu(T[\theta])$, and the left side is
$T\langle x\mapsto\sigma(x)\langle\tau\rangle\rangle$. (c) is ($\nu$2) with
($\nu$3). For formulas use induction on the formula: below a binder the
substitution is lifted, and the lifting of $\nu\circ\sigma$ is $\nu$ composed
with the lifting of $\sigma$ by ($\nu$1) and ($\nu$2).
\end{proof}

\begin{lemma}[A normalizer is the identity in the core]
\label{comp:lem:nusound}
Let $\nu$ be an admissible normalizer. For every raw formula $A$ and every
context $\Gamma$, both $\Gamma,A\proves_{\Nat}\nu(A)$ and
$\Gamma,\nu(A)\proves_{\Nat}A$. For every raw term $t$ and substitution
$\sigma$, $\proves_{\Nat}t[\sigma]=t\langle\sigma\rangle$.
\end{lemma}

\begin{proof}
The term statement is ($\nu$5) for $t[\sigma]$. For formulas the induction
on $A$ in the proof of \cref{comp:lem:conversion} applies, with ($\nu$5) at
the two sides of a literal equation in place of the equation of a reduct.
\end{proof}

The inclusion of $\nu$-normal syntax into raw syntax therefore commutes with
substitution up to provable equality: the term $t\langle\sigma\rangle$ is
by definition the normal form of $t[\sigma]$, the two terms need not be
identical, and their equation is a theorem of $\Nat$.

\subsection{The kernel on normal forms}

\begin{definition}[The kernel $\Nat^{\nu}$]
\label{comp:def:nukernel}
Let $\nu$ be an admissible normalizer. The calculus $\Nat^{\nu}$ has the
$\nu$-normal native formulas as formulas, finite lists of them as contexts,
the substitution of \cref{comp:def:nusubst}, and the following rules, in
which every term is $\nu$-normal:
\begin{itemize}
\item \rulename{Hyp}, \rulename{Tr}, \rulename{QE}, \rulename{DE},
 \rulename{DI} and \rulename{QI}, read with normalizing substitution;
\item \rulename{Comp}$^{\nu}$: $\Gamma\proves\nu(e)$ for every instance $e$
 of \rulename{I0} or of \rulename{IS} on raw terms.
\end{itemize}
\end{definition}

\begin{lemma}
\label{comp:lem:nustruct}
Weakening, cut on an assumption and normalizing substitution are admissible
in $\Nat^{\nu}$: if $\Gamma\proves_{\Nat^{\nu}}A$, then
$\Gamma\langle\sigma\rangle\proves_{\Nat^{\nu}}A\langle\sigma\rangle$ for
every $\sigma$.
\end{lemma}

\begin{proof}
By induction on derivations, with eigenvariables renamed away from $\sigma$.
No rule of $\Nat^{\nu}$ tests a term or a formula for a literal shape other
than the root of the matrix in \rulename{QI}, which substitution preserves.
By \cref{comp:lem:laws}(b) the instance under $\sigma$ of a premise or
conclusion of \rulename{Tr}, \rulename{QE}, \rulename{DE}, \rulename{DI} or
\rulename{QI} is the corresponding premise or conclusion of the same rule at
the substituted parameters. For \rulename{Comp}$^{\nu}$,
$\nu(e)\langle\sigma\rangle\equiv\nu(e[\sigma])$ by \cref{comp:lem:laws}(a),
and $e[\sigma]$ is an instance of the same computation equation. Weakening
and cut are standard.
\end{proof}

\begin{theorem}[Transfer]
\label{comp:thm:transfer}
Let $\nu$ be an admissible normalizer. For every finite context $\Gamma$ of
raw native formulas and every raw native formula $A$,
\[
 \Gamma\proves_{\Nat}A
 \quad\Longleftrightarrow\quad
 \nu(\Gamma)\proves_{\Nat^{\nu}}\nu(A).
\]
Both directions are effective on finite derivations.
\end{theorem}

\begin{proof}
From left to right, by induction on the derivation in $\Nat$, normalizing
every formula. An assumption of $\Gamma$ becomes an assumption of
$\nu(\Gamma)$. A leaf \rulename{I0} or \rulename{IS} with conclusion $e$
becomes \rulename{Comp}$^{\nu}$ with conclusion $\nu(e)$. For \rulename{Tr}
with premises $s=t$ and $h[s/v]=r$, the induction hypotheses are
$\nu(s)=\nu(t)$ and $\nu(h[s/v])=\nu(r)$. By ($\nu$4),
$\nu(h[s/v])\equiv\nu(h)\langle\nu(s)/v\rangle$, so \rulename{Tr} of
$\Nat^{\nu}$ with the template $\nu(h)$ gives
$\nu(r)=\nu(h)\langle\nu(t)/v\rangle$, which is $\nu(r)=\nu(h[t/v])$ by
($\nu$4). In the same way ($\nu$4), with ($\nu$1) for the terms $0$ and
$\Suc z$, turns the premises and the conclusion of \rulename{QE},
\rulename{DI} and \rulename{QI} into those of the same rule at the matrix
$\nu(B)$ and the parameters $\nu(t)$; the matrix of \rulename{QI} remains a
literal equation. For \rulename{DE} the local assumption $\Q{y}B[b/x]$
becomes $\Q{y}\nu(B)[b/x]$ by ($\nu$2). Freshness conditions are preserved,
because $\nu$ introduces no free variable by ($\nu$2).

From right to left, read a derivation of $\Delta\proves B$ in $\Nat^{\nu}$
as raw syntax and show $\Delta\proves_{\Nat}B$ by induction. A leaf
\rulename{Comp}$^{\nu}$ with conclusion $\nu(e)$: the equation $e$ is an
axiom of $\Nat$, and \cref{comp:lem:nusound} gives $\nu(e)$. For
\rulename{Tr} with premises $s=t$ and $h\langle s/v\rangle=r$:
\cref{comp:lem:nusound} gives $h[s/v]=h\langle s/v\rangle$, hence
$h[s/v]=r$; \rulename{Tr} of $\Nat$ gives $r=h[t/v]$, and
\cref{comp:lem:nusound} gives $r=h\langle t/v\rangle$. For \rulename{QE}
convert the premise $B\langle0/y\rangle$ to $B[0/y]$ by
\cref{comp:lem:nusound}, apply \rulename{QE} of $\Nat$, and convert
$B[t/y]$ to $B\langle t/y\rangle$. For \rulename{DE} the local assumption
is unchanged, since $B\langle b/x\rangle\equiv B[b/x]$ by
\cref{comp:lem:laws}(c). For \rulename{DI} work in the context
$\Delta,B[t/x,0/y]$: convert the assumption to
$B\langle t/x,0/y\rangle$, cut with the translated premise, convert the
result to $B[t/x,\Suc z/y]$, and apply \rulename{DI} of $\Nat$.
\rulename{QI} is treated in the same way. The conversions are uniform in the
parameters, so no eigenvariable enters a context or a conclusion. This gives
$\nu(\Gamma)\proves_{\Nat}\nu(A)$ from
$\nu(\Gamma)\proves_{\Nat^{\nu}}\nu(A)$; \cref{comp:lem:nusound} and cut,
applied to each assumption and to the conclusion, give
$\Gamma\proves_{\Nat}A$.
\end{proof}

A derivation with the single introduction must not be normalized rule by
rule: a witness that is not literally $0$ can have the normal form $0$, and
the mode of \rulename{WI} then changes. The theorem is therefore stated for
the kernel with \rulename{DI} and \rulename{QI}, and each calculus with
\rulename{WI} on normal forms is compared with $\Nat^{\nu}$ afterwards.

\begin{corollary}[Calculi equivalent to the kernel on normal forms]
\label{comp:cor:bridge}
Let $\nu$ be an admissible normalizer and let $\mathsf{K}$ be a calculus on
$\nu$-normal native formulas that derives exactly the sequents of
$\Nat^{\nu}$. Let $J$ and $H$ be the translations of formulas from $\Nat$ to
$\iPRA$ and from $\iPRA$ to $\Nat$ of \cref{base:thm:bridge}, and put
$H_{\nu}\defeq\nu\circ H$.
\begin{enumerate}[label=(\alph*)]
\item $\Gamma\proves_{\Nat}A$ if and only if
 $\nu(\Gamma)\proves_{\mathsf K}\nu(A)$, for raw $\Gamma$ and $A$. For
 $\nu$-normal $\Delta$ and $B$, $\Delta\proves_{\mathsf K}B$ if and only if
 $\Delta\proves_{\Nat}B$.
\item Weakening, cut, normalizing substitution and the Leibniz rule, from
 $s=t$ and $F(s)$ infer $F(t)$ for every matrix $F$, are admissible in
 $\mathsf{K}$.
\item For $\nu$-normal $\Delta,B$ and formulas $\Theta,\varphi$ of $\iPRA$,
 \begin{align*}
  \Delta\proves_{\mathsf K}B
   &\quad\Longleftrightarrow\quad J\Delta\proves_{\iPRA}JB,\\
  \Theta\proves_{\iPRA}\varphi
   &\quad\Longleftrightarrow\quad
   H_{\nu}\Theta\proves_{\mathsf K}H_{\nu}\varphi,
 \end{align*}
 the formulas $B$ and $H_{\nu}JB$ are interderivable in $\mathsf{K}$, and
 $\iPRA\proves\varphi\leftrightarrow JH_{\nu}\varphi$. The numerical
 domains are whole, equality is interpreted by equality, and the numerical
 composites have the graph $x=y$. In particular
 $\proves_{\mathsf K}0=\Suc0$ if and only if $\iPRA\proves\bot$.
\end{enumerate}
For a calculus on raw terms that derives exactly the sequents of $\Nat$,
the same statements, with $\nu$ the identity, are
\cref{kern:prop:transfer}.
\end{corollary}

\begin{proof}
(a) The first statement is \cref{comp:thm:transfer}. For the second, apply
it to $\Delta$ and $B$, which are their own normal forms. (b) Weakening, cut
and substitution are \cref{comp:lem:nustruct}. For the Leibniz rule,
$\Nat$ derives $F[t/v]$ from $s=t$ and $F[s/v]$
(\cref{base:thm:leibniz}), and normalizing this sequent gives the rule on
normal forms by ($\nu$4). (c) The first equivalence is (a) composed with
\cref{base:thm:bridge}. For the second, $\Theta\proves_{\iPRA}\varphi$ if
and only if $H\Theta\proves_{\Nat}H\varphi$, if and only if
$\nu(H\Theta)\proves_{\mathsf K}\nu(H\varphi)$ by (a). The formulas $B$ and
$HJB$ are interderivable in $\Nat$; normalizing the two sequents gives the
interderivability of $B$ and $H_{\nu}JB$ in $\mathsf{K}$. By
\cref{comp:lem:nusound}, $H\varphi$ and $\nu(H\varphi)$ are interderivable
in $\Nat$, hence $\iPRA\proves JH\varphi\leftrightarrow JH_{\nu}\varphi$,
and $\iPRA\proves\varphi\leftrightarrow JH\varphi$. The map $\nu$ fixes the
variables and changes neither a domain formula nor the interpretation of
equality: normal forms are representatives of expressions, not a restriction
or a quotient of the domain of numbers. The formula $0=\Suc0$ is
$\nu$-normal by ($\nu$1).
\end{proof}

Every calculus on normal forms in the following sections is compared with
$\Nat^{\nu}$ on the same $\nu$-normal sequents, and \cref{comp:cor:bridge}
is then not repeated. The normalizer and the normalizing substitution are
part of the specification of such a calculus. A checker for it contains
them, and the measurements of \cref{comp:sec:sizes} count them: the deletion
of a rule does not delete the cost of the computation that the rule
expressed.

%%%%%%%%%%%%%%%%%%%%%%%%%%%%%%%%%%%%%%%%%%%%%%%%%%%%%%%%%%%%%%%%%%%%%%%%%%%%
\section{Zero-normal terms}
\label{comp:sec:zero}

This section draws on \pkg{zero\_normal}, which deletes the zero computation
rule by normalizing it in the syntax, and on the restatements of the same
normalization layer in \pkg{initializer\_equality},
\pkg{equality\_from\_iteration}, \pkg{uniform\_step},
\pkg{seed\_evaluation} and \pkg{reflexive\_base}.

\subsection{Zero normalization}

\begin{definition}[Zero normalization]
\label{comp:def:nfzero}
For terms $c,a,h$ the \emph{smart constructor} is
\[
 \overline{I}(c,a;\,u.\,h)\defeq
 \begin{cases}
  a, & c\equiv0,\\
  \It{c}{a}{u}{h}, & c\not\equiv0.
 \end{cases}
\]
The map $\nf_0$ on raw terms is defined by structural recursion:
\begin{align*}
 \nf_0(x)&\defeq x, & \nf_0(0)&\defeq0, & \nf_0(\Suc t)&\defeq\Suc\nf_0(t),\\
 \nf_0(\It{c}{a}{u}{h})&\defeq
 \overline{I}(\nf_0(c),\nf_0(a);\,u.\,\nf_0(h)). &&&&
\end{align*}
A term is \emph{zero-normal} if no iterator in it, at any depth, has the
literal $0$ as its counter.
\end{definition}

The test in the last clause is on the \emph{normalized} counter. For
example $\nf_0\bigl(\It{\It{0}{0}{v}{0}}{a}{u}{h}\bigr)\equiv\nf_0(a)$,
although the outer counter is not literally $0$; the equivalence of
``$\nf_0(t)\equiv0$'' with ``$t\equiv0$'' fails. Only the zero computation
is normalized: $\It{1}{a}{u}{h}$ is zero-normal and is not identified with
$h\subst{a}{u}$. In $\It{c}{a}{u}{h}$ the binder $u$ covers $h$ only, so the
initializer returned by a contraction keeps its scope and no variable index
is adjusted; for instance
$\nf_0\bigl(\It{n}{a}{u}{\It{0}{u}{v}{0}}\bigr)\equiv\It{n}{a}{u}{u}$ for a
variable $n$.

\begin{lemma}
\label{comp:lem:nfzero}
\begin{enumerate}[label=(\alph*)]
\item $\nf_0(t)$ is zero-normal; $\nf_0(t)\equiv t$ if and only if $t$ is
 zero-normal; $\nf_0$ is idempotent.
\item $\nf_0$ commutes with renaming, and $\FV(\nf_0(t))\subseteq\FV(t)$.
\item $\nf_0(t)$ has at most as many nodes as $t$, and $\nf_0$ is a
 primitive recursive operation on codes of finite terms.
\item $t\to^{*}\nf_0(t)$ by zero contractions. Zero contraction, in any
 position and in any order, terminates, and $\nf_0(t)$ is its unique normal
 form.
\end{enumerate}
\end{lemma}

\begin{proof}
(a) By structural induction: the children of a rebuilt iterator are
zero-normal, and the smart constructor returns either the zero-normal
initializer or an iterator whose counter is not $0$. On a zero-normal term
every clause returns its argument. (b) A renaming sends no variable,
successor or iterator to the literal $0$, so it commutes with the test of
the smart constructor; a contraction deletes subterms and creates no
variable. (c) A contraction deletes an iterator node with its counter and
body and keeps the normalized initializer; otherwise the normalized children
are kept below the same constructor. A structural recursion with these size
bounds is primitive recursive. It never performs a successor contraction:
$\It{\Suc n}{a}{u}{h}$ is kept even for a closed numeral $n$, so no
evaluator for arbitrary primitive recursive programs is involved. (d) The
bottom-up computation of $\nf_0$ is a sequence of zero contractions. A zero
contraction decreases the number of nodes, hence terminates. If $t\to t'$
by a zero contraction, then $\nf_0(t)\equiv\nf_0(t')$, by induction on the
position: at the root both sides are the normalized initializer; in a child,
the normal form of the child is unchanged by the induction hypothesis, and
the result is unchanged even if the child is erased. A term without a zero
redex is zero-normal and fixed by $\nf_0$, so every normal form of $t$ is
$\nf_0(t)$.
\end{proof}

\begin{lemma}[Coherence with substitution]
\label{comp:lem:nfzerosubst}
For every raw term $t$ and every capture-free simultaneous substitution
$\sigma$,
\[
 \nf_0(t[\sigma])\equiv\nf_0\bigl(\nf_0(t)[\nf_0\circ\sigma]\bigr).
\]
\end{lemma}

\begin{proof}
By induction on $t$. For a variable both sides are $\nf_0(\sigma(x))$, by
idempotence. The cases $0$ and $\Suc t$ are immediate. Let
$t\equiv\It{c}{a}{u}{h}$ and let $\sigma'$ be the lifting of $\sigma$ below
$u$; by \cref{comp:lem:nfzero}(b) the lifting of $\nf_0\circ\sigma$ is
$\nf_0\circ\sigma'$. The left side is
$\overline{I}(\nf_0(c[\sigma]),\nf_0(a[\sigma]);u.\nf_0(h[\sigma']))$. If
$\nf_0(c)\equiv0$, then $\nf_0(t)\equiv\nf_0(a)$ and the right side is
$\nf_0(\nf_0(a)[\nf_0\circ\sigma])\equiv\nf_0(a[\sigma])$ by the induction
hypothesis for $a$; the induction hypothesis for $c$ gives
$\nf_0(c[\sigma])\equiv\nf_0(0[\nf_0\circ\sigma])\equiv0$, so the left side
is $\nf_0(a[\sigma])$ too. If $\nf_0(c)\not\equiv0$, then
$\nf_0(t)\equiv\It{\nf_0(c)}{\nf_0(a)}{u}{\nf_0(h)}$ and the right side is
the smart constructor applied to
$\nf_0(\nf_0(c)[\nf_0\circ\sigma])$, $\nf_0(\nf_0(a)[\nf_0\circ\sigma])$ and
$\nf_0(\nf_0(h)[\nf_0\circ\sigma'])$, which by the induction hypotheses are
the three arguments of the smart constructor on the left side. The
substituted counter may or may not be $0$; both sides apply the same test
to the same term.
\end{proof}

\begin{proposition}
\label{comp:prop:nfzeroadm}
$\nf_0$ is an admissible normalizer, and the $\nf_0$-normal terms are the
zero-normal terms.
\end{proposition}

\begin{proof}
($\nu$1) holds by definition, ($\nu$2) and ($\nu$3) by
\cref{comp:lem:nfzero}, ($\nu$4) by \cref{comp:lem:nfzerosubst}, and
($\nu$5) by \cref{comp:lem:nfzero}(d) and \cref{comp:lem:conversion}.
\end{proof}

In the rest of this section and in
\cref{comp:sec:reflexive,comp:sec:seed,comp:sec:step,comp:sec:seedeval},
terms and formulas are zero-normal and $t\langle\sigma\rangle$ is
$\nf_0(t[\sigma])$. The substitution is computed by rebuilding the term
bottom-up with the smart constructor: in an iterator the substitution acts
on the counter and the initializer in the current scope and, lifted, on the
body. A substitution into a zero-normal term can create a counter $0$, which
is then contracted. Neither the normalizer nor any rule tests whether a term
\emph{denotes} $0$.

\subsection{The kernel without the zero computation rule}

\begin{definition}[The kernel $\Nat^{0}$]
\label{comp:def:nzero}
The calculus $\Nat^{0}$ has zero-normal native formulas, finite lists of
them as contexts, normalizing substitution, and the seven rules
\rulename{Hyp}, \rulename{Tr}, \rulename{QE}, \rulename{DE}, \rulename{DI},
\rulename{QI} and
\[
 \infrule{}{\Gamma\proves\It{\Suc c}{a}{u}{h}
            =h\langle\overline{I}(c,a;\,u.\,h)/u\rangle}
 \qquad\text{(\rulename{IS})}
\]
for all zero-normal $c,a,h$. The left side is zero-normal because $\Suc c$
is not the literal $0$; at $c\equiv0$ the right side is
$h\langle a/u\rangle$.
\end{definition}

\begin{lemma}
\label{comp:lem:nzerorefl}
$\Nat^{0}$ derives $r=r$ for every zero-normal $r$, and $\Nat^{0}$ derives
exactly the sequents of $\Nat^{\nf_0}$.
\end{lemma}

\begin{proof}
Put $J_r\defeq\It{1}{r}{u}{u}$. \rulename{IS} gives $J_r=r$, because the
right side is $u\langle\overline{I}(0,r;u.u)/u\rangle\equiv r$, and
\cref{comp:lem:eqalg}(a) gives $r=r$. For the second statement compare
\rulename{IS} with \rulename{Comp}$^{\nf_0}$. If $e$ is the instance
$\It{0}{a}{u}{h}=a$ of \rulename{I0}, then $\nf_0(e)$ is
$\nf_0(a)=\nf_0(a)$, which $\Nat^{0}$ derives. If $e$ is an instance of
\rulename{IS} at raw $c,a,h$, then by \cref{comp:lem:nfzerosubst}
$\nf_0(e)$ is the instance of \rulename{IS} of $\Nat^{0}$ at
$\nf_0(c),\nf_0(a),\nf_0(h)$. Conversely an instance of \rulename{IS} of
$\Nat^{0}$ at zero-normal $c,a,h$ is $\nf_0$ of the raw instance at the same
terms. The other rules are common.
\end{proof}

\begin{theorem}[The zero computation rule is deleted]
\label{comp:thm:zero}
For every finite context $\Gamma$ of raw native formulas and every raw
native formula $A$,
\[
 \Gamma\proves_{\Nat}A
 \quad\Longleftrightarrow\quad
 \nf_0(\Gamma)\proves_{\Nat^{0}}\nf_0(A).
\]
$\Nat^{0}$ is closed under weakening, cut and normalizing substitution, and
\cref{comp:cor:bridge} applies to it.
\end{theorem}

\begin{proof}
\Cref{comp:prop:nfzeroadm}, \cref{comp:thm:transfer} and
\cref{comp:lem:nzerorefl,comp:lem:nustruct}.
\end{proof}

\begin{remark}
The theorem is the faithfulness corollary of \pkg{zero\_normal} for its
seven-rule calculus and the raw/canonical equivalence of
\pkg{initializer\_equality}, whose comparison calculus lists reflexivity as
an eighth rule; \pkg{equality\_from\_iteration}, \pkg{uniform\_step},
\pkg{seed\_evaluation} and \pkg{reflexive\_base} cite it. The normalization
layer of \pkg{zero\_normal} is the Lean module \Lean{ZeroNormalization},
with \Lean{nf\_subst}, \Lean{hsub\_comp}, \Lean{nf\_correct} and
\Lean{inclusion\_subst\_coherent}; the theorem is \Lean{nzprf\_iff\_nprf}
for the calculus \Lean{NZPrf}, with \Lean{NZPrf.normal} and
\Lean{nzprf\_iff\_pra}. The Lean calculus admits raw rule parameters and
proves as a theorem that it derives only zero-normal formulas from
zero-normal assumptions; here this holds by the definition of the syntax.
Definitional zero computation is proposed earlier in the package
\pkg{prevalidated\_elimination}, where it does not reduce the checker
because the arity coding places the initializer below a dummy binder and a
contraction must lower indices, and the package \pkg{uniform\_zero} asks
whether another layout of binders makes it cheaper; \pkg{zero\_normal}
carries it out for the present layout, in which the initializer lies outside
the binder.
\end{remark}

\subsection{Six rules with the single introduction}

\begin{lemma}[Zero-normal aliases of zero]
\label{comp:lem:alias}
The closed terms $\It{1}{0}{u}{0}$ and $\It{1}{0}{u}{u}$ are zero-normal and
are not literally $0$, and \rulename{IS} of $\Nat^{0}$ derives
$\It{1}{0}{u}{0}=0$ and $\It{1}{0}{u}{u}=0$.
\end{lemma}

\begin{proof}
The counter $\Suc0$ is not $0$, so neither term is contracted. The right
sides of the two instances of \rulename{IS} are
$0\langle\overline{I}(0,0;u.0)/u\rangle\equiv0$ and
$u\langle\overline{I}(0,0;u.u)/u\rangle\equiv0$.
\end{proof}

The raw alias $\It{0}{0}{u}{0}$ of \cref{comp:thm:eval} is not available: it
normalizes to $0$ and selects the progress mode. The packages use
$\It{1}{0}{u}{0}$ (\pkg{zero\_normal}, \pkg{initializer\_equality},
\pkg{equality\_from\_iteration}) or $\It{1}{0}{u}{u}$ (\pkg{uniform\_step},
\pkg{seed\_evaluation}); the only property used is that the term is provably
$0$ and not literally $0$.

\begin{definition}[The calculus $\mathsf{Z}_6$]
\label{comp:def:zsix}
$\mathsf{Z}_6$ has the syntax and the substitution of $\Nat^{0}$ and the six
rules \rulename{Hyp}, \rulename{IS}, \rulename{Tr}, \rulename{QE},
\rulename{DE} and \rulename{WI}; the test $t\equiv0$ of \rulename{WI} is
applied to the zero-normal witness as supplied. There is no primitive
reflexivity, symmetry, substitution, cut, induction, coverage or separation
rule.
\end{definition}

\begin{theorem}
\label{comp:thm:zsix}
$\mathsf{Z}_6$ and $\Nat^{0}$ derive the same sequents. Hence
$\Gamma\proves_{\Nat}A$ if and only if
$\nf_0(\Gamma)\proves_{\mathsf{Z}_6}\nf_0(A)$.
\end{theorem}

\begin{proof}
Weakening, injective renaming and cut are admissible in $\mathsf{Z}_6$ by
direct induction on derivations: they preserve the root of a matrix and the
literal-zero status of a witness, and cut substitutes no term. Reflexivity
is derived as in \cref{comp:lem:nzerorefl}, and \cref{comp:lem:eqalg}(b)
gives symmetry and transport in literal equations; the instances match by
\cref{comp:lem:laws}. By \cref{comp:lem:single}(i), \rulename{QI} is
admissible in $\mathsf{Z}_6$; by \cref{comp:lem:single}(ii) with the alias
$Z\defeq\It{1}{0}{u}{0}$ of \cref{comp:lem:alias}, so is the one instance of
\rulename{DI} that is not an instance of \rulename{WI}. Conversely
\rulename{WI} is admissible in $\Nat^{0}$ by \cref{comp:lem:single}(iii).
The last statement is \cref{comp:thm:zero}.
\end{proof}

\begin{corollary}[Equation induction]
\label{comp:cor:zsixind}
For a literal equation $E$ the rule that infers $\Gamma\proves E(t)$ from
$\Gamma\proves E(0)$ and $\Gamma,E(z)\proves E(\Suc z)$, with $z$ fresh, is
admissible in $\mathsf{Z}_6$.
\end{corollary}

\begin{proof}
\rulename{QI}, admissible by \cref{comp:thm:zsix}, and \rulename{QE}.
\end{proof}

\begin{remark}
\pkg{zero\_normal} states \cref{comp:thm:zsix} with a seven-rule bridge, the
calculus $\Nat^{0}$, and warns that a derivation of the calculus
$\mathsf{W}$ of \cref{kern:def:W} on raw terms must not be normalized with
its modes kept. The six-rule calculus
is not formalized in Lean in this form: the formalized six-rule calculus on
zero-normal terms is \Lean{NRPrf} of \cref{comp:sec:reflexive}, with a
different mode test, and the calculus $\mathsf{U}$ of \cref{comp:sec:step},
with the present test.
\end{remark}

%%%%%%%%%%%%%%%%%%%%%%%%%%%%%%%%%%%%%%%%%%%%%%%%%%%%%%%%%%%%%%%%%%%%%%%%%%%%
\section{The introduction selected by a reflexive base}
\label{comp:sec:reflexive}

This section draws on \pkg{reflexive\_base}. The single introduction is
kept, and its mode is selected by the formula that the case mode would
assume, in place of the pair of tests on the root of the matrix and on the
witness. No alias of zero is needed.

\subsection{The calculus}

\begin{definition}[The calculus $\mathsf{R}$]
\label{comp:def:reflexive}
For a zero-normal matrix $B(x,y)$ and a zero-normal term $t$, let
$\mathsf{r}(B,t)$ hold if the instance $B(t,0)$ is literally $r=r$ for some
term $r$, and put
\[
 \eta^{\mathsf r}_{B,t}(z)\defeq
 \begin{cases}
  z, & \mathsf{r}(B,t),\\
  0, & \text{otherwise}.
 \end{cases}
\]
The rule \rulename{RI} is
\[
 \infrule{\Gamma,\,B(t,\eta^{\mathsf r}_{B,t}(z))\proves B(t,\Suc z)}
         {\Gamma\proves\D{x}{y}B(x,y)}
 \qquad(z\text{ fresh for }\Gamma,\ t\text{ and the parameters of }B).
\]
The calculus $\mathsf{R}$ has the syntax and the substitution of $\Nat^{0}$
and the six rules \rulename{Hyp}, \rulename{IS}, \rulename{Tr},
\rulename{QE}, \rulename{DE} and \rulename{RI}.
\end{definition}

The predicate $\mathsf{r}$ compares two finite trees. Since the instance
$B(t,0)$ has the root of $B$, the condition $\mathsf{r}(B,t)$ implies that
$B$ is a literal equation. The instance $B(t,0)$ is normalized, so the
predicate depends on the zero normalization of the calculus.

\subsection{The structural rules, without circularity}

The lemmas are proved in the following order: laws of normalizing
substitution; weakening and cut; equality; unrestricted introduction by
cases; substitution. None of them presupposes the equivalence theorem.

\begin{lemma}
\label{comp:lem:rstruct}
Weakening, injective renaming of free variables and cut on an assumption are
admissible in $\mathsf{R}$. Reflexivity, symmetry, transitivity, congruence
and transport in literal equations are derivable in $\mathsf{R}$ from
\rulename{IS} and \rulename{Tr}.
\end{lemma}

\begin{proof}
The first three operations leave the matrix and the witness of each
instance of \rulename{RI} unchanged or apply an injective renaming to them.
A reflexive equation stays reflexive, and an equation that is not reflexive
does not become reflexive under an injective renaming, so no mode changes.
The equality rules are those of \cref{comp:lem:nzerorefl,comp:lem:eqalg};
their derivations use no introduction rule.
\end{proof}

In particular every base selected by $\mathsf{r}$ is derivable in every
context.

\begin{lemma}[Every introduction by cases is admissible]
\label{comp:lem:rcases}
For every zero-normal matrix $B$ and witness $t$, if
$\Gamma,B(t,0)\proves_{\mathsf R}B(t,\Suc z)$ with $z$ fresh, then
$\Gamma\proves_{\mathsf R}\D{x}{y}B$.
\end{lemma}

\begin{proof}
If $\mathsf{r}(B,t)$ fails, this is \rulename{RI}. If it holds, the
assumption $B(t,0)$ is literally $r=r$; derive it by reflexivity and remove
it by cut. Weaken the resulting derivation of $B(t,\Suc z)$ by the
assumption $B(t,z)$ and apply \rulename{RI} in progress mode.
\end{proof}

\begin{lemma}[Substitution]
\label{comp:lem:rsubst}
Normalizing substitution is admissible in $\mathsf{R}$.
\end{lemma}

\begin{proof}
If $B(t,0)$ is literally $r=r$, then for every $\sigma$ its instance is
$r\langle\sigma\rangle=r\langle\sigma\rangle$, and by \cref{comp:lem:laws}
this is the zero instance of the substituted matrix at the substituted
witness. Hence a true guard remains true. Now argue by induction on a
derivation. \rulename{IS}, \rulename{Tr}, \rulename{QE} and \rulename{DE}
commute with normalizing substitution by \cref{comp:lem:laws}, with fresh
eigenvariables. For an instance of \rulename{RI} in progress mode the guard
remains true and the induction hypothesis is the premise of \rulename{RI} in
progress mode at the substituted data. For an instance in case mode the
induction hypothesis is the substituted case premise, and
\cref{comp:lem:rcases} applies whether or not the guard has become true.
\end{proof}

The guard is preserved by substitution and not reflected: if the zero
instance of an equation is $p=q$ for distinct variables $p,q$, the guard
fails, and it holds after the substitution of $0$ for both. A second way in
which a false guard becomes true is the contraction of a counter that the
substitution turns into $0$. The proof repairs exactly this change of mode.

\subsection{Progress recovered by two decoders}

By \cref{comp:lem:rcases}, \rulename{DE} and the construction of
\cref{comp:lem:single}(i), the closure by cases is admissible in
$\mathsf{R}$:
\begin{equation}
\label{comp:eq:rqcase}
 \infrule{\Lambda,\,F(0)\proves F(\Suc z)}{\Lambda\proves\Q{y}F(y)}
 \qquad(z\text{ fresh for }\Lambda\text{ and the parameters of }F).
\end{equation}
Recall the selector $\sel(v,a,b)=\It{v}{a}{u}{b}$, with $u$ not in $b$. In
zero-normal syntax the instance $\sel(0,a,b)$ is $a$, and \rulename{IS}
gives $\sel(\Suc v,a,b)=b$.

\begin{lemma}[Case decoder]
\label{comp:lem:decoder}
Let $r(y)$ be a zero-normal template and $a$ a term without $y$, and let
$\mathrm{Dec}_{r,a}(y)$ be the equation
\[
 \sel\bigl(y,\,a,\,\sel(y,0,r(y))\bigr)=r(y).
\]
Then $\proves_{\mathsf R}\Q{y}\mathrm{Dec}_{r,a}(y)$. Consequently
$\Lambda\proves_{\mathsf R}a=r(0)$ implies
$\Lambda\proves_{\mathsf R}\mathrm{Dec}_{r,a}(v)$ for every term $v$.
\end{lemma}

\begin{proof}
The zero instance of $\mathrm{Dec}_{r,a}$ is literally $a=r(0)$. At $y=\Suc z$ two
instances of \rulename{IS} and transitivity give
\[
 \sel\bigl(\Suc z,a,\sel(\Suc z,0,r(\Suc z))\bigr)
 =\sel(\Suc z,0,r(\Suc z))=r(\Suc z),
\]
a derivation that does not use the assumption of the base. Rule
\eqref{comp:eq:rqcase} gives $\Q{y}\mathrm{Dec}_{r,a}(y)$ in the empty context.
Weaken into $\Lambda$ and apply \rulename{QE} with the given base and the
target $v$.
\end{proof}

The lemma is a numerical case distinction in a fixed context: $\Lambda$ does
not contain the generalized variable, and the conclusion can be instantiated
at a variable that occurs in $\Lambda$, because \rulename{QE} imposes no
freshness condition on its target.

\begin{theorem}[Equation progress]
\label{comp:thm:rprogress}
Let $E(y)$ be the zero-normal literal equation $s(y)=t(y)$. If
$\Gamma,E(z)\proves_{\mathsf R}E(\Suc z)$ with $z$ fresh for $\Gamma$ and
the parameters of $E$, then $\Gamma\proves_{\mathsf R}\Q{y}E(y)$.
\end{theorem}

\begin{proof}
Let $\pi$ be the given derivation, write $s_0\defeq s(0)$ and
$t_0\defeq t(0)$, and keep the base as an assumption:
$\Delta\defeq\Gamma,E(0)$, where $E(0)$ is $s_0=t_0$. The \emph{progressive
lift} of $E$ is
\[
 H(y)\defeq\bigl(\sel(y,0,s(y))=\sel(y,0,t(y))\bigr).
\]
Its zero instance is literally $0=0$. Consider the two decoders
\[
 \mathrm{Dec}_{s,s_0}(y)\colon\ \sel\bigl(y,s_0,\sel(y,0,s(y))\bigr)=s(y),
 \qquad
 \mathrm{Dec}_{t,s_0}(y)\colon\ \sel\bigl(y,s_0,\sel(y,0,t(y))\bigr)=t(y).
\]
The zero instance of the first is $s_0=s_0$, derivable by reflexivity, and
that of the second is $s_0=t_0$, the assumption $E(0)$. By
\cref{comp:lem:decoder}, $\Delta\proves\mathrm{Dec}_{s,s_0}(z)$ and
$\Delta\proves\mathrm{Dec}_{t,s_0}(z)$ at the fresh variable $z$.

Add the assumption $H(z)$. Congruence in the template
$w\mapsto\sel(z,s_0,w)$ turns it into
\[
 \sel\bigl(z,s_0,\sel(z,0,s(z))\bigr)=\sel\bigl(z,s_0,\sel(z,0,t(z))\bigr).
\]
The symmetric form of $\mathrm{Dec}_{s,s_0}(z)$, this equation and
$\mathrm{Dec}_{t,s_0}(z)$ give $s(z)=t(z)$, which is $E(z)$. Weaken $\pi$ into the
context $\Delta,H(z)$ and cut $E(z)$ into its assumption; the result is
$\Delta,H(z)\proves E(\Suc z)$. Nothing is substituted into $\pi$: its
variable is still $z$. \rulename{IS} on each side of $H(\Suc z)$ converts
$E(\Suc z)$ to $H(\Suc z)$, so
\[
 \Delta,\,H(z)\proves H(\Suc z).
\]
Regard $H(y)$ as a matrix with an unused witness coordinate. Its zero
instance at any witness is $0=0$, so \rulename{RI} is in progress mode and
gives $\Delta\proves\D{x}{y}H(y)$; \rulename{DE} returns
$\Delta\proves\Q{y}H(y)$, since the witness is unused. Reflexivity derives
$H(0)$, \rulename{QE} at $\Suc z$ gives $H(\Suc z)$, and \rulename{IS}
gives $E(\Suc z)$. This is a derivation of
$\Gamma,E(0)\proves E(\Suc z)$, and rule \eqref{comp:eq:rqcase} gives
$\Gamma\proves\Q{y}E(y)$.
\end{proof}

The proof uses equality, the closure by cases, weakening, renaming and cut.
It uses no implication between formulas, no substitution of terms into the
given derivation and not \cref{comp:lem:rsubst}.

\subsection{Equivalence}

\begin{theorem}
\label{comp:thm:reflexive}
For every zero-normal context $\Gamma$ and formula $A$,
\[
 \Gamma\proves_{\mathsf R}A
 \quad\Longleftrightarrow\quad
 \Gamma\proves_{\Nat^{0}}A .
\]
Both translations are effective and are the identity on the sequent. Hence
$\Gamma\proves_{\Nat}A$ if and only if
$\nf_0(\Gamma)\proves_{\mathsf R}\nf_0(A)$ for raw $\Gamma,A$.
\end{theorem}

\begin{proof}
From $\mathsf{R}$ to $\Nat^{0}$: five rules are common. An instance of
\rulename{RI} in case mode is an instance of \rulename{DI}. In progress mode
the matrix is a literal equation; put $E(y)\defeq B(t,y)$. \rulename{QI}
gives $\Q{y}E(y)$; in the context $\Gamma,E(0)$, \rulename{QE} gives
$E(\Suc z)$; \rulename{DI} at the witness $t$ concludes $\D{x}{y}B$. From
$\Nat^{0}$ to $\mathsf{R}$: \rulename{DI} is admissible by
\cref{comp:lem:rcases} and \rulename{QI} by \cref{comp:thm:rprogress}. The
last statement is \cref{comp:thm:zero}.
\end{proof}

\begin{remark}
The Lean formalization is \Lean{NRPrf}, with \Lean{nrprf\_iff\_nprf},
\Lean{NRPrf.normal} and \Lean{nrprf\_iff\_pra}; its completeness proof
translates the calculus of the package \pkg{diagonal\_guard}, with the
invariant that a literally reflexive assumption may be dropped. The method
of a guarded introduction is that of the package \pkg{diagonal\_guard}; the
guard on a computed zero instance, the proof of
\cref{comp:lem:rsubst} for a guard that can change from false to true, and
\cref{comp:prop:rawreflexive} are specific to \pkg{reflexive\_base}. The
rule that keeps both $B(t,0)$ and, at a literal equation, $B(t,z)$ as
assumptions of one introduction is a different, older rule.
\end{remark}

\subsection{The same guard on raw terms, and two variants}

\begin{proposition}[Seven rules without normalization]
\label{comp:prop:rawreflexive}
Let $\mathsf{R}^{\mathrm{raw}}$ be the calculus on raw terms with ordinary
substitution and the rules \rulename{Hyp}, \rulename{I0}, \rulename{IS},
\rulename{Tr}, \rulename{QE}, \rulename{DE} and \rulename{RI}, where
$\mathsf{r}(B,t)$ tests the literal syntax of the raw instance $B(t,0)$.
Then $\mathsf{R}^{\mathrm{raw}}$ derives exactly the sequents of $\Nat$.
\end{proposition}

\begin{proof}
Reflexivity and the equality rules are derived from \rulename{I0} and
\rulename{Tr} as in \cref{base:lem:eq}. Weakening, injective renaming and cut are direct
inductions; \cref{comp:lem:rcases} and rule \eqref{comp:eq:rqcase} are
obtained as before, and a true guard is preserved by ordinary substitution,
so \cref{comp:lem:rsubst} holds with the same proof. The translation from
$\mathsf{R}^{\mathrm{raw}}$ to $\Nat$ is that of \cref{comp:thm:reflexive}.

Only the lift in the proof of progress changes, because the raw zero
instance of $H$ is not reflexive. For $E(y)\equiv(s(y)=t(y))$ use the
asymmetric equation
\[
 K(y)\defeq\bigl(\sel(y,t(y),s(y))=\sel(0,t(y),s(y))\bigr).
\]
Its raw zero instance is literally
$\sel(0,t(0),s(0))=\sel(0,t(0),s(0))$, so it satisfies the guard before any
computation. For every $v$, \rulename{IS} computes the left side of
$K(\Suc v)$ to $s(\Suc v)$ and \rulename{I0} computes its right side to
$t(\Suc v)$, so $K(\Suc v)$ and $E(\Suc v)$ are interderivable. In the
context $\Gamma,E(0)$ the two decoder equations are
\[
 \sel\bigl(y,s(0),\sel(y,t(y),s(y))\bigr)=s(y),\qquad
 \sel\bigl(y,s(0),\sel(0,t(y),s(y))\bigr)=t(y).
\]
Their zero instances follow from \rulename{I0} and, respectively,
reflexivity and the assumption $s(0)=t(0)$; their successor instances
follow from \rulename{IS} for the outer selector and \rulename{IS},
respectively \rulename{I0}, for the inner one. The closure by cases and
\rulename{QE} give both decoders at every term. Congruence of $K(z)$ in the
template $w\mapsto\sel(z,s(0),w)$ and the decoders give $E(z)$. Cut it into
the given derivation of $E(\Suc z)$, convert $E(\Suc z)$ to $K(\Suc z)$, and
conclude with \rulename{RI} in progress mode, \rulename{DE}, \rulename{QE}
and the closure by cases as in \cref{comp:thm:rprogress}.
\end{proof}

The asymmetric lift also serves on zero-normal terms, where its base
contracts to $t(0)=t(0)$. The guard is thus independent of the
normalization of the zero computation. The Lean formalization of
$\mathsf{R}^{\mathrm{raw}}$ is \Lean{RRPrf}, with \Lean{rrprf\_iff\_nprf};
its structural rules are weakening, injective renaming and cut.

\begin{proposition}[Two variants of the guard]
\label{comp:prop:rvariants}
\Cref{comp:thm:reflexive} remains true if $\mathsf{r}(B,t)$ is replaced by
either of the following predicates:
\begin{enumerate}[label=(\alph*)]
\item $B(t,0)$ is literally $0=0$;
\item $B(x,0)$, with $x$ still a variable, is literally $r=r$.
\end{enumerate}
\end{proposition}

\begin{proof}
The proofs of this section use three properties of the guard: a selected
base is derivable; a true guard is preserved by substitution; and the lift
$H$ of \cref{comp:thm:rprogress} satisfies the guard. For (a) the base is
$0=0$, which is preserved by substitution, and the zero instance of $H$ is
$0=0$. For (b), if $B(x,0)$ is $r=r$ then every $B(t,0)$ is reflexive, the
guard is preserved by the substitution of the parameters, and the zero
instance of $H$ with an unused coordinate $x$ is $0=0$.
\end{proof}

Variant (a) compares with one fixed formula and recognizes fewer matrices.
Variant (b) removes the witness from the test and costs a checker one more
normalizing substitution, since the instance $B(t,0)$ is needed for the
assumption of the case mode in any event. The proof of
\cref{comp:prop:rawreflexive} does not establish variant (a) on raw terms,
where the base of the asymmetric lift is reflexive and in general not
$0=0$.

%%%%%%%%%%%%%%%%%%%%%%%%%%%%%%%%%%%%%%%%%%%%%%%%%%%%%%%%%%%%%%%%%%%%%%%%%%%%
\section{Equality as replacement of an initializer}
\label{comp:sec:seed}

This section draws on \pkg{initializer\_equality}. Successor computation is
kept as a rule, and transport in an arbitrary template is replaced by a rule
that changes only the initializer of an iterator. Because an iterator with
the counter $0$ is its initializer, the same rule replaces at the root.

\subsection{The rule Seed and the calculus}

\begin{definition}[Seed replacement]
\label{comp:def:rho}
For zero-normal terms $s,t,b$ the partial operation $\rho_{s,t}(b)$ is
\[
 \rho_{s,t}(b)\defeq
 \begin{cases}
  t, & b\equiv s,\\
  \It{n}{t}{u}{h}, & b\equiv\It{n}{s}{u}{h},\\
  \text{undefined}, & \text{otherwise}.
 \end{cases}
\]
\end{definition}

The two branches are disjoint, since a finite term is not its own proper
subterm. The second branch inspects one constructor and its initializer and
leaves the counter and the body unchanged; its result is zero-normal because
the counter is the same term, which is not $0$. No bound variable is
substituted or renamed.

\begin{definition}[The calculus $\mathsf{S}$]
\label{comp:def:seedcalc}
$\mathsf{S}$ has the syntax and the substitution of $\Nat^{0}$ and the six
rules \rulename{Hyp}, \rulename{IS} (as in \cref{comp:def:nzero}),
\rulename{QE}, \rulename{DE}, \rulename{WI} and
\[
 \infrule{\Gamma\proves s=t\qquad\Gamma\proves b=r}
         {\Gamma\proves\rho_{s,t}(b)=r}
 \quad(\rho_{s,t}(b)\text{ defined})
 \qquad\text{(\rulename{Seed})}.
\]
\end{definition}

The conclusion of \rulename{Seed} is not crossed. The rule has no parameter:
the two premises determine the conclusion or the failure of the rule.
Substitution is not removed from the calculus, since \rulename{IS} and the
binder rules use it.

\begin{lemma}[Schematic form]
\label{comp:lem:seedschema}
\rulename{Seed} is exactly the schema
\[
 \infrule{\Gamma\proves s=t\qquad
          \Gamma\proves\overline{I}(n,s;\,u.\,h)=r}
         {\Gamma\proves\overline{I}(n,t;\,u.\,h)=r}
\]
for arbitrary zero-normal $n$ and $h$. Its instance at $n\equiv0$ is the
rule of \emph{root replacement}: from $s=t$ and $s=r$ infer $t=r$.
\end{lemma}

\begin{proof}
At $n\equiv0$ the schema is the first branch of $\rho$; the erased body is
irrelevant. At any other $n$ the second premise displays the iterator from
which the counter and the body are read, and the schema is the second
branch. Conversely the first branch is the schema at $n\equiv0$, $h\equiv0$,
and the second branch is the schema at the counter and body of $b$.
\end{proof}

\begin{lemma}[\rulename{Seed} is stable under substitution]
\label{comp:lem:seedsubst}
If $\rho_{s,t}(b)$ is defined, then for every substitution $\sigma$ the term
$\rho_{s\langle\sigma\rangle,t\langle\sigma\rangle}(b\langle\sigma\rangle)$
is defined and equals $\rho_{s,t}(b)\langle\sigma\rangle$.
\end{lemma}

\begin{proof}
In the root branch this is immediate. Let $b\equiv\It{n}{s}{u}{h}$. Then
$b\langle\sigma\rangle$ is the smart constructor applied to the substituted
counter, initializer and body. If the substituted counter is $0$, both
$b\langle\sigma\rangle$ and $\rho_{s,t}(b)\langle\sigma\rangle$ lose their
iterator node and become $s\langle\sigma\rangle$ and
$t\langle\sigma\rangle$, and the rule applies in its root branch. Otherwise
both keep the node, with the same counter and body, and the rule applies in
its second branch. The rule is stable although normalization can change the
branch.
\end{proof}

\subsection{Equality recovered}

The derivations of this subsection hold in $\mathsf{S}$ over an arbitrary
context and use only \rulename{IS} and \rulename{Seed}, apart from the given
premises.

\begin{lemma}[Root equality]
\label{comp:lem:seedalg}
Reflexivity, symmetry and transitivity are derivable in $\mathsf{S}$.
\end{lemma}

\begin{proof}
\rulename{IS} gives $\It{1}{a}{u}{u}=a$. Root replacement with this
equation as both premises gives $a=a$. For $s=t$, root replacement with the
second premise $s=s$ gives $t=s$. For $s=t$ and $t=r$, symmetry gives
$t=s$, and root replacement with the premises $t=s$ and $t=r$ gives $s=r$.
\end{proof}

\begin{lemma}[One-step application]
\label{comp:lem:beta}
For every zero-normal template $f(u)$ and term $a$, \rulename{IS} derives
$\It{1}{a}{u}{f}=f\langle a/u\rangle$.
\end{lemma}

\begin{proof}
The right side of \rulename{IS} at the counter $0$ is
$f\langle\overline{I}(0,a;u.f)/u\rangle\equiv f\langle a/u\rangle$. This is
the point at which the definitional zero computation is used.
\end{proof}

\begin{theorem}[Congruence in every template]
\label{comp:thm:seedcong}
For every zero-normal template $f(u)$ the rule that infers
$f\langle s/u\rangle=f\langle t/u\rangle$ from $s=t$ is admissible in
$\mathsf{S}$. If the derivation of $s=t$ has $P$ inference nodes, the
derivation of the conclusion has $P+4$ nodes, independently of the depth and
the number of occurrences of $u$ in $f$.
\end{theorem}

\begin{proof}
Write $f(a)$ for $f\langle a/u\rangle$ and put
$A_a\defeq\It{1}{a}{u}{f}$. The derivation is
\[
 \infrule{\infrule{s=t\qquad A_s=f(s)}{A_t=f(s)}\qquad A_t=f(t)}
         {f(s)=f(t)}
\]
The equations $A_s=f(s)$ and $A_t=f(t)$ are instances of \rulename{IS}
(\cref{comp:lem:beta}). The upper inference is \rulename{Seed} in its second
branch, which replaces the initializer $s$ of $A_s$ by $t$. The lower
inference is root replacement. The derivation adds two instances of
\rulename{IS} and two of \rulename{Seed}; it uses neither reflexivity nor
symmetry nor transitivity. The occurrences of $u$ in $f$, including those
below iterator binders, are handled by the substitution of \rulename{IS}.
\end{proof}

\begin{corollary}[Transport]
\label{comp:cor:seedtr}
The following rules are admissible in $\mathsf{S}$, for all terms and
contexts: from $s=t$ and $f(s)=r$ infer $f(t)=r$; the rule \rulename{Tr};
and transport in a literal equation, from $s=t$ and $a(s)=b(s)$ infer
$a(t)=b(t)$.
\end{corollary}

\begin{proof}
\Cref{comp:thm:seedcong} gives $f(s)=f(t)$, and root replacement with the
second premise $f(s)=r$ gives $f(t)=r$. Symmetry gives \rulename{Tr}. Two
uses of \rulename{Tr}, as in \cref{comp:lem:eqalg}(b), give transport in an
equation.
\end{proof}

\begin{remark}
The node count is for the displayed tree; its labels contain copies of $f$,
$s$ and $t$, so it bounds neither the size of a serialized derivation nor
the time of checking. The package \pkg{witness\_rule} has a congruence rule
restricted to the counter position of a different recursor, with a separate
root rule and a construction with $P+5$ nodes; here the lifting and the
first root replacement are one inference, and the instance of the same rule
at the counter $0$ is the root replacement. The two term languages differ,
so the comparison is not a uniform shortening of derivations.
\end{remark}

\subsection{Equivalence}

\begin{theorem}
\label{comp:thm:seed}
For every zero-normal context $\Gamma$ and formula $A$,
\[
 \Gamma\proves_{\mathsf S}A
 \quad\Longleftrightarrow\quad
 \Gamma\proves_{\Nat^{0}}A .
\]
Hence $\Gamma\proves_{\Nat}A$ if and only if
$\nf_0(\Gamma)\proves_{\mathsf S}\nf_0(A)$ for raw $\Gamma,A$.
\end{theorem}

\begin{proof}
Weakening, injective renaming and cut are admissible in $\mathsf{S}$ by
direct induction: they leave the terms of each rule instance unchanged up
to an injective renaming, so neither the mode of \rulename{WI} nor the
branch of \rulename{Seed} changes.

From $\Nat^{0}$ to $\mathsf{S}$: \rulename{Tr} is admissible by
\cref{comp:cor:seedtr}, which with \cref{comp:lem:seedalg} also supplies the
symmetry and the transport in literal equations required by
\cref{comp:lem:single}. By \cref{comp:lem:single}(i), \rulename{QI} is
admissible, and by \cref{comp:lem:single}(ii) with the alias
$Z\defeq\It{1}{0}{u}{0}$, for which \rulename{IS} derives $Z=0$, so is the
instance of \rulename{DI} that \rulename{WI} lacks. The other rules are
common.

From $\mathsf{S}$ to $\Nat^{0}$: by \cref{comp:lem:seedschema} an instance
of \rulename{Seed} has the premises $s=t$ and
$\overline{I}(n,s;u.h)=r$. \rulename{Tr} of $\Nat^{0}$ in the template
$\overline{I}(n,v;u.h)$, which is $v$ if $n\equiv0$ and $\It{n}{v}{u}{h}$
otherwise, gives $r=\overline{I}(n,t;u.h)$, and symmetry gives the
conclusion. \rulename{WI} is admissible in $\Nat^{0}$ by
\cref{comp:lem:single}(iii). The last statement is \cref{comp:thm:zero}.
\end{proof}

By \cref{comp:cor:bridge}(b), substitution and the Leibniz rule are
admissible in $\mathsf{S}$. \pkg{initializer\_equality} also proves the
Leibniz rule directly in $\mathsf{S}$, by the induction on the matrix of
\cref{base:thm:leibniz} with the binder steps of \cref{base:lem:bindercong}:
at $\Q{y}B$ the closure by cases of \cref{comp:lem:single}(i) replaces
\cref{base:lem:qI}, and at $\D{x}{y}B$ the target is introduced by
\rulename{WI} at the variable opened by \rulename{DE}, which is in case
mode. Only the literal-equation case uses \cref{comp:cor:seedtr}.

\begin{remark}
The Lean formalization of the rule is \Lean{NIPrf}, with
\Lean{niprf\_iff\_nprf}; it keeps \rulename{DI} and \rulename{QI} and
therefore has seven rules. A successor-seeded canonical recursor, as in the
package \pkg{canonical\_recursors}, erases the body of a recursor at the
counter $0$ and keeps the node; it does not make the zero step the identity
of the initializer, and the root replacement is then not an instance of
replacement of an initializer.
\end{remark}

%%%%%%%%%%%%%%%%%%%%%%%%%%%%%%%%%%%%%%%%%%%%%%%%%%%%%%%%%%%%%%%%%%%%%%%%%%%%
\section{One successor step}
\label{comp:sec:step}

This section draws on \pkg{uniform\_step}, which presents the rule
\rulename{Step} with clause contexts, and on the companion calculus of
\pkg{equality\_from\_iteration}, where the same rule is called UE and the
contexts are lists of native formulas. On zero-normal terms the case
\rulename{Eval}$_0$ of \cref{comp:def:e3} is absorbed by the syntax, and
reflexivity with the one remaining case gives all computation and all
equality reasoning.

\subsection{The rule Step}

\begin{definition}[The rule \rulename{Step} and the calculus $\mathsf{U}$]
\label{comp:def:step}
On zero-normal terms the rule \rulename{Step} is
\[
 \infrule{\Gamma\proves\overline{I}(c,a;\,u.\,h)=t\qquad
          \Gamma\proves\It{\Suc c}{a}{u}{h}=r}
         {\Gamma\proves r=h\langle t/u\rangle}.
\]
The calculus $\mathsf{U}$ has the syntax and the substitution of $\Nat^{0}$
and the six rules \rulename{Hyp}, \rulename{Refl}, \rulename{Step},
\rulename{QE}, \rulename{DE} and \rulename{WI}.
\end{definition}

The second premise shows $c$, $a$ and $h$: its left side is an iterator
whose counter is a successor. The left side of the first premise is compared
with the predecessor iterate built by the smart constructor; at $c\equiv0$
it is $a$. The rule has one case, takes no template and compares no term
with an arbitrary context. On raw terms the same rule is
\rulename{Eval}$_{\Suc}$.

\begin{lemma}[Computation and equality from \rulename{Refl} and
\rulename{Step}]
\label{comp:lem:stepderived}
The following are derivable from \rulename{Refl} and \rulename{Step} over
an arbitrary context.
\begin{enumerate}[label=(\alph*)]
\item Successor computation: $\It{\Suc c}{a}{u}{h}
 =h\langle\overline{I}(c,a;u.h)/u\rangle$, by one \rulename{Step} on two
 reflexivities.
\item The pivot: from $s=t$ and $s=r$ infer $r=t$, by two instances of
 \rulename{Step} and one reflexivity.
\item Symmetry and transitivity.
\item Congruence: from $s=t$ infer $h\langle s/v\rangle=h\langle
 t/v\rangle$, by two instances of \rulename{Step} and two reflexivities.
\item \rulename{Tr}: from $s=t$ and $h\langle s/v\rangle=r$ infer
 $r=h\langle t/v\rangle$, by four instances of \rulename{Step} and three
 reflexivities; and transport in a literal equation.
\end{enumerate}
The derivations (b)--(e) use \rulename{Step} at the counter $c\equiv0$ only.
\end{lemma}

\begin{proof}
(a) Apply \rulename{Step} to the reflexivities of $\overline{I}(c,a;u.h)$
and of $\It{\Suc c}{a}{u}{h}$.

(b) Put $J_s\defeq\It{1}{s}{v}{v}$; its predecessor iterate is
$\overline{I}(0,s;v.v)\equiv s$. \rulename{Step} with the premises $s=r$ and
$J_s=J_s$ gives $J_s=r$. \rulename{Step} with the premises $s=t$ and $J_s=r$
gives $r=t$:
\[
 \infrule{s=t\qquad\infrule{s=r\qquad J_s=J_s}{J_s=r}}{r=t}
\]
This is \cref{comp:lem:frame} with the frame $K(s)\defeq J_s$, which is not
normalized away. The derivation is valid for every $s$, including $0$.

(c) The pivot applied to $s=s$ and $s=t$ gives $t=s$. From $s=t$ and $t=r$,
symmetry gives $t=s$, and the pivot applied to $t=r$ and $t=s$ gives $s=r$.

(d) Put $H_s\defeq\It{1}{s}{v}{h}$. \rulename{Step} with the premises $s=s$
and $H_s=H_s$ gives $H_s=h\langle s/v\rangle$. \rulename{Step} with the
premises $s=t$ and $H_s=h\langle s/v\rangle$ gives
$h\langle s/v\rangle=h\langle t/v\rangle$:
\[
 \infrule{s=t\qquad
          \infrule{s=s\qquad H_s=H_s}{H_s=h\langle s/v\rangle}}
         {h\langle s/v\rangle=h\langle t/v\rangle}
\]
The first result is reused as the second premise; it is not necessary to
derive both $H_s=h\langle s/v\rangle$ and $H_s=h\langle t/v\rangle$ and to
pivot them.

(e) The pivot applied to the conclusion of (d) and to
$h\langle s/v\rangle=r$ gives $r=h\langle t/v\rangle$. Transport in a
literal equation follows by \cref{comp:lem:eqalg}(b).
\end{proof}

\begin{remark}
The short congruence (d) is that of \pkg{seed\_evaluation} and of the
package \pkg{clausal\_single\_evaluation}. \pkg{uniform\_step} and
\pkg{equality\_from\_iteration} derive both $H_s=h\langle s/v\rangle$ and
$H_s=h\langle t/v\rangle$ and pivot them; their transport has six instances
of \rulename{Step} and five reflexivities. None of the derivations uses
equation induction, a numerical case distinction, the injectivity of the
successor, a binder rule or the substitution of terms in derivations.
\end{remark}

\begin{proposition}[\rulename{Step} is stable under substitution]
\label{comp:prop:stepsubst}
The instance under a normalizing substitution $\sigma$ of the premises and
the conclusion of an instance of \rulename{Step} are the premises and the
conclusion of the instance of \rulename{Step} at the substituted counter,
initializer and body.
\end{proposition}

\begin{proof}
The counter of the second premise stays a successor. For zero-normal
$c,a,h$ the term $\overline{I}(c,a;u.h)$ is $\nf_0(\It{c}{a}{u}{h})$, so by
\cref{comp:lem:laws}(a) its instance under $\sigma$ is the smart
constructor applied to the substituted counter, initializer and body,
including when the substituted counter is $0$. By \cref{comp:lem:laws}(b)
the instance of $h\langle t/u\rangle$ is the instance of the substituted
body at $t\langle\sigma\rangle$.
\end{proof}

\begin{theorem}
\label{comp:thm:step}
For every zero-normal context $\Gamma$ and formula $A$,
\[
 \Gamma\proves_{\mathsf U}A
 \quad\Longleftrightarrow\quad
 \Gamma\proves_{\Nat^{0}}A .
\]
Hence $\Gamma\proves_{\Nat}A$ if and only if
$\nf_0(\Gamma)\proves_{\mathsf U}\nf_0(A)$ for raw $\Gamma,A$. The
equivalence holds also for the calculus with the seven rules \rulename{Hyp},
\rulename{Refl}, \rulename{Step}, \rulename{QE}, \rulename{DE},
\rulename{DI} and \rulename{QI}.
\end{theorem}

\begin{proof}
Weakening, injective renaming and cut are admissible in $\mathsf{U}$ by
direct induction, as in \cref{comp:thm:zsix}. From $\Nat^{0}$ to
$\mathsf{U}$: \rulename{IS} and \rulename{Tr} are derived by
\cref{comp:lem:stepderived}. The term $Z\defeq\It{1}{0}{u}{0}$ is not
literally $0$, and \rulename{Step} on the reflexivities of $0$ and of $Z$
derives $Z=0$; \cref{comp:lem:single}(i) and (ii) give \rulename{QI} and
\rulename{DI}. From $\mathsf{U}$ to $\Nat^{0}$: reflexivity is
\cref{comp:lem:nzerorefl}. For an instance of \rulename{Step},
\rulename{IS} gives
$\It{\Suc c}{a}{u}{h}=h\langle\overline{I}(c,a;u.h)/u\rangle$; congruence
along the first premise gives
$h\langle\overline{I}(c,a;u.h)/u\rangle=h\langle t/u\rangle$, and with the
second premise the equality rules give $r=h\langle t/u\rangle$.
\rulename{WI} is admissible by \cref{comp:lem:single}(iii). For the
seven-rule calculus the arithmetic rules are translated in the same way and
the introductions are common.
\end{proof}

\begin{remark}
$\mathsf{U}$ is the companion calculus of \pkg{equality\_from\_iteration},
formalized as \Lean{ZEPrf} with \Lean{zeprf\_iff\_nprf}; the missing
introduction by cases is derived there with the alias $\It{1}{0}{u}{0}$,
with no weakening and no cut. The seven-rule calculus is the formalization
\Lean{NSPrf} of the arithmetic interface of \pkg{uniform\_step}, with
\Lean{nsprf\_iff\_nprf}, \Lean{NSPrf.normal} and \Lean{nsprf\_iff\_pra}. The
theorem is an equivalence of consequence relations under arbitrary
assumptions; it is more than the agreement of the two interfaces on true
equations between numerals.
\end{remark}

\subsection{Five patterns with clause contexts}
\label{comp:sec:clausal}

\pkg{uniform\_step} combines \rulename{Step} with the clause contexts of
\cref{base:sec:clauses}: the formulas are the equations and the formulas
$\D{x}{y}B$, a context is a finite list of clauses, and an assumption of a
formula $A$ is the constant clause $\operatorname{ord}(A)$ of
\cref{claus:sec:origins}.

\begin{definition}[The calculus $\mathsf{C}_{\mathrm{Step}}$]
\label{comp:def:cstep}
$\mathsf{C}_{\mathrm{Step}}$ is the calculus $\mathsf{C}^{0}$ of
\cref{claus:def:c0}: zero-normal terms, normalizing substitution, clause
contexts, and the five rules \rulename{Use}, \rulename{Refl}, the rule
\rulename{Step} of \cref{comp:def:step} (called \rulename{Eval} there), the
elimination \rulename{DE} that retains the clause $[B(b,0);\,y.B(b,y)]$
(called \rulename{Elim} there), and \rulename{WI} with its assumption
stored as a constant clause (called \rulename{Intro} there).
\end{definition}

\begin{theorem}[Faithfulness of the five patterns]
\label{comp:thm:clausal}
For every finite context $\Gamma$ of raw native formulas and every raw
native formula $A$,
\[
 \Gamma\proves_{\Nat}A
 \quad\Longleftrightarrow\quad
 \operatorname{ord}\bigl(\nf_0(\tau\Gamma)\bigr)
 \proves_{\mathsf{C}_{\mathrm{Step}}}\nf_0(\tau A),
\]
where $\tau$ is the translation of \cref{claus:def:tau}, which replaces each
$\Q{y}B$ by $\D{w}{y}B$ with an unused $w$. The ordinary realization of
$\mathsf{C}_{\mathrm{Step}}$ is bi-interpretable and equiconsistent with
$\iPRA$, with whole numerical domains, equality interpreted by equality and
numerical composites with the graph $x=y$.
\end{theorem}

\begin{proof}
\Cref{claus:thm:zero}(iv), with \cref{claus:cor:composites,claus:cor:bridge}
for the composites and the bridge.
\end{proof}

The derived rules of the calculus are those of \cref{claus:lem:c0}; its
items (ii)--(v) are the derivations of \cref{comp:lem:stepderived}.

\begin{remark}
\label{comp:rem:clausal}
The calculus of \pkg{uniform\_step} is $\mathsf{C}^{0}$; the package differs
from \cref{claus:sec:zero} in its proof of \cref{comp:thm:clausal}, in three
respects. First, from $\Nat$ it translates the derivations of $\Nat^{0}$
(\cref{comp:thm:zero}) rule by rule, where \cref{claus:thm:zero}(i)
normalizes derivations of the raw clause calculus: \rulename{Hyp} is
\rulename{Use} applied to the reflexivity of $0$, \rulename{IS} and
\rulename{Tr} are \cref{comp:lem:stepderived}, \rulename{QE} and
\rulename{QI} are the closure interfaces of \cref{claus:lem:closure}, and
for \rulename{DE} the retained clause derives the translated closure
(\cref{claus:lem:closure}(iii)) and a cut supplies the branch. Second, its
alias of zero is $q\defeq\It{1}{0}{u}{u}$, and it recovers the introduction
by cases of a literal equation $E(x,y)$ at the witness $0$ without a cut:
the given premise is closed at the unused witness $q$ into $\D{w}{y}E(0,y)$,
the elimination retains the clause $[E(0,0);\,y.E(0,y)]$, and under the
assumption $E(q,0)$ a transport along $q=0$, a \rulename{Use} of that clause
and a transport along $0=q$ give $E(q,\Suc z)$, the premise of \rulename{WI}
at $q$ in case mode. Third, toward $\Nat$ it argues by the soundness of the
five rules in $\Teq$ under the ordinary reading, recorded after
\cref{claus:cor:bridge}, and by \cref{base:thm:bridge}, in place of the
translation of annotated contexts in \cref{claus:thm:zero}(ii). The five
patterns keep the choice between cases and progress inside \rulename{WI};
the distinction between coverage and equation induction is not removed. The
Lean formalization of the package is the arithmetic interface with native
contexts (\Lean{NSPrf}, \cref{comp:thm:step}).
\end{remark}

%%%%%%%%%%%%%%%%%%%%%%%%%%%%%%%%%%%%%%%%%%%%%%%%%%%%%%%%%%%%%%%%%%%%%%%%%%%%
\section{Seed evaluation and suspended application}
\label{comp:sec:seedeval}

This section draws on \pkg{seed\_evaluation}: its rule \rulename{SeedEval},
which fuses the replacement of an initializer (\cref{comp:sec:seed}) with
the successor step (\cref{comp:sec:step}), and its second construction, in
which every iteration at a successor counter is factored into a suspended
one-step application.

\subsection{The rule SeedEval}

\begin{definition}
\label{comp:def:seedeval}
On zero-normal terms the rule \rulename{SeedEval} is
\[
 \infrule{\Gamma\proves a=b\qquad\Gamma\proves\It{\Suc c}{a}{u}{h}=r}
         {\Gamma\proves r=h\langle\overline{I}(c,b;\,u.\,h)/u\rangle}.
\]
The calculus $\mathsf{C}_{\mathrm{Seed}}$ is
$\mathsf{C}_{\mathrm{Step}}$ (\cref{comp:def:cstep}) with \rulename{SeedEval}
in place of \rulename{Step}.
\end{definition}

The counter, the initializer and the body are read off the second premise,
and the left side of the first premise is compared with the initializer $a$
alone. In comparison with \rulename{Step}, the first premise no longer
mentions the preceding iterate, and the conclusion changes at the same time:
it iterates from the equal seed $b$ before the last application of the
body. At $c\equiv0$ the two rules are the same inference,
\[
 \infrule{a=b\qquad\It{1}{a}{u}{h}=r}{r=h\langle b/u\rangle},
\]
and by \cref{comp:lem:stepderived} the pivot, symmetry, transitivity,
congruence and \rulename{Tr} are derived from reflexivity and this shared
instance. They are therefore available with either rule.

\begin{lemma}[The two rules are interderivable]
\label{comp:lem:seedstep}
Over an arbitrary context and with reflexivity:
\begin{enumerate}[label=(\alph*)]
\item \rulename{SeedEval} derives successor computation, by one instance on
 two reflexivities;
\item every instance of \rulename{SeedEval} is derivable from
 \rulename{Step}, by two instances at the counter $0$ and one further
 instance;
\item every instance of \rulename{Step} is derivable from
 \rulename{SeedEval}, by seven instances.
\end{enumerate}
At $c\equiv0$ no expansion is needed in (b) and (c). None of the derivations
uses a binder rule, a case distinction or induction.
\end{lemma}

\begin{proof}
Put $P\defeq\overline{I}(c,a;u.h)$ and $B\defeq\It{\Suc c}{a}{u}{h}$.

(a) \rulename{SeedEval} applied to $a=a$ and $B=B$ concludes
$B=h\langle P/u\rangle$. Its first premise does not contain $P$.

(b) Let $a=b$ and $B=r$ be given. Congruence
(\cref{comp:lem:stepderived}(d)) in the template
$v\mapsto\overline{I}(c,v;u.h)$, with $v$ fresh for $c$ and $h$, gives
$\overline{I}(c,a;u.h)=\overline{I}(c,b;u.h)$ by two instances of
\rulename{Step} at the counter $0$. \rulename{Step} with this equation and
$B=r$ gives $r=h\langle\overline{I}(c,b;u.h)/u\rangle$. The body $h$ is
fixed throughout; no agreement of two different bodies is proved.

(c) Let $P=t$ and $B=r$ be given. By (a), $B=h\langle P/u\rangle$. The pivot
applied to $B=r$ and $B=h\langle P/u\rangle$ gives $h\langle P/u\rangle=r$.
\rulename{Tr} along $P=t$ in the template $h$ gives
$r=h\langle t/u\rangle$. The derivation has one instance for the
computation, two for the pivot and four for \rulename{Tr}.
\end{proof}

\begin{theorem}
\label{comp:thm:seedeval}
For every finite context $\Gamma$ of clauses and every formula $A$ of the
clause calculi,
\[
 \Gamma\proves_{\mathsf{C}_{\mathrm{Step}}}A
 \quad\Longleftrightarrow\quad
 \Gamma\proves_{\mathsf{C}_{\mathrm{Seed}}}A .
\]
The sequent is literally unchanged. The same equivalence holds for the
calculi with native contexts: the calculus with the rules \rulename{Hyp},
\rulename{Refl}, \rulename{SeedEval}, \rulename{QE}, \rulename{DE},
\rulename{DI} and \rulename{QI}, and the calculus with \rulename{WI} in
place of the last two, derive the sequents of $\Nat^{0}$. Consequently
\cref{comp:thm:clausal} holds for
$\mathsf{C}_{\mathrm{Seed}}$.
\end{theorem}

\begin{proof}
Replace each arithmetic inference of a derivation by its expansion in
\cref{comp:lem:seedstep} and keep every other inference. The two calculi
have the same terms, formulas, clauses, contexts, substitution and freshness
conditions, and an expansion has the same premises and conclusion as the
inference that it replaces. No witness is changed, so the mode of each
\rulename{WI} is unchanged. The statement for native contexts follows in
the same way from \cref{comp:thm:step}.
\end{proof}

The seven-rule calculus with native contexts is formalized as
\Lean{NVPrf}, with \Lean{nvprf\_iff\_nprf}, \Lean{NVPrf.normal} and
\Lean{nvprf\_iff\_pra}. The fusion differs from the rule of
\pkg{initializer\_equality}, which keeps successor computation as a separate
rule, and from the package \pkg{counter\_evaluation}, which fuses evaluation
with an equation between counters for a recursor whose body binds the
counter and the state.

\begin{proposition}[Size of the computation certificate]
\label{comp:prop:seedcert}
Let $c\not\equiv0$. The derivation of
$B=h\langle P/u\rangle$ by \rulename{Step} is
$\mathrm{Step}(\mathrm{Refl}(P),\mathrm{Refl}(B))$, and by
\rulename{SeedEval} it is
$\mathrm{SeedEval}(\mathrm{Refl}(a),\mathrm{Refl}(B))$. The two have the
same number of inference nodes, and the second contains
$|P|-|a|=1+|c|+|h|$ fewer nodes of term syntax, where $|\cdot|$ counts one
node for each variable, constant and constructor. In a serialization that
writes $\It{c}{a}{u}{h}$ as the JSON list \texttt{["i",c,a,h]} without
optional whitespace, the second derivation is shorter by
$8+\ell(c)+\ell(h)$ bytes, where $\ell$ is the length of the serialized
term. At $c\equiv0$ the two derivations coincide.
\end{proposition}

\begin{proof}
$|P|=1+|c|+|a|+|h|$, and the only difference between the two derivations is
the term of the first reflexivity. The brackets, the tag and the separators
of the list contribute eight bytes, and the counter and the body are no
longer repeated. The conclusion is the same in both derivations.
\end{proof}

The saving grows without bound with the body. It is established for these
derivations of a single successor computation; the translation of
\cref{comp:lem:seedstep}(c) expands other derivations, and a format that
shares subterms changes the comparison.

\subsection{Suspended application}
\label{comp:sec:app}

The second construction of \pkg{seed\_evaluation} removes the counter from
the arithmetic rule altogether. The price is a normalization between the
zero normalization and the full normalization of \cref{comp:sec:full}: an
iterator at a successor counter is factored into a one-step application
whose body is never executed by the normalizer.

\begin{definition}[Unit normalization]
\label{comp:def:unit}
Extend the raw terms by the constructor $U(a;\,u.\,h)$, which binds $u$ in
$h$ and not in $a$. Its \emph{expansion} is
$U(a;u.h)^{\circ}\defeq\It{1}{a^{\circ}}{u}{h^{\circ}}$, and $(\cdot)^{\circ}$
commutes with the other constructors. On terms $n,a,h$ define
\[
 K(n,a;\,u.\,h)\defeq
 \begin{cases}
  a, & n\equiv0,\\
  U(K(m,a;\,u.\,h);\,u.\,h), & n\equiv\Suc m,\\
  \It{n}{a}{u}{h}, & \text{otherwise},
 \end{cases}
\]
by recursion on $n$, and define $\nf_U$ by structural recursion:
$\nf_U(x)\defeq x$, $\nf_U(0)\defeq0$, $\nf_U(\Suc t)\defeq\Suc\nf_U(t)$,
\[
 \nf_U(U(a;u.h))\defeq U(\nf_U(a);u.\nf_U(h)),\qquad
 \nf_U(\It{n}{a}{u}{h})\defeq K(\nf_U(n),\nf_U(a);\,u.\,\nf_U(h)).
\]
A term is \emph{unit-normal} if it is fixed by $\nf_U$: its subterms are
unit-normal and the counter of each of its iterators has neither $0$ nor
$\Suc$ at its root.
\end{definition}

For example, for a variable $x$,
\[
 \nf_U\bigl(\It{\Suc\Suc x}{a}{u}{h}\bigr)\equiv
 U\bigl(U(\It{x}{\nf_U(a)}{u}{\nf_U(h)};u.\nf_U(h));u.\nf_U(h)\bigr).
\]
Each displayed $u$ binds its own copy of the body. The result is obtained by
inspecting the leading successors of the counter; neither $h\subst{a}{u}$
nor any other execution of $h$ occurs. A suspended application is always
kept, also when its body is constant or the identity. In $\Teq$ the
expansion of $U(a;u.h)$ equals $h^{\circ}\subst{a^{\circ}}{u}$, so the
constructor adds no numerical freedom and no strength.

\begin{lemma}[Termination and uniqueness]
\label{comp:lem:unitrewrite}
Consider the two contractions
\[
 \It{0}{a}{u}{h}\to a,\qquad
 \It{\Suc n}{a}{u}{h}\to U(\It{n}{a}{u}{h};\,u.\,h),
\]
closed under every context, with no contraction for $U$. Every reduction
sequence is finite, every term reduces to $\nf_U(t)$, and $\nf_U(t)$ is the
unique normal form of $t$. In particular the reduction is confluent, and
$\nf_U$ is idempotent and primitive recursive on codes of terms.
\end{lemma}

\begin{proof}
Assign positive integer weights by $w(x)=w(0)=1$, $w(\Suc t)=w(t)+1$,
\[
 w(U(a;u.h))=w(a)+w(h)+1,\qquad
 w(\It{n}{a}{u}{h})=(w(n)+1)(w(a)+w(h)+1).
\]
Write $A=w(a)$, $H=w(h)$ and $q=w(n)$. The zero contraction decreases the
weight from $2(A+H+1)$ to $A$. The second contraction changes it from
$(q+2)(A+H+1)$ to $(q+1)(A+H+1)+H+1$, a decrease by exactly $A>0$; the
duplication of the body is accounted for. Every constructor is strictly
increasing in the weight of each child, so the decrease persists in every
context, including bodies. Hence every reduction decreases a positive
integer.

$\nf_U$ is invariant under both contractions at the root: for the first by
its definition, and for the second because
\[
 \nf_U(\It{\Suc n}{a}{u}{h})
 \equiv U\bigl(K(\nf_U(n),\nf_U(a);u.\nf_U(h));u.\nf_U(h)\bigr)
 \equiv\nf_U\bigl(U(\It{n}{a}{u}{h};u.h)\bigr).
\]
A structural induction extends the invariance to contractions in context.
The irreducible terms are exactly the unit-normal terms. By termination
every term has an irreducible reduct, which by invariance is $\nf_U(t)$. Any
two reducts of $t$ reduce to this term.

With a primitive recursive coding of terms, the weight and a deterministic
choice of one contraction are primitive recursive, and fewer than $w(t)$
contractions occur; the bounded iteration of the step computes $\nf_U$. The
normalizer is not an evaluator of primitive recursive programs: it never
substitutes an argument into the body of $U$.
\end{proof}

\begin{lemma}[Coherence with substitution]
\label{comp:lem:unitsubst}
$\nf_U(t[\sigma])\equiv\nf_U\bigl(\nf_U(t)[\nf_U\circ\sigma]\bigr)$ for all
terms $t$ and substitutions $\sigma$. Hence the normalizing substitution
$t\langle\sigma\rangle_U\defeq\nf_U(t[\sigma])$ satisfies the laws of
\cref{comp:lem:laws}; $\nf_U$ commutes with renaming and introduces no free
variable. On terms without $U$, $\nf_U\circ\nf_0=\nf_U$.
\end{lemma}

\begin{proof}
Both contractions are preserved by every capture-free substitution, and
contracting inside the terms substituted into a fixed term is a reduction of
the instance. Hence $t[\sigma]$ reduces to $\nf_U(t)[\nf_U\circ\sigma]$, and
both have the normal form of the left side, by
\cref{comp:lem:unitrewrite}. Below $U$ and the body of an iterator the
substitution is lifted. The laws follow as in \cref{comp:lem:laws}, whose
proof uses only ($\nu$1)--($\nu$4). The last statement holds because
$t\to^{*}\nf_0(t)$ by zero contractions, which are among the two
contractions.
\end{proof}

\begin{definition}[The rule \rulename{App} and the calculus
$\mathsf{C}_{\mathrm{App}}$]
\label{comp:def:app}
On unit-normal terms the rule \rulename{App} is
\[
 \infrule{\Gamma\proves a=b\qquad\Gamma\proves U(a;\,u.\,h)=r}
         {\Gamma\proves r=h\langle b/u\rangle_U}.
\]
The calculus $\mathsf{C}_{\mathrm{App}}$ has unit-normal terms, the formulas
and clauses of \cref{comp:sec:clausal} over them, the substitution
$\langle\cdot\rangle_U$, and the five rules \rulename{Use}, \rulename{Refl},
\rulename{App}, \rulename{DE} and \rulename{WI}, the last three logical
rules as in \cref{comp:def:cstep}.
\end{definition}

\rulename{App} inspects no counter and builds no predecessor iterate; that
work is done by the constructor $K$.

\begin{lemma}
\label{comp:lem:appderived}
From \rulename{Refl} and \rulename{App}, over an arbitrary context:
\begin{enumerate}[label=(\alph*)]
\item the pivot, symmetry, transitivity, congruence in every unit-normal
 template and \rulename{Tr} are derivable, without induction or case
 distinction;
\item for unit-normal $c,a,h$,
 $K(\Suc c,a;u.h)=h\langle K(c,a;u.h)/u\rangle_U$ is derivable by one
 \rulename{App} on two reflexivities;
\item the image under $\nf_U$ of an instance of \rulename{Step} on
 zero-normal terms is an instance of \rulename{App}.
\end{enumerate}
\end{lemma}

\begin{proof}
(a) As in \cref{comp:lem:stepderived}, with the frame
$J_s\defeq U(s;v.v)$ for the pivot and $H_s\defeq U(s;v.h)$ for congruence:
\rulename{App} on $s=r$ and $J_s=J_s$ gives $J_s=r$, and \rulename{App} on
$s=t$ and $J_s=r$ gives $r=t$; \rulename{App} on $s=s$ and $H_s=H_s$ gives
$H_s=h\langle s/v\rangle_U$, and \rulename{App} on $s=t$ and this equation
gives $h\langle s/v\rangle_U=h\langle t/v\rangle_U$.

(b) By definition $K(\Suc c,a;u.h)\equiv U(K(c,a;u.h);u.h)$. Apply
\rulename{App} to the reflexivities of $K(c,a;u.h)$ and of this term.

(c) Let the instance have the premises $\overline{I}(c,a;u.h)=t$ and
$\It{\Suc c}{a}{u}{h}=r$ and the conclusion $r=\nf_0(h[t/u])$. Write
$c',a',h'$ for the images of $c,a,h$ under $\nf_U$. Then
$\nf_U(\overline{I}(c,a;u.h))\equiv K(c',a';u.h')$, in both cases of the
smart constructor, and
$\nf_U(\It{\Suc c}{a}{u}{h})\equiv U(K(c',a';u.h');u.h')$. \rulename{App}
concludes $\nf_U(r)=h'\langle\nf_U(t)/u\rangle_U$, and by
\cref{comp:lem:unitsubst} this term is
$\nf_U(h[t/u])\equiv\nf_U(\nf_0(h[t/u]))$.
\end{proof}

\begin{lemma}
\label{comp:lem:unitzero}
For every term $t$ without $U$, $\nf_U(t)\equiv0$ if and only if
$\nf_0(t)\equiv0$. In particular, for a zero-normal $t$,
$\nf_U(t)\equiv0$ if and only if $t\equiv0$.
\end{lemma}

\begin{proof}
By structural induction. A variable and a successor do not normalize to
$0$ under either map. For an iterator, $\nf_U$ yields $0$ exactly when the
normalized counter and the normalized initializer are both $0$: otherwise
the result is the initializer, which is not $0$, or has $U$ or an iterator
at its root. The same holds for $\nf_0$, and the induction hypotheses for
the counter and the initializer apply.
\end{proof}

\begin{theorem}
\label{comp:thm:app}
Let $\Gamma$ be a context of clauses and $A$ a formula of
$\mathsf{C}_{\mathrm{Step}}$.
\begin{enumerate}[label=(\alph*)]
\item If $\Gamma\proves_{\mathsf{C}_{\mathrm{Step}}}A$, then
 $\nf_U(\Gamma)\proves_{\mathsf{C}_{\mathrm{App}}}\nf_U(A)$.
\item $\mathsf{C}_{\mathrm{App}}$ is sound for $\Teq$ under the expansion
 $(\cdot)^{\circ}$ and the readings of clauses and formulas, and $\Teq$
 proves $t^{\circ}=\nf_U(t)^{\circ}$ for every term $t$.
\item If $\Gamma$ consists of assumption clauses and
 $\nf_U(\Gamma)\proves_{\mathsf{C}_{\mathrm{App}}}\nf_U(A)$, then
 $\Gamma\proves_{\mathsf{C}_{\mathrm{Step}}}A$.
\item For a unit-normal $t$, $\nf_U(t^{\circ})\equiv t$.
\end{enumerate}
Consequently the ordinary realization of $\mathsf{C}_{\mathrm{App}}$ is
bi-interpretable and equiconsistent with $\iPRA$, with whole domains,
equality interpreted by equality and numerical composites with the graph
$x=y$.
\end{theorem}

\begin{proof}
(a) By induction on the derivation, translating terms by $\nf_U$, formulas
termwise and contexts clause by clause. \rulename{Refl} is immediate, and
each \rulename{Step} becomes one \rulename{App} by
\cref{comp:lem:appderived}(c). \rulename{Use} and \rulename{DE} are
translated by \cref{comp:lem:unitsubst}. For \rulename{WI} the root of the
matrix is preserved, and by \cref{comp:lem:unitzero} the witness is
literally $0$ after the translation exactly when it was before, so the mode
is the same; freshness is preserved because $\nf_U$ introduces no variable.

(b) In $\Teq$ the zero contraction is the zero computation equation. For
the second contraction both sides equal
$h\subst{\It{n}{a}{u}{h}}{u}$: the left by successor computation, the
expansion of the right by the successor and the zero computation at the
counter $1$. Contractions in the counter or the initializer are lifted by
congruence, and contractions in a body by body congruence
(\cref{base:cor:rot}(ii), transferred to $\Teq$), which uses the induction for
literal equations. Hence $\Teq\proves t^{\circ}=\nf_U(t)^{\circ}$.
\rulename{App} is sound by the computation at the counter $1$ and equality,
and the logical rules are sound under the ordinary reading as recorded
after \cref{claus:cor:bridge}: \rulename{Use} and \rulename{DE} are
ordinary eliminations, \rulename{WI} in case mode uses coverage and in
progress mode the induction for a literal equation under the assumed base,
and a normalized instance and the ordinary instance are provably equal.

(c) By (b), $\Teq$ proves the reading of $A$ from the readings of $\Gamma$.
If $\Gamma$ is $\operatorname{ord}(\Delta)$, then $\Delta$ and $A$ are
zero-normal native formulas without $\Q{y}$, so
$\Delta\proves_{\Nat}A$ by \cref{base:thm:bridge}, and
\cref{claus:thm:zero}(iii) gives
$\Gamma\proves_{\mathsf{C}_{\mathrm{Step}}}A$.
% ed.: the notes of seed_evaluation assert (c) for every clause context by
% reference to the baseline; the argument given there, and here, yields it
% for contexts of assumption clauses.

(d) By structural induction: the expansion $\It{1}{a^{\circ}}{u}{h^{\circ}}$
of an application normalizes to $U(a;u.h)$, and an iterator of a
unit-normal term has a counter with neither $0$ nor $\Suc$ at its root,
whose normalized expansion is the same counter.

The last statement follows from (a)--(d) and
\cref{comp:thm:clausal}: the constructor $U$ is an abbreviation of a
fixed term, and normalization changes representatives of expressions only.
\end{proof}

\begin{proposition}[The variant without a new constructor]
\label{comp:prop:unitabbrev}
Let $U(a;u.h)$ abbreviate $\It{1}{a}{u}{h}$ in the unchanged term grammar,
and define $K$ and $\nf_U$ by the clauses of \cref{comp:def:unit}, without
the clause for $U$. The fixed points of $\nf_U$ are the terms in which
every iterator counter is the literal $1$ or has a root other than $0$ and
$\Suc$. The normal forms are those of the contractions
\[
 \It{0}{a}{u}{h}\to a,\qquad
 \It{\Suc n}{a}{u}{h}\to\It{1}{\It{n}{a}{u}{h}}{u}{h}\quad(n\not\equiv0),
\]
which terminate; \cref{comp:lem:unitsubst,comp:lem:unitzero} hold; the rule
\rulename{App} is the instance of \rulename{Step} at the counter $0$,
followed by normalization; and \cref{comp:thm:app} holds for the resulting
calculus.
\end{proposition}

\begin{proof}
The side condition on the second contraction is necessary: without it
$\It{1}{a}{u}{h}$ would be contracted to
$\It{1}{\It{0}{a}{u}{h}}{u}{h}$, which the zero contraction contracts back.
For termination use the weights $w(x)=w(0)=1$, $w(\Suc t)=w(t)+1$ and, with
$A=w(a)$ and $H=w(h)$,
\[
 w(\It{n}{a}{u}{h})=
 \begin{cases}
  A+H+1, & n\equiv1,\\
  (w(n)+1)(A+H+1), & \text{otherwise}.
 \end{cases}
\]
The zero contraction decreases the weight from $2(A+H+1)$ to $A$. In the
second contraction the left counter $\Suc n$ is not $1$, so the left weight
is $(w(n)+2)(A+H+1)$. If $n\not\equiv1$ the right weight is
$(w(n)+1)(A+H+1)+H+1$, smaller by $A$. If $n\equiv1$ the right weight is
$A+2H+2$ and the left weight is $4(A+H+1)$. The weight is strictly
increasing in the initializer and the body. In the counter position the one
exceptional case is a contraction that makes the counter the literal $1$,
which replaces the product by the smaller sum; a counter that is the
irreducible term $1$ is not contracted. Hence every contraction in context
decreases the weight.

Both contractions preserve $\nf_U$; for the second this is the identity
$K(\Suc n,a;u.h)\equiv U(K(n,a;u.h);u.h)$, which holds also at $n\equiv0$.
For coherence with substitution the conditional contraction is not closed
under substitution, since a substitution can turn $n$ into a term with the
normal form $0$. Argue in two steps. First, replacing each substituted term
by its normal form does not change the normal form of the instance, by the
invariance in context. Second, if $t\to t'$ by one contraction, then
$\nf_U(t[\sigma])\equiv\nf_U(t'[\sigma])$: for the zero contraction
literally, and for the second by the displayed identity for $K$, which is
valid at $0$. Following a reduction of $t$ to $\nf_U(t)$ gives the
statement. The remaining proofs are those of
\cref{comp:lem:appderived,comp:lem:unitzero,comp:thm:app}; the unit
iterator is never removed and its body is never exposed.
\end{proof}

%%%%%%%%%%%%%%%%%%%%%%%%%%%%%%%%%%%%%%%%%%%%%%%%%%%%%%%%%%%%%%%%%%%%%%%%%%%%
\section{Atomic-seed zero contraction with raw rule metadata}
\label{comp:sec:atomic}

This section draws on the mathematical result of
\pkg{atomic\_seed\_refinement}: a six-rule presentation of $\Nat$ whose
normalizer contracts $\It{0}{a}{u}{h}$ only when the normalized initializer
$a$ is a variable or the literal $0$, and whose rules keep their raw
instances. That normalizer is weaker than the zero normalization of
\cref{comp:sec:zero} and is not an admissible normalizer: it violates the
coherence ($\nu$4), and \cref{comp:thm:transfer} does not apply to it. The
first subsection gives the form of the transfer argument that needs no
coherence, at the price of raw metadata in every rule; it applies to every
normalizer that satisfies the other four conditions of
\cref{comp:def:normalizer}. The remaining subsections
apply it to the atomic-seed normalizer.

\subsection{Normalization after raw instantiation}

\begin{definition}[Conversion normalizer]
\label{comp:def:convnorm}
A \emph{conversion normalizer} is a map $\nu$ from raw terms to raw terms
that satisfies ($\nu$1), ($\nu$2), ($\nu$3) and ($\nu$5) of
\cref{comp:def:normalizer}, extended to formulas and contexts as there.
Every admissible normalizer is a conversion normalizer.
\end{definition}

The proof of \cref{comp:lem:nusound} uses ($\nu$5) and the extension of
$\nu$ to formulas only. The lemma therefore holds for every conversion
normalizer: $\Gamma,A\proves_{\Nat}\nu(A)$ and
$\Gamma,\nu(A)\proves_{\Nat}A$ for every raw formula $A$ and every context
$\Gamma$.

\begin{definition}[Raw-metadata reading]
\label{comp:def:rawmeta}
Let $\nu$ be a conversion normalizer. A \emph{sequent over $\nu$} is a
sequent $\Sigma\proves C$ in which $\Sigma$ is a finite list of $\nu$-normal
formulas and $C$ is $\nu$-normal. For each rule $R$ among the eight rules of
$\Nat$ and \rulename{WI} (\cref{comp:def:wi}, on raw terms), the rule
$R^{[\nu]}$ on sequents over $\nu$ is the following. $\rulename{Hyp}^{[\nu]}$
derives $\Sigma\proves C$ for $C\in\Sigma$. For every other $R$, take an
instance of $R$ on raw syntax in the context $\Sigma$. It is fixed by its
\emph{raw metadata}: the raw terms, templates, matrices and the witness that
the schema of $R$ mentions, and for \rulename{DE} the conclusion $A$, which is
required to be $\nu$-normal. Let its premises be $\Sigma,L_i\proves P_i$,
where $L_i$ is the local assumption of the premise if there is one, and let
its conclusion be $\Sigma\proves C_0$. Then $R^{[\nu]}$ is the inference
\[
 \infrule{\Sigma,\nu(L_i)\proves\nu(P_i)\quad(\text{all }i)}
         {\Sigma\proves\nu(C_0)}.
\]
The side conditions are tested on the raw instance. The eigenvariable ($b$
of \rulename{DE}, $z$ of \rulename{DI}, \rulename{QI} and \rulename{WI}) is
not free in $\Sigma$, in the conclusion, or in any of the raw metadata; the
test of \rulename{QI} and \rulename{WI} for a literal equation is applied to
the raw matrix, which has the root of its normal form; and the test
$t\equiv0$ of \rulename{WI} is applied to the raw witness. A \emph{raw-metadata
calculus over $\nu$} is a calculus on sequents over $\nu$ whose rules are
rules $R^{[\nu]}$. The calculus $\Nat^{[\nu]}$ has the eight rules
$R^{[\nu]}$ for the rules $R$ of $\Nat$.
\end{definition}

\begin{lemma}[Soundness of raw-metadata rules]
\label{comp:lem:rawsound}
Let $\nu$ be a conversion normalizer and let $R$ be one of the eight rules of
$\Nat$ or \rulename{WI}. If the premises of an instance of $R^{[\nu]}$ are
derivable in $\Nat$, read as raw sequents, then so is its conclusion. Hence
$\Sigma\proves_{\mathsf K}C$ implies $\Sigma\proves_{\Nat}C$ for every
raw-metadata calculus $\mathsf K$ over $\nu$.
\end{lemma}

\begin{proof}
For \rulename{Hyp} there is nothing to prove. Otherwise let
$\Sigma,\nu(L_i)\proves_{\Nat}\nu(P_i)$ for every premise. By
\cref{comp:lem:nusound}, $\Sigma,L_i\proves_{\Nat}\nu(L_i)$ and
$\Sigma,L_i,\nu(P_i)\proves_{\Nat}P_i$, and two cuts (\cref{base:lem:cut})
give $\Sigma,L_i\proves_{\Nat}P_i$. These are the premises of the raw
instance of $R$. Its side conditions hold in $\Nat$: the conditions of
$\Nat$ and of \cref{kern:def:W} require that the eigenvariable is not free
in the context, in the conclusion, in the witness or in the formula
$\D{x}{y}B$ or the parameters of the matrix, and all of these are among the
data tested. For a rule of $\Nat$ the rule gives $\Sigma\proves_{\Nat}C_0$.
For \rulename{WI} the mode is that of the raw instance, and \rulename{WI}
on raw terms is admissible in $\Nat$ by \cref{comp:lem:single}(iii). For
\rulename{DE} the conclusion $C_0$ is $\nu$-normal; in the other cases
\cref{comp:lem:nusound} and cut give $\Sigma\proves_{\Nat}\nu(C_0)$. The
second statement follows by induction on derivations.
\end{proof}

\begin{theorem}[Transfer with raw metadata]
\label{comp:thm:rawtransfer}
Let $\nu$ be a conversion normalizer.
\begin{enumerate}[label=(\alph*)]
\item For every finite context $\Gamma$ of raw native formulas and every raw
 native formula $A$,
 \[
  \Gamma\proves_{\Nat}A
  \quad\Longleftrightarrow\quad
  \nu(\Gamma)\proves_{\Nat^{[\nu]}}\nu(A).
 \]
\item Let $\mathsf K$ be a raw-metadata calculus over $\nu$ in which every
 rule of $\Nat^{[\nu]}$ is admissible. Then $\mathsf K$ derives exactly the
 sequents over $\nu$ that $\Nat$ derives, and (a) holds with $\mathsf K$ in
 place of $\Nat^{[\nu]}$.
\item If $\nu$ is an admissible normalizer, the rules of $\Nat^{[\nu]}$ and
 the rules of $\Nat^{\nu}$ (\cref{comp:def:nukernel}) have the same
 instances.
\end{enumerate}
Both directions of (a) and (b) are effective on finite derivations.
\end{theorem}

\begin{proof}
(a) From left to right, by induction on the derivation in $\Nat$, replace
each sequent $\Gamma'\proves A'$ by $\nu(\Gamma')\proves\nu(A')$ and each
inference of a rule $R$ by the inference of $R^{[\nu]}$ with the same raw
metadata. An assumption $A'\in\Gamma'$ gives $\nu(A')\in\nu(\Gamma')$. The
premises, the local assumptions and the conclusion of the inference of
$R^{[\nu]}$ are by definition the normal forms of those of the raw
inference; for \rulename{DE} the conclusion $\nu(A')$ is $\nu$-normal by
($\nu$3). An eigenvariable that is not free in $\Gamma'$, in $A'$ or in the
raw metadata is not free in $\nu(\Gamma')$ or in $\nu(A')$ by ($\nu$2). No
property of $\nu$ under substitution is used. From right to left,
\cref{comp:lem:rawsound} gives $\nu(\Gamma)\proves_{\Nat}\nu(A)$, and
\cref{comp:lem:nusound} with cut gives $\Gamma\proves_{\Nat}A$.

(b) A sequent over $\nu$ that $\mathsf K$ derives is derivable in $\Nat$ by
\cref{comp:lem:rawsound}. Conversely let $\Delta\proves_{\Nat}C$ with
$\Delta$ and $C$ $\nu$-normal. By (a) and ($\nu$3),
$\Delta\proves_{\Nat^{[\nu]}}C$, and induction on this derivation, with the
admissibility of each rule in $\mathsf K$, gives $\Delta\proves_{\mathsf K}C$.
For (a) with $\mathsf K$: $\Gamma\proves_{\Nat}A$ holds if and only if
$\nu(\Gamma)\proves_{\Nat}\nu(A)$, by \cref{comp:lem:nusound} and cut, and
the latter holds if and only if $\nu(\Gamma)\proves_{\mathsf K}\nu(A)$.

(c) Let $m$ be the raw metadata of an instance of a rule of $\Nat^{[\nu]}$.
By ($\nu$4) and \cref{comp:lem:laws}(a), the normal form of each premise,
local assumption and conclusion is the corresponding formula of the instance
at the metadata $\nu(m)$, formed with normalizing substitution; for example
$\nu(h[s/v])\equiv\nu(h)\langle\nu(s)/v\rangle$, and
$\nu(\Q{y}B[b/x])\equiv\Q{y}\nu(B)[b/x]$ by ($\nu$2). By ($\nu$2) an
eigenvariable that is fresh for $m$ is fresh for $\nu(m)$, and the root of
the matrix of \rulename{QI} is unchanged. Hence the instance is an instance of
the same rule of $\Nat^{\nu}$; the instances of $\rulename{I0}^{[\nu]}$ and
$\rulename{IS}^{[\nu]}$ are exactly those of \rulename{Comp}$^{\nu}$.
Conversely an instance of a rule of $\Nat^{\nu}$ is the instance of the
corresponding rule of $\Nat^{[\nu]}$ whose metadata are $\nu$-normal.
\end{proof}

For an admissible normalizer, (a) and (c) give a second proof of
\cref{comp:thm:transfer}; in it, ($\nu$4) is used only to show that normal
metadata suffice. For a conversion normalizer without ($\nu$4), normal
metadata need not suffice, and the instances with raw metadata are the
reason why the rule-by-rule translation in (a) works.

\begin{remark}
\label{comp:rem:rawwi}
The rule $\rulename{WI}^{[\nu]}$ selects its mode from the raw witness and
is sound for every conversion normalizer by \cref{comp:lem:rawsound}, also
when the raw witness has the normal form $0$. The mode change of
\cref{comp:ex:modechange} arises only when the witness is normalized before
the mode is selected. Read as an inference of $\rulename{WI}^{[\nf]}$, the
introduction of that example at the raw witness $\It{0}{0}{u}{0}$ stays in
case mode, and its normalized premise is $0=1\proves\Suc z=1$.
\end{remark}

\begin{corollary}[Interpretation of raw-metadata calculi]
\label{comp:cor:rawbridge}
Let $\nu$ be a conversion normalizer and let $\mathsf K$ be a calculus as in
\cref{comp:thm:rawtransfer}(b). Then statements (a) and (c) of
\cref{comp:cor:bridge} hold for $\mathsf K$. Weakening and cut are admissible
in $\mathsf K$, and so are the substitution of terms, from
$\Delta\proves C$ infer $\nu(\Delta[\sigma])\proves\nu(C[\sigma])$, and the
Leibniz rule with raw metadata, from $\Sigma\proves\nu(s=t)$ and
$\Sigma\proves\nu(F[s/v])$ infer $\Sigma\proves\nu(F[t/v])$, for raw
$s,t,F$.
\end{corollary}

\begin{proof}
Statement (a) of \cref{comp:cor:bridge} is \cref{comp:thm:rawtransfer}(b).
The proof of statement (c) given there uses of $\nu$ only statement (a),
\cref{comp:lem:nusound}, ($\nu$1) and ($\nu$2), and applies unchanged. For
the admissible rules pass through $\Nat$. A derivation in $\mathsf K$ gives
one in $\Nat$ by \cref{comp:lem:rawsound}. Weakening and substitution in
$\Nat$ are \cref{base:lem:struct}, cut is \cref{base:lem:cut}, and the
Leibniz rule is \cref{base:thm:leibniz}, with \cref{comp:lem:nusound} to
convert its premises from and its conclusion to normal form. The resulting
sequent of $\Nat$ is returned to $\mathsf K$ by
\cref{comp:thm:rawtransfer}(b), applied to its normal form.
\end{proof}

Unlike \cref{comp:lem:nustruct}, the corollary obtains substitution in
derivations through $\Nat$ and not by a direct induction on derivations of
$\mathsf K$; no coherence of $\nu$ with substitution is needed.

\subsection{The atomic-seed normalizer}

\begin{definition}[Atomic-seed zero contraction]
\label{comp:def:nfatomic}
A raw term is \emph{atomic} if it is a variable or the literal $0$. The map
$\nf_{\mathrm a}$ on raw terms is defined by structural recursion:
$\nf_{\mathrm a}(x)\defeq x$, $\nf_{\mathrm a}(0)\defeq0$,
$\nf_{\mathrm a}(\Suc t)\defeq\Suc\nf_{\mathrm a}(t)$ and, with
$n'\defeq\nf_{\mathrm a}(n)$, $a'\defeq\nf_{\mathrm a}(a)$ and
$h'\defeq\nf_{\mathrm a}(h)$,
\[
 \nf_{\mathrm a}(\It{n}{a}{u}{h})\defeq
 \begin{cases}
  a', & n'\equiv0\text{ and }a'\text{ is atomic},\\
  \It{n'}{a'}{u}{h'}, & \text{otherwise}.
 \end{cases}
\]
It is extended to formulas and contexts as in \cref{comp:def:normalizer}. A
raw term is \emph{atomic-seed normal} if no iterator in it, at any depth, has
the literal $0$ as its counter and an atomic initializer.
\end{definition}

The recursion is that of $\nf_0$ (\cref{comp:def:nfzero}) with one more
test, on the normalized initializer. For example
$\nf_{\mathrm a}(\It{0}{x}{u}{h})\equiv x$ and
$\nf_{\mathrm a}(\It{\It{0}{0}{v}{0}}{y}{u}{h})\equiv y$, while
$\It{0}{\Suc0}{u}{0}$ is atomic-seed normal and is not zero-normal.

\begin{lemma}[Properties of the atomic-seed normalizer]
\label{comp:lem:nfatomic}
\begin{enumerate}[label=(\alph*)]
\item $\nf_{\mathrm a}(t)$ is atomic-seed normal; $\nf_{\mathrm a}(t)\equiv t$
 if and only if $t$ is atomic-seed normal; $\nf_{\mathrm a}$ is idempotent.
\item $\nf_{\mathrm a}(t[\rho])\equiv\nf_{\mathrm a}(t)[\rho]$ for every
 renaming $\rho$ of variables, and
 $\FV(\nf_{\mathrm a}(t))\subseteq\FV(t)$. $\nf_{\mathrm a}(t)$ has at most
 as many nodes as $t$, and $\nf_{\mathrm a}$ is a primitive recursive
 operation on codes of finite terms.
\item $t\to^{*}\nf_{\mathrm a}(t)$ by zero contractions. Hence
 $\nf_0(\nf_{\mathrm a}(t))\equiv\nf_0(t)$; every zero-normal term is
 atomic-seed normal, and $\nf_{\mathrm a}(\nf_0(t))\equiv\nf_0(t)$.
\item \emph{Fixed-context congruence.} Let $C$ be a raw term or formula, $x$
 a variable, and $s,s'$ raw terms none of whose free variables is bound in
 $C$ at a free occurrence of $x$. If
 $\nf_{\mathrm a}(s)\equiv\nf_{\mathrm a}(s')$, then
 $\nf_{\mathrm a}(C[s/x])\equiv\nf_{\mathrm a}(C[s'/x])$.
\end{enumerate}
\end{lemma}

\begin{proof}
(a) By structural induction. The children of a rebuilt iterator are
atomic-seed normal and the iterator is rebuilt only when the test fails; a
returned initializer is atomic-seed normal. On an atomic-seed normal term
every clause returns its argument: the children are fixed by the induction
hypothesis, and the test fails at a normal iterator.
(b) A renaming sends variables to variables and fixes $0$, so it sends
atomic terms to atomic terms, other terms to other terms, and a counter to
the literal $0$ exactly when the counter is $0$; by induction it commutes
with the test and with the recursion.
% ed.: the package states commutation with injective renamings; the
% argument covers every renaming of variables, as ($\nu$2) requires.
A contraction deletes the counter and the body and keeps the normalized
initializer, so it creates no variable and adds no node. The structural
recursion with these size bounds is primitive recursive, as in
\cref{comp:lem:nfzero}(c); no successor contraction is performed.
(c) By structural induction: the children reduce to their normal forms by
zero contractions in their positions, since the reduction of
\cref{comp:def:reduction} is closed under all term contexts including the
bodies of iterators, and the step at the root, if it is taken, is a zero
contraction. By \cref{comp:lem:nfzero}(d) the zero normal form is unique,
so $\nf_0(\nf_{\mathrm a}(t))\equiv\nf_0(t)$. A zero-normal term has no
iterator with the counter $0$ and is atomic-seed normal, and (a) gives the
last identity.
(d) By induction on $C$. For each constructor, $\nf_{\mathrm a}$ of a node is
a function of the values of $\nf_{\mathrm a}$ at its children; for an
iterator the defining clause depends on $n'$, $a'$ and $h'$ only, and the
formula constructors are kept. The substitution commutes with the
constructors, because no free variable of $s$ or $s'$ is captured, and at
an occurrence of $x$ the two values are $\nf_{\mathrm a}(s)$ and
$\nf_{\mathrm a}(s')$.
\end{proof}

\begin{proposition}
\label{comp:prop:nfatomic}
$\nf_{\mathrm a}$ is a conversion normalizer and is not an admissible
normalizer: it violates ($\nu$4).
\end{proposition}

\begin{proof}
($\nu$1) holds by definition, ($\nu$2) and ($\nu$3) by
\cref{comp:lem:nfatomic}(a),(b), and ($\nu$5) by
\cref{comp:lem:nfatomic}(c) and \cref{comp:lem:conversion}. For ($\nu$4) let
$t\equiv\It{0}{x}{u}{0}$ and $\sigma\defeq[\Suc0/x]$. Then
\[
 \nf_{\mathrm a}(t[\sigma])\equiv\It{0}{\Suc0}{u}{0},
 \qquad
 \nf_{\mathrm a}\bigl(\nf_{\mathrm a}(t)[\nf_{\mathrm a}\circ\sigma]\bigr)
 \equiv\nf_{\mathrm a}(x[\Suc0/x])\equiv\Suc0,
\]
since $\Suc0$ is not atomic and $\nf_{\mathrm a}(t)\equiv x$. The two terms
are different; their equation is derivable in $\Nat$ by ($\nu$5), and it is
not an identity of syntax.
\end{proof}

The failure concerns substitution into a term that is not normal. If $t$ and
the terms $\sigma(x)$ are atomic-seed normal, both sides of ($\nu$4) are
$\nf_{\mathrm a}(t[\sigma])$. It matters for the rules: the calculus
$\Nat^{\nf_{\mathrm a}}$ defined as in \cref{comp:def:nukernel} with
$\nu\defeq\nf_{\mathrm a}$, whose metadata are normal, has only part of the
instances of $\Nat^{[\nf_{\mathrm a}]}$. For the raw body $B(y)\defeq(\It{0}{y}{u}{h}=y)$,
with $y$ not free in $h$, the normal body is $\nf_{\mathrm a}(B)\equiv(y=y)$,
whose instance at a target $t$ is $\nf_{\mathrm a}(t)=\nf_{\mathrm a}(t)$,
while the normal form of the raw instance $B(t)$ is
$\It{0}{\nf_{\mathrm a}(t)}{u}{\nf_{\mathrm a}(h)}=\nf_{\mathrm a}(t)$ when
$\nf_{\mathrm a}(t)$ is not atomic. The normal form of a premise therefore
does not determine the normal form of the conclusion; the raw body does.

\begin{remark}[Why the initializer is atomic]
\label{comp:rem:rebasing}
In the native syntax the binder $u$ of $\It{n}{a}{u}{h}$ covers $h$ only,
and a zero contraction returns the initializer without changing an index
(\cref{comp:sec:zero}). In an arity coding that places the initializer below
a slot, a contraction must rebase the initializer (see the remark after
\cref{comp:thm:zero}). \pkg{atomic\_seed\_refinement} gives an example for an
arity coding with de Bruijn levels, which it calls the compact
protected-level codec. A position has a depth. The children of an iterator array
$[n,a,h]$ at depth $d$ lie at depth $d+2$, and the slot of depth $d$ is the
state $u$ for $h$ and a dummy for $n$ and $a$. The child of a successor array
$[t]$ at depth $d$ lies at depth $d+4$, and the slots of depths $d$ and $d+2$
are dummies. The leaf $0$ is $0$, a leaf $k$ with $3\le k<d$ names the slot
of depth $k$, and a dummy reads as $0$. At depth $3$ the array
$[0,[[0],[0],[5]],0]$ is $\It{0}{\It{1}{1}{v}{\Suc v}}{u}{0}$, with value
$2$: the inner body $[5]$ lies at depth $7$, and its leaf, at depth $11$,
names the state $v$ of depth $5$. Replacing the outer array by its
initializer array gives $[[0],[0],[5]]$ at depth $3$. Its body lies at depth
$5$, its leaf at depth $9$ names the dummy of depth $5$ introduced by the
successor, and the array is $\It{1}{1}{v}{\Suc0}$, with value $1$. The
correct contraction is $[[0],[0],[3]]$, with value $2$. Both arrays are
well-formed integer and list data, so a successful traversal does not detect
the change of value. With an atomic initializer no binder below the
initializer moves: a leaf that names an outer slot keeps its name, and a leaf
that names the dummy of the contracted iterator denotes $0$. A correct sort
test and the handling of dummy slots are still required; the article does
not give a smaller implementation in such a coding.
\end{remark}

\subsection{Six rules}

\begin{definition}[The calculus $\mathsf{A}_6$]
\label{comp:def:asix}
$\mathsf{A}_6$ is the raw-metadata calculus over $\nf_{\mathrm a}$ with the
six rules $\rulename{Hyp}^{[\nf_{\mathrm a}]}$,
$\rulename{IS}^{[\nf_{\mathrm a}]}$, $\rulename{Tr}^{[\nf_{\mathrm a}]}$,
$\rulename{QE}^{[\nf_{\mathrm a}]}$, $\rulename{DE}^{[\nf_{\mathrm a}]}$ and
$\rulename{WI}^{[\nf_{\mathrm a}]}$. Explicitly, $\Sigma$ is a list of
atomic-seed normal formulas, $C$ is atomic-seed normal, $\overline{F}$ denotes
$\nf_{\mathrm a}(F)$, and the letters $n,a,h,s,t,r,H,B$ denote raw metadata:
\begin{gather*}
 \infrule{}{\Sigma\proves C}\ (C\in\Sigma),
 \qquad
 \infrule{}{\Sigma\proves
  \overline{\It{\Suc n}{a}{u}{h}=h\subst{\It{n}{a}{u}{h}}{u}}},
 \\[4pt]
 \infrule{\Sigma\proves\overline{s=t}\qquad\Sigma\proves\overline{H(s)=r}}
         {\Sigma\proves\overline{r=H(t)}},
 \qquad
 \infrule{\Sigma\proves\overline{\Q{y}B(y)}\qquad\Sigma\proves\overline{B(0)}}
         {\Sigma\proves\overline{B(t)}},
 \\[4pt]
 \infrule{\Sigma\proves\overline{\D{x}{y}B(x,y)}\qquad
          \Sigma,\overline{\Q{y}B(b,y)}\proves C}
         {\Sigma\proves C},
 \\[4pt]
 \infrule{\Sigma,\overline{B(t,\eta_{B,t}(z))}\proves\overline{B(t,\Suc z)}}
         {\Sigma\proves\overline{\D{x}{y}B(x,y)}}.
\end{gather*}
They are \rulename{Hyp}, \rulename{IS}, \rulename{Tr}, \rulename{QE},
\rulename{DE} and \rulename{WI}. Every bar is taken after the displayed raw
substitution. In \rulename{DE} the variable $b$ is not free in $\Sigma$, in
$C$ or in $\D{x}{y}B$. In \rulename{WI} the variable $z$ is not free in
$\Sigma$, in $t$ or in $\D{x}{y}B$, and $\eta_{B,t}(z)=z$ exactly when $B$ is
a literal equation and the raw witness $t$ is literally $0$, otherwise
$\eta_{B,t}(z)=0$. There is no restriction on $a$ in \rulename{IS} or on $s$
in \rulename{Tr}. There is no primitive zero computation, reflexivity,
symmetry, transitivity, general transport, induction or cut. The normalizer
and the raw metadata are part of the specification.
\end{definition}

\begin{lemma}[Reflexivity without zero computation]
\label{comp:lem:asrefl}
$\Sigma\proves_{\mathsf{A}_6}\overline{t=t}$ for every raw term $t$ and every
$\Sigma$.
\end{lemma}

\begin{proof}
Let $u$ be a variable not free in $t$ and put $L\defeq\It{\Suc0}{\Suc0}{u}{t}$.
The raw instance of \rulename{IS} at $n\equiv0$, $a\equiv\Suc0$ and $h\equiv t$
is $L=t\subst{\It{0}{\Suc0}{u}{t}}{u}\equiv(L=t)$, so \rulename{IS} gives
$\Sigma\proves\overline{L=t}$. \rulename{Tr} with $s\equiv L$, the target
$t$, the identity template $H(v)\equiv v$ and $r\equiv t$ has the premises
$\overline{L=t}$ and $\overline{H(L)=t}\equiv\overline{L=t}$, and the
conclusion $\overline{t=H(t)}\equiv\overline{t=t}$.
\end{proof}

The derivation does not use a zero computation, and its instance of
\rulename{IS} has a nonzero initializer. The body of $L$ is constant. The
term $\It{1}{t}{u}{u}$ of \cref{comp:lem:nzerorefl} does not serve here: the
right side of its instance of \rulename{IS} is $\It{0}{t}{u}{u}$, whose normal
form is $\It{0}{\nf_{\mathrm a}(t)}{u}{u}$ and not $\nf_{\mathrm a}(t)$ when
$\nf_{\mathrm a}(t)$ is not atomic.

\begin{lemma}[A reflexive closure]
\label{comp:lem:asqrefl}
$\Sigma\proves_{\mathsf{A}_6}\Q{y}(y=y)$ for every $\Sigma$.
\end{lemma}

\begin{proof}
Apply \rulename{WI} to the matrix $B(x,y)\defeq(y=y)$, with $x$ unused, at
the raw witness $0$. The matrix is a literal equation, so the mode is
progress, and the premise, with $z$ not free in $\Sigma$, is
$\Sigma,z=z\proves\Suc z=\Suc z$, which holds by \cref{comp:lem:asrefl}. This gives $\Sigma\proves\D{x}{y}(y=y)$. Apply
\rulename{DE} with a fresh $b$: the local assumption is
$\overline{\Q{y}B(b,y)}\equiv\Q{y}(y=y)$, and \rulename{Hyp} derives
$\Sigma,\Q{y}(y=y)\proves\Q{y}(y=y)$. The conclusion $\Q{y}(y=y)$ is closed,
so $b$ does not occur in it.
\end{proof}

\begin{theorem}[Zero computation from atomic contraction]
\label{comp:thm:asizero}
For all raw terms $t$, $h$ and every $\Sigma$,
\[
 \Sigma\proves_{\mathsf{A}_6}\overline{\It{0}{t}{u}{h}=t}.
\]
\end{theorem}

\begin{proof}
Let $y$ be a variable different from $u$ and not free in $t$ or $h$, and
apply \rulename{QE} with the raw body
\[
 B(y)\defeq\bigl(\It{0}{y}{u}{h}=y\bigr).
\]
Since the counter is $0$ and the normalized initializer $y$ is atomic,
\[
 \overline{\Q{y}B(y)}\equiv\Q{y}(y=y),
 \qquad
 \overline{B(0)}\equiv(0=0).
\]
The first premise is \cref{comp:lem:asqrefl} and the second is
\cref{comp:lem:asrefl}. At the target $t$ the conclusion is
$\overline{B(t)}$; the initializer lies outside the binder $u$ and $y$ is not
free in $h$, so $B(t)\equiv(\It{0}{t}{u}{h}=t)$.
\end{proof}

If the raw body were replaced by its normal form $y=y$ before the target is
substituted, the same inference would conclude $\overline{t=t}$ and not the
zero computation. The raw body of \rulename{QE} is therefore necessary for
this derivation; that a calculus with normal metadata does not derive the
zero computation in another way is not claimed.

\begin{lemma}[The two introductions]
\label{comp:lem:asintro}
The rules $\rulename{DI}^{[\nf_{\mathrm a}]}$ and
$\rulename{QI}^{[\nf_{\mathrm a}]}$ are admissible in $\mathsf{A}_6$.
\end{lemma}

\begin{proof}
An instance of $\rulename{DI}^{[\nf_{\mathrm a}]}$ is an instance of
\rulename{WI} unless its matrix $B$ is a literal equation and its raw witness
is literally $0$. In that case let
$\Sigma,\overline{B(0,0)}\proves\overline{B(0,\Suc z)}$ be derivable, and use
\rulename{WI} at the raw witness $j\defeq\It{0}{0}{u}{0}$. It is not
literally $0$, so the mode is case. Since $\nf_{\mathrm a}(j)\equiv0$ and $j$
is closed, \cref{comp:lem:nfatomic}(d) gives
$\overline{B(j,0)}\equiv\overline{B(0,0)}$ and
$\overline{B(j,\Suc z)}\equiv\overline{B(0,\Suc z)}$, so the premise of
\rulename{WI} is the given sequent, and the conclusion
$\overline{\D{x}{y}B}$ does not depend on the witness. The variable $z$ is
fresh for $j$, which is closed. No derivation is rewritten and no equation is
transported.

For $\rulename{QI}^{[\nf_{\mathrm a}]}$ let $E(y)$ be a literal equation and
$\Sigma,\overline{E(z)}\proves\overline{E(\Suc z)}$ derivable, with $z$ not
free in $\Sigma$ or in $\Q{y}E$. Let $x$ be a variable not free in $E$ and put
$B(x,y)\defeq E(y)$. \rulename{WI} at the raw witness $0$ is in progress mode,
and its premise $\Sigma,\overline{B(0,z)}\proves\overline{B(0,\Suc z)}$ is the
given sequent; it derives $\Sigma\proves\overline{\D{x}{y}E(y)}$.
\rulename{DE} at a fresh $b$ has the local assumption
$\overline{\Q{y}B(b,y)}\equiv\overline{\Q{y}E(y)}$, and \rulename{Hyp}
derives the conclusion $\overline{\Q{y}E(y)}$ from it; $b$ is not free in
$\Sigma$, in that conclusion or in $\D{x}{y}E$.
\end{proof}

The alias $j$ is the raw alias of \cref{kern:lem:alias}. On zero-normal terms
with normal metadata it is unavailable, and \cref{comp:lem:single}(ii) uses
the alias $\It{1}{0}{u}{0}$ together with transport along its equation with
$0$ (\cref{comp:thm:zsix}). With raw metadata the witness keeps its raw shape
for the mode test, and its normal form makes the normalized branch literally
the branch at $0$.

\begin{theorem}[Six rules with atomic-seed contraction]
\label{comp:thm:asix}
$\mathsf{A}_6$ derives exactly the sequents over $\nf_{\mathrm a}$ that $\Nat$
derives. For every finite context $\Gamma$ of raw native formulas and every
raw native formula $A$,
\[
 \Gamma\proves_{\Nat}A
 \quad\Longleftrightarrow\quad
 \nf_{\mathrm a}(\Gamma)\proves_{\mathsf{A}_6}\nf_{\mathrm a}(A),
\]
effectively on finite derivations, and \cref{comp:cor:rawbridge} applies to
$\mathsf{A}_6$. In particular $\mathsf{A}_6$ is equiconsistent and
bi-interpretable with $\iPRA$, with whole numerical domains, equality
interpreted by equality, and numerical composites with the graph $x=y$.
\end{theorem}

\begin{proof}
$\mathsf{A}_6$ is a raw-metadata calculus over the conversion normalizer
$\nf_{\mathrm a}$ (\cref{comp:prop:nfatomic}). By
\cref{comp:thm:rawtransfer}(b) it suffices that each of the eight rules of
$\Nat^{[\nf_{\mathrm a}]}$ is admissible in $\mathsf{A}_6$. Five of them are
rules of $\mathsf{A}_6$. An instance of $\rulename{I0}^{[\nf_{\mathrm a}]}$
is $\Sigma\proves\overline{\It{0}{a}{u}{h}=a}$ for raw $a,h$, which is
\cref{comp:thm:asizero}. $\rulename{DI}^{[\nf_{\mathrm a}]}$ and
$\rulename{QI}^{[\nf_{\mathrm a}]}$ are admissible by \cref{comp:lem:asintro}.
The last statements are \cref{comp:cor:rawbridge} and statement (c) of
\cref{comp:cor:bridge}; the contradiction $0=\Suc0$ is atomic-seed normal.
\end{proof}

\begin{remark}
\label{comp:rem:asix}
\pkg{atomic\_seed\_refinement} proves the equivalence directly against
$\Nat$, in place of \cref{comp:thm:rawtransfer}: a derivation of $\Nat$ is
translated rule by rule with its raw metadata kept, and a derivation of
$\mathsf{A}_6$ is simulated in $\Nat$ by converting each normalized premise
to its raw form, which uses that every raw term is provably equal to its
partial normal form, by the zero computation and the congruence of iterator
bodies (\cref{comp:lem:conversion}). Both directions of the comparison use
only the derived rules of $\Nat$, so the argument is not circular. The
package composes the equivalence with the bridge of \cref{base:thm:bridge}
for the bi-interpretation with $\iPRA$, as in \cref{comp:cor:rawbridge}. The
package also observes that \rulename{IS} can be restricted to raw
initializers that are not literally $0$, and \rulename{Tr} to raw sources
that are not literally $0$, without changing a normalized instance: by
\cref{comp:lem:nfatomic}(d) the raw term $\It{0}{0}{u}{0}$ in place of an
excluded $0$ gives the same normalized premises and conclusion. This
concerns tests on raw data; it does not justify a restriction on normalized
initializers or sources, and the primary calculus has neither restriction.
Compared with $\mathsf{Z}_6$, $\mathsf{A}_6$ has the same six rule names, a
weaker normalizer and raw metadata in every rule; the zero computation is
not normalized away but derived. The calculus is not formalized in Lean; the
package records finite tests of its reference checker
(\cref{comp:sec:sizes}).
\end{remark}

%%%%%%%%%%%%%%%%%%%%%%%%%%%%%%%%%%%%%%%%%%%%%%%%%%%%%%%%%%%%%%%%%%%%%%%%%%%%
\section{Full normalization}
\label{comp:sec:full}

This section draws on \pkg{full\_normalization} and
\pkg{computation\_normal}. The two packages arrive at the same six-rule
calculus on full normal forms; they differ in the proof of termination, in
the repair of the single introduction and in their measurements. Both
computations are moved into the syntax, and no computation rule remains.

\subsection{Termination and unique normal forms}

The reduction $\to$ is that of \cref{comp:def:reduction}. It is stable under
substitution: if $t\to t'$, then $t[\sigma]\to t'[\sigma]$; and if
$\sigma(x)\to^{*}\sigma'(x)$ for every $x$, then
$t[\sigma]\to^{*}t[\sigma']$. A term is \emph{strongly normalizing} if it
has no infinite reduction sequence. A term has finitely many one-step
reducts, so for a strongly normalizing term $t$ the reduction sequences from
$t$ have a maximal length $d(t)$. The number of nodes of $t$ is $|t|$.

\begin{lemma}[Iterator closure]
\label{comp:lem:snclosure}
Let $n$ and $a$ be strongly normalizing, and let $h\subst{v}{u}$ be strongly
normalizing for every strongly normalizing term $v$. Then $\It{n}{a}{u}{h}$
is strongly normalizing.
\end{lemma}

\begin{proof}
With $v\equiv u$ the body $h$ is strongly normalizing. Argue by induction on
the triple $(d(n),|n|,d(a)+d(h))$ in the lexicographic order; it suffices to
show that every one-step reduct of $\It{n}{a}{u}{h}$ is strongly
normalizing. A reduction in $n$ decreases the first component. A reduction
in $a$ or in $h$ leaves the first two components and decreases the third;
if $h\to h'$, then $h'\subst{v}{u}$ is a reduct of $h\subst{v}{u}$ for every
$v$, so the hypothesis on the body persists. If $n\equiv0$, the root reduct
is $a$. If $n\equiv\Suc p$, then $d(p)=d(n)$ and $|p|<|n|$, so
$\It{p}{a}{u}{h}$ is strongly normalizing by the induction hypothesis, and
the root reduct $h\subst{\It{p}{a}{u}{h}}{u}$ is strongly normalizing by
the hypothesis on the body.
\end{proof}

\begin{theorem}[Strong normalization]
\label{comp:thm:sn}
If $\sigma$ assigns strongly normalizing terms to the variables, then
$t[\sigma]$ is strongly normalizing for every raw term $t$. In particular
every raw term is strongly normalizing.
\end{theorem}

\begin{proof}
By induction on $t$. Variables, $0$ and successors are immediate. For
$\It{n}{a}{u}{h}$, with $u$ renamed away from $\sigma$, the instances
$n[\sigma]$ and $a[\sigma]$ are strongly normalizing by the induction
hypotheses. The induction hypothesis for $h$ applies to every extension of
$\sigma$ by an assignment of a strongly normalizing $v$ to $u$; this is the
hypothesis on the body in \cref{comp:lem:snclosure} for the term
$\It{n[\sigma]}{a[\sigma]}{u}{h[\sigma]}$. The last statement is the case
of the identity substitution.
\end{proof}

\begin{theorem}[Confluence]
\label{comp:thm:confluence}
The reduction $\to$ is confluent. Hence every raw term $t$ has a unique
normal form $\nf(t)$, up to the renaming of bound variables.
\end{theorem}

\begin{proof}
Define the \emph{parallel reduction} $\Rightarrow$ inductively: $x\Rightarrow
x$; $0\Rightarrow0$; if $t\Rightarrow t'$ then $\Suc t\Rightarrow\Suc t'$;
and if $n\Rightarrow n'$, $a\Rightarrow a'$ and $h\Rightarrow h'$, then
\[
 \It{n}{a}{u}{h}\Rightarrow\It{n'}{a'}{u}{h'},\qquad
 \It{0}{a}{u}{h}\Rightarrow a',\qquad
 \It{\Suc n}{a}{u}{h}\Rightarrow h'\subst{\It{n'}{a'}{u}{h'}}{u}.
\]
Then $t\Rightarrow t$, and ${\to}\subseteq{\Rightarrow}\subseteq{\to^{*}}$.
By induction on $h\Rightarrow h'$, if also $r\Rightarrow r'$ then
$h\subst{r}{u}\Rightarrow h'\subst{r'}{u}$; the case of a contracted
successor redex uses the composition law of substitution.

Define the \emph{complete development} $t^{\bullet}$ by
$x^{\bullet}\defeq x$, $0^{\bullet}\defeq0$,
$(\Suc t)^{\bullet}\defeq\Suc t^{\bullet}$,
$(\It{0}{a}{u}{h})^{\bullet}\defeq a^{\bullet}$,
$(\It{\Suc n}{a}{u}{h})^{\bullet}\defeq
h^{\bullet}\subst{\It{n^{\bullet}}{a^{\bullet}}{u}{h^{\bullet}}}{u}$, and
$(\It{n}{a}{u}{h})^{\bullet}\defeq
\It{n^{\bullet}}{a^{\bullet}}{u}{h^{\bullet}}$ for every other $n$. The
claim is that $t\Rightarrow t'$ implies $t'\Rightarrow t^{\bullet}$; it is
proved by induction on $t\Rightarrow t'$. Consider
$\It{n}{a}{u}{h}\Rightarrow\It{n'}{a'}{u}{h'}$, with
$n'\Rightarrow n^{\bullet}$, $a'\Rightarrow a^{\bullet}$ and
$h'\Rightarrow h^{\bullet}$ by the induction hypotheses. If $n\equiv0$, then
$n'\equiv0$ and $\It{0}{a'}{u}{h'}\Rightarrow a^{\bullet}$. If
$n\equiv\Suc m$, then $n'\equiv\Suc m'$ with $m'\Rightarrow m^{\bullet}$, and
$\It{\Suc m'}{a'}{u}{h'}\Rightarrow
h^{\bullet}\subst{\It{m^{\bullet}}{a^{\bullet}}{u}{h^{\bullet}}}{u}$.
Otherwise the congruence rule applies. For
$\It{0}{a}{u}{h}\Rightarrow a'$ the induction hypothesis is
$a'\Rightarrow a^{\bullet}$. For
$\It{\Suc n}{a}{u}{h}\Rightarrow h'\subst{\It{n'}{a'}{u}{h'}}{u}$ the
substitution property, applied to $h'\Rightarrow h^{\bullet}$ and
$\It{n'}{a'}{u}{h'}\Rightarrow\It{n^{\bullet}}{a^{\bullet}}{u}{h^{\bullet}}$,
gives the reduct $(\It{\Suc n}{a}{u}{h})^{\bullet}$. The remaining cases are
immediate.

By the claim, two parallel reducts of $t$ both reduce in parallel to
$t^{\bullet}$, so $\Rightarrow$ has the diamond property, and its transitive
closure $\to^{*}$ is confluent. With \cref{comp:thm:sn}, every term has
exactly one normal form.
\end{proof}

\begin{remark}
The proof of \cref{comp:thm:sn} is that of \pkg{full\_normalization}, which
obtains confluence from local confluence and termination.
\pkg{computation\_normal} obtains termination by an embedding into the
simply typed calculus with a recursor at the type of numbers (G\"odel's
system $T$): the iterator is translated as a recursor whose step ignores its
counter, a zero contraction becomes one recursor contraction, a successor
contraction becomes one recursor contraction followed by two
$\beta$-contractions, and the translation respects contexts and
substitution, so an infinite reduction would contradict the strong
normalization of $T$, for which the package cites the mechanized proof of
S.~Urciuoli (EPTCS 376, 2023); it obtains confluence by parallel reduction,
as above. The Lean formalization defines the normalizer by evaluation: the term
is evaluated in an environment of terms by structural recursion
(\Lean{Tm.nbe}), the iterator by a recursion on the value of its counter
that unfolds a numeral prefix and leaves a neutral counter in place
(\Lean{Tm.recV}), and \Lean{Tm.full} is the value in the identity
environment. Totality and uniqueness of normal forms then need no proof.
The system $T$ and the evaluation are used in the metatheory; neither is
added to the arithmetic.
\end{remark}

\begin{lemma}[Full normal forms]
\label{comp:lem:fullgrammar}
The normal terms are generated by
\[
 \mathit{ne}::=x\mid\It{\mathit{ne}}{\mathit{nt}}{u}{\mathit{nt}},\qquad
 \mathit{nt}::=\mathit{ne}\mid0\mid\Suc\,\mathit{nt}.
\]
The terms $\mathit{ne}$ are the \emph{neutral} terms. The counter of every
iterator of a normal term is neutral; it has neither $0$ nor $\Suc$ at its
root. Every closed normal term is a numeral.
\end{lemma}

\begin{proof}
A term is normal exactly when its subterms are normal and no contraction
applies at its root, which gives the grammar. A neutral term has a free
variable: following the counters from its root leads, outside every binder,
to a variable. Hence a closed normal term is built from $0$ and $\Suc$.
\end{proof}

For example $\It{x}{0}{u}{0}$ is normal, although its value is $0$ for every
$x$, while $\It{\Suc x}{0}{u}{0}$ has the normal form $0$. Normalization
does not decide equality.

\subsection{The kernel without computation rules}

\begin{proposition}
\label{comp:prop:nfadm}
$\nf$ is an admissible normalizer.
\end{proposition}

\begin{proof}
($\nu$1): variables and $0$ are normal, and the reducts of $\Suc t$ are the
terms $\Suc t'$ with $t'$ a reduct of $t$. ($\nu$2): reduction commutes with
renaming, the renaming of a normal term is normal by
\cref{comp:lem:fullgrammar}, and a contraction creates no free variable.
($\nu$3) is immediate. ($\nu$4): by the two stability properties,
$t[\sigma]\to^{*}\nf(t)[\sigma]\to^{*}\nf(t)[\nf\circ\sigma]$, so both
sides of ($\nu$4) are the normal form of $t[\sigma]$. ($\nu$5) is
\cref{comp:lem:conversion}.
\end{proof}

In the rest of this section terms and formulas are in full normal form and
$t\langle\sigma\rangle$ is $\nf(t[\sigma])$. An expression
$\It{n}{a}{u}{h}$ in a rule denotes the normal form of the displayed term,
also when a substitution exposes a redex. Formula normalization changes
terms and never removes $\Q{y}$ or $\D{x}{y}$.

\begin{definition}[The kernel $\Nat^{\mathrm{f}}$]
\label{comp:def:nfull}
$\Nat^{\mathrm{f}}$ has native formulas over full normal forms, normalizing
substitution, and the seven rules \rulename{Hyp}, \rulename{Refl} (for every
normal term), \rulename{Tr}, \rulename{QE}, \rulename{DE}, \rulename{DI} and
\rulename{QI}.
\end{definition}

\begin{theorem}[Both computation rules are deleted]
\label{comp:thm:full}
For every finite context $\Gamma$ of raw native formulas and every raw
native formula $A$,
\[
 \Gamma\proves_{\Nat}A
 \quad\Longleftrightarrow\quad
 \nf(\Gamma)\proves_{\Nat^{\mathrm f}}\nf(A).
\]
$\Nat^{\mathrm f}$ is closed under weakening, cut and normalizing
substitution, and \cref{comp:cor:bridge} applies to it.
\end{theorem}

\begin{proof}
By \cref{comp:prop:nfadm} and \cref{comp:thm:transfer} it suffices that
$\Nat^{\mathrm f}$ and $\Nat^{\nf}$ have the same rules. The left side of
an instance $e$ of \rulename{I0} or \rulename{IS} reduces to its right side,
so $\nf(e)$ is literally $r=r$ for a normal $r$, an instance of
\rulename{Refl}. Conversely, for a normal $r$ the reflexivity $r=r$ is
$\nf(e)$ for the instance $\It{0}{r}{u}{0}=r$ of \rulename{I0}.
\end{proof}

\begin{remark}
The theorem is formalized as \Lean{fnprf\_iff\_nprf} for the calculus
\Lean{FNPrf}; the laws of the normalizer are \Lean{Tm.full\_subst\_congr}
and \Lean{NFm.full\_subst\_congr}, and ($\nu$5) is \Lean{NPrf.eqFull}. Both
packages state the theorem with the auxiliary seven-rule calculus as a
bridge to their six rules, and both stress that the comparison uses
properties of the raw calculus that are established before it: the equality
rules, the Leibniz rule, cut and equation induction of $\Nat$.
\end{remark}

\subsection{Six rules on full normal forms}

\begin{definition}[The calculus $\mathsf{F}$]
\label{comp:def:fsix}
$\mathsf{F}$ has the syntax and the substitution of $\Nat^{\mathrm f}$ and
the six rules \rulename{Hyp}, \rulename{Refl}, \rulename{Tr},
\rulename{QE}, \rulename{DE} and \rulename{WI}. The template of
\rulename{Tr} is any normal term with one distinguished variable, used any
number of times; there is no restriction on the terms of its first premise.
The test $t\equiv0$ of \rulename{WI} is applied to the normal witness.
\end{definition}

There is no computation axiom: both equations are enforced by the
normalizing substitution. Weakening, injective renaming and cut are
admissible in $\mathsf{F}$ by direct induction on derivations, since they
preserve normality, the literal $0$, the root of a matrix and freshness, and
no guard depends on the context. The equality rules and transport in literal
equations are derivable by \cref{comp:lem:eqalg}(b), the instances matching
by \cref{comp:lem:laws}.

On full normal forms \cref{comp:lem:single}(ii) has no closed term $Z$ to
use: by \cref{comp:lem:fullgrammar} a closed normal term is a numeral, and
the aliases of \cref{comp:lem:alias} have the normal form $0$. The packages
give two repairs.

\begin{lemma}[A fresh name for zero]
\label{comp:lem:freshzero}
Let $Z(x,y)$ be the equation $\It{y}{0}{u}{x}=0$. Then
$\proves_{\mathsf F}\D{x}{y}Z(x,y)$, and
$\Q{y}Z(b,y)\proves_{\mathsf F}b=0$ for a variable $b$.
\end{lemma}

\begin{proof}
Introduce $\D{x}{y}Z$ by \rulename{WI} at the witness $0$. The matrix is a
literal equation, so the mode is progress. The conclusion of the premise is
$Z(0,\Suc z)$, the normal form of $\It{\Suc z}{0}{u}{0}=0$, which is $0=0$;
reflexivity derives it without the assumption. For the second statement the
base $Z(b,0)$ is the normal form of $\It{0}{0}{u}{b}=0$, which is $0=0$.
\rulename{QE} with this base, derived by reflexivity, at the target $1$
gives $Z(b,1)$, the normal form of $\It{1}{0}{u}{b}=0$, which is $b=0$.
\end{proof}

The witness opened by \rulename{DE} is a variable, hence not literally $0$,
and it is provably $0$. It is discharged before the conclusion; no constant
and no element of the domain is added.

\begin{theorem}
\label{comp:thm:fsix}
For every context $\Gamma$ and formula $A$ in full normal form,
\[
 \Gamma\proves_{\mathsf F}A
 \quad\Longleftrightarrow\quad
 \Gamma\proves_{\Nat^{\mathrm f}}A .
\]
Hence $\Gamma\proves_{\Nat}A$ if and only if
$\nf(\Gamma)\proves_{\mathsf F}\nf(A)$ for raw $\Gamma,A$, and substitution
and the Leibniz rule are admissible in $\mathsf{F}$.
\end{theorem}

\begin{proof}
From $\Nat^{\mathrm f}$ to $\mathsf{F}$: \rulename{QI} is admissible by
\cref{comp:lem:single}(i). An instance of \rulename{DI} is an instance of
\rulename{WI} unless its matrix is a literal equation $E(x,y)$ and its
witness is $0$. In that case let $\Gamma,E(0,0)\proves E(0,\Suc z)$ be
given. Choose $b$ fresh and work in the context
$\Gamma'\defeq\Gamma,\Q{y}Z(b,y)$, where $\Gamma'\proves b=0$ by
\cref{comp:lem:freshzero}. By \cref{comp:lem:single}(ii) with the term $b$,
applied to the weakened premise, $\Gamma'\proves\D{x}{y}E$. \rulename{DE}
with the first premise $\proves\D{x}{y}Z$ closes the scope of $b$, which
occurs neither in $\Gamma$ nor in $\D{x}{y}E$. From $\mathsf{F}$ to
$\Nat^{\mathrm f}$, \rulename{WI} is admissible by
\cref{comp:lem:single}(iii). The remaining statements are
\cref{comp:thm:full,comp:cor:bridge}.
\end{proof}

Only transport in literal equations, weakening, injective renaming and cut
are used in the repair; the Leibniz rule and the substitution of terms in
derivations of $\mathsf{F}$ are consequences of the theorem and not
premises of it.

\begin{lemma}[The compound envelope]
\label{comp:lem:envelope}
For a literal equation $E\equiv(s=t)$ in full normal form put
\[
 G(v,E)\defeq\bigl(\It{v}{0}{u}{s}=\It{v}{0}{u}{t}\bigr),\qquad
 L(E)\defeq\Q{v}G(v,E),
\]
with $u,v$ not free in $E$. Then $G(0,E)\equiv(0=0)$ and
$G(\Suc w,E)\equiv E$, and for every context $\Gamma$ both
$\Gamma,E\proves_{\mathsf F}L(E)$ and $\Gamma,L(E)\proves_{\mathsf F}E$.
\end{lemma}

\begin{proof}
The two identities are computations of normal forms; in the second, $u$ is
not free in $s$ and $t$. For the first derivation apply the closure by cases
of \cref{comp:lem:single}(i) to the matrix $G(v,E)$: its premise asks for
$G(\Suc z,E)$, which is the assumption $E$. For the second apply
\rulename{QE} to $L(E)$ with the base $0=0$ at the target $1$. Both are
uniform in the parameters of $E$ and use no induction.
\end{proof}

The envelope yields a second proof of the missing case in
\cref{comp:thm:fsix}, without \cref{comp:lem:freshzero}. Let
$\Gamma,E(0,0)\proves E(0,\Suc z)$ and put $F(x,y)\defeq L(E(x,y))$, a
matrix with the root $\Q{v}$. In the context $\Gamma,F(0,0)$ the envelope
gives $E(0,0)$, cut gives $E(0,\Suc z)$, and the envelope gives
$F(0,\Suc z)$. \rulename{WI} at the witness $0$ is in case mode, because
$F$ is not a literal equation, and derives $\D{x}{y}F$. Open it by
\rulename{DE} at a fresh $b$, retaining $\Q{y}F(b,y)$. To introduce
$\D{x}{y}E$ at the witness $b$, in case mode since a variable is not $0$,
assume $E(b,0)$; the envelope gives $F(b,0)$, \rulename{QE} gives
$F(b,\Suc z)$, and the envelope gives $E(b,\Suc z)$. \rulename{DE} closes
the derivation.

\begin{remark}
The fresh-zero lemma is the repair of \pkg{computation\_normal}, and the
compound envelope that of \pkg{full\_normalization}. The Lean formalization
of the six rules, \Lean{F6Prf} with \Lean{f6prf\_iff\_nprf}, uses a third
repair, \Lean{F6Prf.dIZero}: with $T(x)\defeq\sel(x,0,0)$, the formula
$\D{x}{y}B(T(x),y)$ is introduced at the witness $1$, in case mode, from
the given premise, because $T(1)$ has the normal form $0$; it is opened at
a fresh $b$; and $\D{x}{y}B$ is introduced at the neutral witness $T(b)$
from the opened closure by \rulename{QE}.
\end{remark}

\begin{example}
\label{comp:ex:modechange}
The formula $\D{x}{y}(y=1)$ is derivable: under its reading the antecedent
is $0=1$. On raw terms it is introduced in case mode at the witness
$\It{0}{0}{u}{0}$, from the premise $0=1\proves\Suc z=1$. Full
normalization turns this witness into $0$, and the same inference would be
in progress mode, with the premise $z=1\proves\Suc z=1$, which is not
derivable. A derivation normalized rule by rule is therefore rejected, and
the derivation produced by either repair is accepted. The package
\pkg{full\_normalization} ships both certificates.
\end{example}

\subsection{What normalization costs}

\begin{proposition}[No primitive recursive uniform normalizer]
\label{comp:prop:notpr}
The map $t\mapsto\nf(t)$ on codes of raw terms is total recursive and not
primitive recursive.
\end{proposition}

\begin{proof}
Totality is \cref{comp:thm:sn,comp:thm:confluence}. Suppose that the map is
primitive recursive. Enumerate the raw terms with one free variable,
replacing a code that is not well formed by the code of the term $0$. The
check of well-formedness, the formation of the closed instance of a term at
a numeral and the reading of the value of a numeral are primitive recursive.
By \cref{comp:lem:fullgrammar} and ($\nu$5) the normal form of a closed
instance is the numeral of its value. Hence the function $V(e,n)$, the
value at $n$ of the term with the code $e$, is primitive recursive. By
\cref{base:thm:primrec} every primitive recursive unary function is the
value function of such a term. In particular $n\mapsto V(n,n)+1$ has a code
$e$, and $V(e,e)=V(e,e)+1$, a contradiction.
\end{proof}

Each fixed term denotes a primitive recursive function; the obstruction is
to one evaluator for all descriptions. A checker for $\mathsf{F}$ or
$\Nat^{\mathrm f}$ is a total recursive program that contains a normalizer,
and its running time and the size of its intermediate terms are not bounded
by a primitive recursive function of the input: by
\cref{comp:lem:fullgrammar} every closed term is expanded to a numeral, so
a short closed expression can have a normal form of enormous size. No
primitive recursive bound and no proof in $\iPRA$ of a uniform normalization
theorem is claimed. The statements of this section are external statements
about each finite derivation; both interpretations of
\cref{comp:cor:bridge} act on fixed finite expressions.

\begin{proposition}[Fuel]
\label{comp:prop:fuel}
Let the checker of $\mathsf{F}$ be modified so that a certificate carries a
number $f$, each successor contraction spends one unit, and exhaustion
rejects. The relation ``the certificate with fuel $f$ is accepted'' is
primitive recursive, and a formula has an accepted certificate for some
fuel exactly when it has an accepted certificate of the unbounded checker.
\end{proposition}

\begin{proof}
Let $m$ be the size of the certificate. An instance of a term or formula of
size at most $M$ at terms of size at most $M$ has size at most $M^{2}$. A
successor contraction replaces a term of size at most $M$ by a term of size
at most $M^{2}$, since the body and the recursive iterator are subterms of
the redex and the iterator is copied once for each occurrence of the bound
variable. A zero contraction and the traversals between contractions do not
increase the size, and at most $M$ zero contractions occur in a row. The
checker performs at most $m$ inferences, each with a number of
substitutions bounded by a constant $k$, and at most $f$ successor
contractions. Hence every term that occurs has size at most
$m^{2^{km+f}}$, and the length of the computation is bounded by a primitive
recursive function of $m$ and $f$; a computation with such bounds on its
length and on the size of its configurations is primitive recursive.
% ed.: computation_normal states the primitive recursive bound in one
% sentence; the squaring argument is supplied here.
An accepting run of the unbounded checker is finite and performs finitely
many successor contractions, so it is accepted with that much fuel; a run
accepted with fuel is a run of the unbounded checker.
\end{proof}

\pkg{computation\_normal} ships the fuelled checker and measures it
(\cref{comp:sec:sizes}); \pkg{full\_normalization} observes that a relation
of this kind can also be obtained from explicit reduction traces checked
step by step, and does not build that checker.

%%%%%%%%%%%%%%%%%%%%%%%%%%%%%%%%%%%%%%%%%%%%%%%%%%%%%%%%%%%%%%%%%%%%%%%%%%%%
\section{A smaller checker for the unchanged core}
\label{comp:sec:scope}

This section draws on the second result of \pkg{computation\_normal}, which
concerns a program and no calculus. The rules, the certificates and the
kernel function of the checker are those of the package
\pkg{validated\_presentation}; only the validator of the public goal
changes, by a composition of that validator with the reuse of the scope
traversal of the package \pkg{scope\_kernel}.

The input is a JSON array. Its last element is the closed public goal, in
the arity coding of the native formulas, and the earlier elements are proof
packets consumed from the end; an unused prefix is allowed. A lexical gate
accepts the whole input only if every value in it is an integer or a list.
The public validator of the baseline carries a depth $d$, even at formula
positions and odd at numerical positions; a node with $r$ children changes
the depth by $5-r$, a formula node has one, two or three children, of which
the node with two children is the equation and has numerical children, and
a numerical node has one or three children. A numerical leaf $a$ is accepted
if $d$ is odd and $a=0$ or $a$ is odd with $2<a<d$. The kernel, before it
reads any packet, opens the goal by a scope traversal whose substitution
argument is, in this initial call, the kernel function itself and not a
term: at a leaf within the scope the traversal returns the leaf, and at a
leaf outside the scope it attempts to use the substitution argument as a
term and fails.

The new validator keeps only the Boolean ``numerical position'', which
becomes true below a node with two children and is otherwise inherited, and
accepts a numerical leaf $a$ if $a=0$ or $a$ is odd with $a>2$.

\begin{theorem}[Same accepted inputs]
\label{comp:thm:scope}
In the execution model without resource limits, the new checker and the
baseline accept exactly the same inputs.
\end{theorem}

\begin{proof}
The lexical gate is the same. On inputs of integers and lists, induction on
the tree shows that the Boolean of the new validator is the parity of the
depth of the old one wherever the shape is valid, by the table of arities
above. All old tests remain except the upper bound $a<d$ at numerical
leaves, so every goal accepted by the old validator is accepted by the new
one. Let a goal be accepted by the new validator and violate an upper
bound. The kernel is unchanged and computes the initial opening of the goal
before it tests membership in the context or reads a packet. Every
numerical position has an odd depth at least three, and the leaf is $0$ or
an odd integer greater than two; at the violating leaf the opening uses its
substitution argument, which is a function, as a term, and the operation
raises an exception. An exception is a rejection in both programs, and no
packet can prevent it. A formula-position integer is rejected by both
validators. On goals accepted by the old validator the initial opening is
the identity, and all later steps of the two programs are the same.
\end{proof}

Hence every soundness, completeness and interpretation result for the
certificate system of the baseline holds for the new checker. The result is
an implementation change for an unchanged calculus; its sizes are in
\cref{comp:sec:sizes}.

%%%%%%%%%%%%%%%%%%%%%%%%%%%%%%%%%%%%%%%%%%%%%%%%%%%%%%%%%%%%%%%%%%%%%%%%%%%%
\section{Checker sizes}
\label{comp:sec:sizes}

This section collects the measurements of all ten packages, by the protocol
of \cref{base:sec:measure}. Every row of \cref{comp:tab:sizes} is a complete
standalone program, with its input handling, its validation of syntax and
scope, its substitution and normalization, all rules and its output. The
gzip column, at level 9 with the timestamp zero, is reported by the packages
for information and is not an objective. The reference interpreter is
CPython 3.12.14. \pkg{uniform\_step}, \pkg{seed\_evaluation} and
\pkg{initializer\_equality} remove docstrings before normalization; the
other packages do not. A dash marks a number that the package does not
state.

\begin{longtable}{@{}p{0.55\linewidth}rrrr@{}}
\caption{Complete checkers and their matched controls, as stated by the
packages.}\label{comp:tab:sizes}\\
\toprule
Complete program & Tokens & AST nodes & Bytes & Gzip-9\\
\midrule
\endfirsthead
\toprule
Complete program & Tokens & AST nodes & Bytes & Gzip-9\\
\midrule
\endhead
\bottomrule
\endfoot
\multicolumn{5}{@{}l}{\pkg{computational\_equality} (tagged, native contexts)}\\
Archived seven-rule control & 1024 & 1080 & 3145 & 964\\
Matched seven-rule control & 1008 & 1064 & 3121 & 957\\
$\mathsf{E}_3$ & 959 & 1026 & 3039 & 966\\
\midrule
\multicolumn{5}{@{}l}{\pkg{equality\_from\_iteration} (same format)}\\
$\mathsf{E}_3$, upstream & 959 & 1026 & 3039 & 966\\
$\mathsf{E}_2$ & 954 & 1019 & 2996 & 955\\
Zero-normal control with \rulename{Root} kept & 987 & 1059 & 3064 & 982\\
$\mathsf{U}$ & 979 & 1047 & 3011 & 965\\
\midrule
\multicolumn{5}{@{}l}{\pkg{uniform\_step} (clause contexts)}\\
Matched zero-normal control, computation and transport & 860 & 949 & 3048 & 1018\\
Matched raw control, three-case \rulename{Eval} & 853 & 969 & 3094 & 1026\\
Incoming canonical clausal checker & 946 & 1072 & 3305 & 1073\\
$\mathsf{C}_{\mathrm{Step}}$ & 829 & 902 & 2950 & 1005\\
Native contexts: matched zero-normal control & 1010 & 1080 & 3118 & --\\
Native contexts: \rulename{Refl} and \rulename{Step} & 981 & 1051 & 3029 & --\\
\midrule
\multicolumn{5}{@{}l}{\pkg{seed\_evaluation} (clause contexts)}\\
Computation and transport control & 860 & 949 & 3048 & --\\
\rulename{Step}, helper style & 829 & 902 & 2950 & --\\
\rulename{SeedEval}, helper style & 829 & 902 & 2950 & --\\
\rulename{Step}, shared forward traversal & 758 & 815 & 2804 & --\\
\rulename{SeedEval}, shared forward traversal & 756 & 815 & 2802 & --\\
\rulename{Step}, goal-directed & 741 & 826 & 2283 & --\\
\rulename{SeedEval}, goal-directed & 741 & 826 & 2283 & --\\
\rulename{Step}, shared goal-directed traversal & 709 & 787 & 2265 & --\\
\rulename{SeedEval}, shared goal-directed traversal & 709 & 787 & 2265 & --\\
\rulename{App} with the constructor $U$, shared forward traversal & 803 & 864 & 2902 & --\\
\midrule
\multicolumn{5}{@{}l}{\pkg{zero\_normal} (tagged, native contexts)}\\
Raw, two introductions (8 rules) & 1084 & 1141 & 3322 & 984\\
$\Nat^{0}$ (7 rules) & 1073 & 1136 & 3272 & 996\\
Raw, single introduction (7 rules) & 1024 & 1080 & 3145 & 964\\
$\mathsf{Z}_6$ (6 rules) & 1013 & 1075 & 3095 & 977\\
\midrule
\multicolumn{5}{@{}l}{\pkg{reflexive\_base} (same format)}\\
$\Nat^{0}$ (7 rules) & 1073 & 1136 & 3272 & 996\\
$\mathsf{Z}_6$, the witness-selected control (6 rules) & 1013 & 1075 & 3095 & 977\\
$\mathsf{R}$ (6 rules) & 1013 & 1073 & 3084 & 981\\
\midrule
\multicolumn{5}{@{}l}{\pkg{initializer\_equality} (closed goals)}\\
One-template control & 1202 & 1360 & 4029 & 1173\\
$\mathsf{S}$, selected & 1192 & 1347 & 3979 & 1175\\
$\mathsf{S}$, explicit shape tests & 1236 & 1404 & 4156 & 1226\\
Canonical recursor checker, inherited & 1338 & 1485 & 4611 & 1205\\
Canonical recursor checker, shared traversal & 1224 & 1378 & 4079 & 1188\\
\midrule
\multicolumn{5}{@{}l}{\pkg{full\_normalization} (goal-directed)}\\
Raw computations, matched control & 895 & 930 & 2569 & 941\\
Zero normalization, matched control & 890 & 930 & 2532 & 945\\
$\mathsf{F}$, block form & 904 & 962 & 2546 & 960\\
$\mathsf{F}$, expression form & 908 & 965 & 2529 & 963\\
\midrule
\multicolumn{5}{@{}l}{\pkg{computation\_normal}}\\
Compact arity envelope (reference) & 265 & 306 & 516 & 332\\
Public goal, baseline & 383 & 442 & 870 & 507\\
Public goal, scope reuse (\cref{comp:sec:scope}) & 377 & 433 & 860 & 508\\
$\mathsf{Z}_6$, original & 1013 & 1075 & 3095 & 977\\
$\mathsf{Z}_6$, term as computation certificate & 991 & 1049 & 3029 & 971\\
$\mathsf{F}$, initial & 993 & 1046 & 3048 & 965\\
$\mathsf{F}$, neutral-counter test & 992 & 1046 & 3044 & 964\\
$\mathsf{F}$, explicit fuel & 1024 & 1086 & 3204 & 1003\\
$\mathsf{Z}_6$, goal-directed & 972 & 1004 & 2982 & 965\\
$\mathsf{F}$, goal-directed & 982 & 1008 & 3063 & 961\\
$\mathsf{F}$, goal-directed, fuel & 1010 & 1044 & 3189 & 1001\\
\midrule
\multicolumn{5}{@{}l}{\pkg{atomic\_seed\_refinement} (tagged, closed native goals)}\\
Matched raw control: identity normalizer, \rulename{I0} as seventh rule & 1319 & 1507 & 5348 & --\\
$\mathsf{A}_6$ & 1333 & 1521 & 5373 & --\\
\end{longtable}

\paragraph{The matched differences.}
The effect of each change of rules, against the control that shares
everything else, is as follows.
\begin{itemize}
\item \pkg{computational\_equality}: $\mathsf{E}_3$ against the matched
 seven-rule control saves 82 bytes, 49 tokens and 38 AST nodes; gzip grows
 by 9 bytes. The matched control differs from the archived one only by the
 same use of the invariant on proved formulas.
\item \pkg{equality\_from\_iteration}: $\mathsf{E}_2$ against
 $\mathsf{E}_3$ saves 5 tokens, 7 AST nodes, 43 bytes and 11 gzip bytes.
 $\mathsf{U}$ against its control with \rulename{Root} saves 8 tokens, 12
 AST nodes and 53 bytes. $\mathsf{U}$ against $\mathsf{E}_3$ has 28 fewer
 bytes, 20 more tokens and 21 more AST nodes.
\item \pkg{uniform\_step}: $\mathsf{C}_{\mathrm{Step}}$ against the matched
 zero-normal control saves 98 bytes, 31 tokens and 47 AST nodes. The
 delivered file with documentation has 3507 bytes and 1274 gzip bytes. With
 native contexts the checker with \rulename{Refl} and \rulename{Step} is
 smaller than its matched control and larger in tokens and AST nodes than
 the checker of $\mathsf{E}_3$ (959, 1026, 3039).
\item \pkg{seed\_evaluation}: the change from \rulename{Step} to
 \rulename{SeedEval} ties in the helper style and in both goal-directed
 pairs, and saves 2 bytes and 2 tokens with the shared forward traversal;
 of the 148 bytes between the 2950-byte and the 2802-byte checker, 146 come
 from the shared traversal, which benefits both rules. The package reports
 gzip sizes of 818 and 822 bytes for the goal-directed pair and of 872 and
 877 bytes for the shared-traversal pair, \rulename{Step} first. The
 checker with the constructor $U$ is larger than its \rulename{Step}
 control in all three coordinates.
\item \pkg{zero\_normal}: the deletion of the zero computation rule saves
 50 bytes, 11 tokens and 5 AST nodes in each of the two pairs; gzip grows
 by 12 and 13 bytes. The frozen upstream checkers have 3358 and 3382 bytes.
\item \pkg{reflexive\_base}: $\mathsf{R}$ against $\mathsf{Z}_6$ saves 11
 bytes and 2 AST nodes with the same number of tokens; gzip grows by 4
 bytes.
\item \pkg{initializer\_equality}: $\mathsf{S}$ against the one-template
 control saves 50 bytes, 10 tokens and 13 AST nodes; gzip grows by 2 bytes.
 The two rows for the canonical recursor checker measure a refactoring of
 an inherited program and no change of rules.
\item \pkg{full\_normalization}: the expression form of $\mathsf{F}$ against
 the control with zero normalization has 3 fewer bytes, 18 more tokens, 35
 more AST nodes and 18 more gzip bytes, and 40 fewer bytes than the raw
 control. Its auxiliary forward checker has 3100 bytes against 3095 for the
 forward checker of $\mathsf{Z}_6$.
\item \pkg{computation\_normal}: the scope reuse saves 10 bytes, 6 tokens
 and 9 AST nodes; gzip grows by 1 byte. For the calculi, the apparent gain
 of 47 bytes of $\mathsf{F}$ over the original checker of $\mathsf{Z}_6$
 disappears once the latter also takes a term as its computation
 certificate: $\mathsf{F}$ with the neutral-counter test is then 15 bytes
 and 1 token larger, 3 AST nodes and 7 gzip bytes smaller. The goal-directed
 checker of $\mathsf{F}$ is larger than that of $\mathsf{Z}_6$, and fuel
 adds to both.
\item \pkg{atomic\_seed\_refinement}: $\mathsf{A}_6$ against its matched raw
 control, which has the same syntax, raw metadata, validation and
 substitution, is larger by 14 tokens, 14 AST nodes and 25 bytes. The
 deletion of the zero computation rule does not pay for the normalizer.
\end{itemize}

\paragraph{The checkers.}
All checkers are single Python modules that import only \texttt{json} and
\texttt{sys}, read one JSON value, use explicit tests in place of
\texttt{assert}, and accept by printing \texttt{ACCEPT} with exit status
zero. Proof translators and test runners are not imported and are not
trusted. The checkers of \pkg{computational\_equality},
\pkg{equality\_from\_iteration}, \pkg{zero\_normal} and
\pkg{reflexive\_base} take an object with the fields \texttt{variables},
\texttt{context}, \texttt{goal} and \texttt{proof}, with tagged terms in
de Bruijn notation and six proof tags; they check open sequents. Their
\rulename{Eval} instruction takes two premise proofs and nothing else. It
relies on the invariant that the inference function returns a well-scoped
native formula, so that a proved formula of length three is an equation
and a term of length four is an iterator; the matched controls use the same
invariant. The checkers on zero-normal terms reject every supplied iterator
with the literal counter $0$, so that a malformed subterm cannot be erased
before it is validated, and rebuild terms with the smart constructor during
substitution; in \pkg{zero\_normal} this is one conditional after the
rebuilding of a node. The test of \pkg{reflexive\_base} compares the slices
of a validated formula after its tag, which is exact because an equation
has three fields and a binder formula two.

The clause checkers of \pkg{uniform\_step} and \pkg{seed\_evaluation} take
the fields \texttt{goal} and \texttt{proof} and check closed goals. Bound
variables are positive indices and eigenvariables are negative levels;
every substituted term has only negative levels free, so the opening of a
binder needs no shift. In the helper style the change from \rulename{Step}
to \rulename{SeedEval} is two lines: the comparison of the first source
with the initializer in place of the predecessor iterate, and the
predecessor iterate from the new seed inside the conclusion. The checker of
\pkg{initializer\_equality} implements \rulename{Seed} by one comparison of
the two left sides and, failing that, one comparison with the initializer of
an iterator; no substitution is performed by that instruction.

The selected checker of \pkg{full\_normalization} is goal-directed: one
recursive traversal checks arities and sorts, renames variables and
contracts exposed redexes, and validation is the requirement that the
identity action returns its input, with all children checked before a
redex is erased. A goal that is a retained assumption or a reflexive
equation closes without reading its proof field, in the checker and in its
controls. In \pkg{computation\_normal} a normal term is its own
reflexivity certificate, and the fuelled variant adds the field
\texttt{fuel}.

The checker of $\mathsf{A}_6$ in \pkg{atomic\_seed\_refinement} takes the
fields \texttt{goal} and \texttt{proof} and checks closed goals with $=$,
$\Q{}$ and $\D{}{}$. Terms and formulas are tagged, with zero-based
de Bruijn indices: the body of an iterator binds the index $0$, $\Q{y}$ binds
$y$ at $0$, and $\D{x}{y}$ binds $y$ at $0$ and $x$ at $1$. The six proof
tags carry the raw metadata of \cref{comp:def:asix}; the checker validates
every raw field before it uses or normalizes it, since normalization can
erase an ill-scoped subterm, forms the raw instance by capture-avoiding
substitution, and normalizes afterwards. Contexts hold normalized formulas.
The matched control has the identity normalizer and a seventh tag for
\rulename{I0}; it accepts a primitive zero computation and rejects the
derivation of \cref{comp:thm:asizero}, and the checker of $\mathsf{A}_6$
rejects the tag for \rulename{I0}.

\paragraph{Executed checks.}
The packages record the following runs, none with an unexpected outcome;
they support the implementations and are not proofs of the theorems.
\pkg{computational\_equality}: 38 inherited fixtures, 14 translation round
trips within the bound $9N$, 29 further round trips, 51 adversarial cases
and 206 process executions. \pkg{equality\_from\_iteration}: 460 process
checks, 62 translation checks, 1280 normalization checks and 320
substitution checks. \pkg{uniform\_step}: 376 process checks of the clause
checker, 100 of the native checkers with 38 translation round trips, and
2104 comparisons of the opening operation. \pkg{seed\_evaluation}: 398
process checks, 166 translation round trips and 92 rejections, among
others, and 83 certificates for the calculus with $U$.
\pkg{zero\_normal}: 500 process checks, the translation of 64 inherited
certificates and 4832 checks of the coherence with substitution.
\pkg{reflexive\_base}: 13 translation cases, 16 malformed inputs and 86
process checks, including two substitutions that turn the guard from false
to true. \pkg{initializer\_equality}: 600 process checks and 31 translation
round trips. \pkg{full\_normalization}: 680 process checks of the
goal-directed family and 132 of the forward checker, with the two
certificates of \cref{comp:ex:modechange}. \pkg{computation\_normal}: 272
runs and 9000 generated goals for \cref{comp:thm:scope}, and a recovered
case introduction at the witness $0$ with the context $a=b$ and the matrix
$\It{y}{a}{u}{b}=\It{y}{x}{u}{0}$. \pkg{atomic\_seed\_refinement}: 20
persisted certificates, among them an equation induction that uses its
hypothesis, the zero computation at a compound initializer, and a case
introduction at the raw witness $\It{0}{0}{u}{0}$ that uses its case
assumption, rejected when the witness is changed to $0$; with the
malformed inputs and the control checks, 54 executions in ordinary and
optimized Python.

%%%%%%%%%%%%%%%%%%%%%%%%%%%%%%%%%%%%%%%%%%%%%%%%%%%%%%%%%%%%%%%%%%%%%%%%%%%%
\section{Further questions}
\label{comp:sec:questions}

This section collects, from all ten packages, the questions that they
leave open, the claims that they state without proof and the results that
they report as not achieved.

\subsection{What is not achieved}

\begin{itemize}
\item No package proves a lower bound or a minimality statement, for rule
 counts or for checker sizes. Every package states that its checker does not
 improve on the compact arity-coded checker of the repository (516 bytes;
 521 bytes at the snapshot used by \pkg{zero\_normal}) or on the fully
 validated profiles of 866, 870 and 788 bytes that have different
 certificate formats and validation contracts. A port of
 \rulename{Eval}, of the identity frame, of \rulename{Seed} or of zero
 normalization to those encodings is not established.
\item Moving a computation into the syntax does not by itself reduce the
 size of a complete checker. The measurements of \cref{comp:sec:sizes} for
 $\mathsf{U}$ against $\mathsf{E}_2$, for the constructor $U$, and for
 $\mathsf{F}$ against $\mathsf{Z}_6$ are mixed or negative, and gzip sizes
 grow in most matched pairs.
\item \pkg{seed\_evaluation} reports that \rulename{SeedEval} gives no
 convincing reduction of checker size: two of its three matched pairs tie.
\item \pkg{zero\_normal} examines the rule of common-source congruence,
 from $s=t$ and $s=u$ infer $h(t)=h(u)$, in place of \rulename{Tr}. With
 reflexivity it is interderivable with \rulename{Tr}: the instance whose
 first premise is a reflexivity is congruence, and the instance with the
 identity template composes two equations with a common left side. An
 apparent saving of 36 bytes
 disappears after the same factoring is applied to the control; both
 programs measure 3322 bytes, 1084 tokens and 1141 AST nodes.
\item \pkg{computation\_normal} recommends that $\mathsf{F}$ not replace
 the zero-normal kernel on the ground of its shorter rule table: the
 matched checker size does not improve uniformly, normal forms can be
 enormous, and the uniform normalizer is not primitive recursive
 (\cref{comp:prop:notpr}).
\item \pkg{atomic\_seed\_refinement} reports that its target, a complete
 checker of a calculus bi-interpretable with $\iPRA$ that is smaller than
 the compact checkers of the repository, is not reached. The checker of
 $\mathsf{A}_6$ is larger than its matched control
 (\cref{comp:sec:sizes}), and the compact guarded checkers, with their
 different interfaces, attain the separate minima of 248 tokens, 270 AST
 nodes and 479 bytes in three different programs. No minimality of six
 rules is claimed: the chapter already contains other six-rule calculi with
 stronger normalization.
\item The Lean formalization covers the calculi as recorded in the remarks
 of this chapter. It does not cover, in the forms stated here: the clause
 calculi $\mathsf{C}_{\mathrm{Step}}$ and $\mathsf{C}_{\mathrm{Seed}}$ with
 their clause contexts, the calculus $\mathsf{C}_{\mathrm{App}}$ and the
 unit normalization, $\mathsf{Z}_6$ and $\mathsf{S}$ with the single
 introduction, the fresh-zero lemma and the compound envelope, the fuelled
 relation of \cref{comp:prop:fuel}, the raw-metadata transfer of
 \cref{comp:thm:rawtransfer} and the calculus $\mathsf{A}_6$, and any of the
 Python checkers.
\end{itemize}

\subsection{Claims stated without proof}

\begin{itemize}
\item \pkg{seed\_evaluation} states that the interpretation in which every
 literal equation is false obstructs the deletion of reflexivity from the
 clause calculi, and that moving reflexivity into an initial universal
 clause still leaves a reflexivity axiom to count.
 \Cref{comp:prop:reflsharp} proves the corresponding statement for
 $\mathsf{E}_2$; the statement for the clause calculi is not proved in the
 package.
\item \Cref{comp:thm:app}(c) is asserted in the notes of
 \pkg{seed\_evaluation} for arbitrary clause contexts, by reference to the
 completeness of the clause calculus for $\Teq$; the argument given covers
 contexts of assumption clauses. The general statement requires the
 completeness of $\mathsf{C}_{\mathrm{Step}}$ for arbitrary clauses.
\item \pkg{full\_normalization} argues that its executable traversal, which
 fuses substitution with normalization and passes the unevaluated recursive
 iterator into the substitution of the body, computes $\nf(t[\sigma])$:
 every step contracts a redex of the materialized instance, termination is
 \cref{comp:thm:sn}, and the result is normal. The argument is given in
 prose and is not carried out as a proof about the program.
\end{itemize}

\subsection{Open questions}

\begin{question}[Lower bounds]
For a fixed term grammar, certificate language, validation contract and set
of primitive operations, what is the least size of a checker of a calculus
bi-interpretable with $\iPRA$? No package fixes such a class.
\pkg{computational\_equality} asks for an order of simplicity that separates
rule counts, constructors of the syntax, tokens, bytes and certificate
sizes; \pkg{initializer\_equality} and \pkg{full\_normalization} ask for
lower bounds once a class of representations is declared.
\end{question}

\begin{question}[A smaller basis for arithmetic and equality]
Can the interface of reflexivity and one evaluation rule be reduced
further while a small normalization layer is counted in full?
\Cref{comp:prop:evalsharp,comp:prop:reflsharp} show that no rule of
$\mathsf{E}_2$ can be deleted from that inventory; they are not lower
bounds across term languages. \pkg{equality\_from\_iteration} notes that a
different basis has to account for the freedom of the initializer, for the
reflection of equality (\cref{base:thm:reflect}) and for both round trips;
\pkg{reflexive\_base} asks whether another arithmetic or logical rule can be
removed with a net reduction of the complete checking algorithm.
\end{question}

\begin{question}[An introduction that is stable under substitution]
\rulename{Eval}, \rulename{Step} and \rulename{Seed} are stable under
substitution (\cref{comp:prop:phisubst,comp:prop:stepsubst,comp:lem:seedsubst}),
and \rulename{WI} is not (\cref{comp:rem:single}). Is there a calculus of
six rules in which substitution in derivations is proved by a direct
induction, with a single introduction that needs no test on the witness,
without enlarging the complete checker and without induction beyond literal
equations? The guard of \cref{comp:sec:reflexive} is preserved and not
reflected. \pkg{uniform\_step} asks whether the logical rules can be merged
further while the distinction between coverage and equation induction and
the eigenvariable conditions are kept.
\end{question}

\begin{question}[Sizes of derivations]
The congruence of \cref{comp:thm:evalcong} adds eight nodes, that of
\cref{comp:thm:seedcong} four, that of \cref{comp:lem:stepderived} four,
and the elimination of \rulename{Root} adds two per instance
(\cref{comp:thm:root}); in each case the labels contain copies of the
template and of the endpoints. What are the least numbers of nodes for
these fragments, for trees and for derivations with shared subproofs, and
what are the sizes of serialized derivations and the checking times on an
arithmetic library? Can sharing reduce the cost without a decoder that is
larger than the rule it removes? Can an interface based on the initializer
shorten broader families of derivations than the single computations of
\cref{comp:prop:seedcert}, without the sevenfold expansion of
\cref{comp:lem:seedstep}(c)?
\end{question}

\begin{question}[Normalization between zero and full]
Is there a normalization between the zero normalization and the full
normalization that reduces both the arithmetic interface and the complete
trusted implementation? The unit normalization of \cref{comp:sec:app} has a
primitive recursive normalizer and a rule without a counter, and its
measured checker is larger than its control.
\end{question}

\begin{question}[Normalization below zero normalization]
For the atomic-seed normalizer, which violates ($\nu$4)
(\cref{comp:prop:nfatomic}), \cref{comp:thm:asix} is proved with raw rule
metadata. Does the calculus $\Nat^{\nf_{\mathrm a}}$, defined as in
\cref{comp:def:nukernel} with atomic-seed normal metadata and normalizing
substitution, still derive exactly the normalized sequents of $\Nat$, and is
$\mathsf{A}_6$ complete when its metadata are restricted to normal terms?
\pkg{atomic\_seed\_refinement} does not claim that the atomic-seed normal
terms with normalizing substitution form a free term language, and its
derivation of the zero computation needs a body that is not normal. It also
leaves open a compact implementation of atomic contraction in a level coding
(\cref{comp:rem:rebasing}), with its sort test and the treatment of dummy
slots, that is smaller than the present compact checkers.
\end{question}

\begin{question}[Classes of provable bases]
Which decidable classes of bases that are derivable in every context give
the best exchange between the size of a checker and the length of
derivations for the guard of \cref{comp:def:reflexive}? The classes ``$r=r$''
and ``$0=0$'' of \cref{comp:prop:rvariants} already differ. Does
$\mathsf{R}^{\mathrm{raw}}$, which needs no normalization, lower the size of
a complete compact checker, and does the guard ``$0=0$'' suffice on raw
terms?
\end{question}

\begin{question}[Other positions of a recursor]
Which arithmetic primitives support both the identity and the application
of an arbitrary template through one distinguished argument, as the
initializer of the scalar iterator does at the counters $0$ and $1$
(\cref{comp:lem:seedschema,comp:lem:beta})? \pkg{initializer\_equality}
requires equality interpreted by equality and both formula round trips in
every such comparison.
\end{question}

\begin{question}[Narrower induction]
\pkg{seed\_evaluation} asks whether the approach by rigidity of the
successor, which derives induction for equations with the right side $0$,
can be completed to the reflection of equality without importing general
equation induction at an intermediate step, and notes that a recursor whose
body sees the counter gives a simpler predecessor (\cref{base:thm:pred}) and
does not by itself supply the reflection argument.
\end{question}

\begin{question}[Formalization and a common benchmark]
Formalize the statements listed above as not covered, in particular the
conversion relation with finite witnesses and the fuelled relation, the
fresh-zero lemma and the compound envelope next to \Lean{F6Prf.dIZero}, and
the clause calculi with \rulename{Step}, \rulename{SeedEval} and
\rulename{App}. Separately, measure the compact, the goal-directed and the
forward checkers under common contracts for the validity of syntax, open
sequents, sharing in certificates and the behavior on errors, so that a
bound for one encoding is not read as a bound for all presentations.
\end{question}
